\documentclass[11pt]{article}
\usepackage[margin=0.85in]{geometry}
\usepackage{fontspec}
\setmainfont{Times New Roman}
\newfontfamily\EmojiFont[Renderer=HarfBuzz]{Apple Color Emoji}
\newfontfamily\SymbolFont{Arial Unicode MS}
\newcommand{\EmojiGlyph}[1]{{\normalfont\EmojiFont #1}}
\newcommand{\SymbolGlyph}[1]{{\normalfont\SymbolFont #1}}
\usepackage{newunicodechar}
\usepackage{microtype}
\usepackage{setspace}
\setstretch{1.08}
\usepackage{hyperref}
\usepackage{xurl}
\urlstyle{same}
\hypersetup{colorlinks=true,linkcolor=blue,urlcolor=blue}
\usepackage{enumitem}
\setlist{leftmargin=*,itemsep=2pt,topsep=2pt}
\usepackage{array}
\usepackage{longtable}
\usepackage{booktabs}
\usepackage{amssymb}
\usepackage{ragged2e}
\renewcommand{\arraystretch}{1.15}
\setlength{\parskip}{0.45em}
\setlength{\parindent}{0pt}
\setlength{\emergencystretch}{3em}
\sloppy
\newcommand{\UncheckedBox}{$\square$}
\newcommand{\CheckedBox}{$\blacksquare$}
\title{Partner Built Prompt Library\\LaTeX Translation and Dependency Map}
\author{financial-services-plugins repository}
\date{Generated on 2026-03-05 09:31 }
\newunicodechar{⚠}{{\EmojiGlyph{⚠}}}
\newunicodechar{✅}{{\EmojiGlyph{✅}}}
\newunicodechar{❌}{{\EmojiGlyph{❌}}}
\newunicodechar{🔗}{{\EmojiGlyph{🔗}}}
\newunicodechar{▸}{{\SymbolGlyph{▸}}}
\newunicodechar{✓}{{\SymbolGlyph{✓}}}
\begin{document}
\maketitle
\tableofcontents
\newpage

\section{Repository Structure Overview}

\subsection{Folder Tree}
\textbf{Root directory: partner-built}
\begin{itemize}[leftmargin=*,itemsep=1pt,topsep=2pt]
\item \hspace*{0.0em}lseg
\item \hspace*{0.0em}spglobal
\item \hspace*{0.9em}lseg/commands
\item \hspace*{0.9em}lseg/CONNECTORS.md
\item \hspace*{0.9em}lseg/README.md
\item \hspace*{0.9em}lseg/skills
\item \hspace*{0.9em}spglobal/README.md
\item \hspace*{0.9em}spglobal/skills
\item \hspace*{1.8em}lseg/commands/analyze-bond-basis.md
\item \hspace*{1.8em}lseg/commands/analyze-bond-rv.md
\item \hspace*{1.8em}lseg/commands/analyze-fx-carry.md
\item \hspace*{1.8em}lseg/commands/analyze-option-vol.md
\item \hspace*{1.8em}lseg/commands/analyze-swap-curve.md
\item \hspace*{1.8em}lseg/commands/macro-rates.md
\item \hspace*{1.8em}lseg/commands/research-equity.md
\item \hspace*{1.8em}lseg/commands/review-fi-portfolio.md
\item \hspace*{1.8em}lseg/skills/bond-futures-basis
\item \hspace*{1.8em}lseg/skills/bond-relative-value
\item \hspace*{1.8em}lseg/skills/equity-research
\item \hspace*{1.8em}lseg/skills/fixed-income-portfolio
\item \hspace*{1.8em}lseg/skills/fx-carry-trade
\item \hspace*{1.8em}lseg/skills/macro-rates-monitor
\item \hspace*{1.8em}lseg/skills/option-vol-analysis
\item \hspace*{1.8em}lseg/skills/swap-curve-strategy
\item \hspace*{1.8em}spglobal/skills/earnings-preview-beta
\item \hspace*{1.8em}spglobal/skills/funding-digest
\item \hspace*{1.8em}spglobal/skills/tear-sheet
\item \hspace*{2.7em}lseg/skills/bond-futures-basis/SKILL.md
\item \hspace*{2.7em}lseg/skills/bond-relative-value/SKILL.md
\item \hspace*{2.7em}lseg/skills/equity-research/SKILL.md
\item \hspace*{2.7em}lseg/skills/fixed-income-portfolio/SKILL.md
\item \hspace*{2.7em}lseg/skills/fx-carry-trade/SKILL.md
\item \hspace*{2.7em}lseg/skills/macro-rates-monitor/SKILL.md
\item \hspace*{2.7em}lseg/skills/option-vol-analysis/SKILL.md
\item \hspace*{2.7em}lseg/skills/swap-curve-strategy/SKILL.md
\item \hspace*{2.7em}spglobal/skills/earnings-preview-beta/report-template.md
\item \hspace*{2.7em}spglobal/skills/earnings-preview-beta/SKILL.md
\item \hspace*{2.7em}spglobal/skills/funding-digest/references
\item \hspace*{2.7em}spglobal/skills/funding-digest/SKILL.md
\item \hspace*{2.7em}spglobal/skills/tear-sheet/references
\item \hspace*{2.7em}spglobal/skills/tear-sheet/SKILL.md
\item \hspace*{3.6em}spglobal/skills/funding-digest/references/sector-seeds.md
\item \hspace*{3.6em}spglobal/skills/tear-sheet/references/corp-dev.md
\item \hspace*{3.6em}spglobal/skills/tear-sheet/references/equity-research.md
\item \hspace*{3.6em}spglobal/skills/tear-sheet/references/ib-ma.md
\item \hspace*{3.6em}spglobal/skills/tear-sheet/references/sales-bd.md
\end{itemize}

\section{Prompt Dependency Structure}

\subsection{Skill-to-Resource Dependencies}

\begingroup
\small
\setlength{\tabcolsep}{2pt}
\begin{longtable}{>{\RaggedRight\arraybackslash\hspace{0pt}}p{0.480\textwidth}>{\RaggedRight\arraybackslash\hspace{0pt}}p{0.480\textwidth}}
\toprule
{} \textbf{Skill Prompt} & \textbf{Referenced Prompt File} \\
\midrule
{} spglobal/ skills/ earnings-preview-beta/ SKILL.md & spglobal/ skills/ earnings-preview-beta/ report-template.md \\
{} spglobal/ skills/ funding-digest/ SKILL.md & spglobal/ skills/ funding-digest/ references/ sector-seeds.md \\
{} spglobal/ skills/ tear-sheet/ SKILL.md & spglobal/ skills/ tear-sheet/ references/ corp-dev.md \\
{} spglobal/ skills/ tear-sheet/ SKILL.md & spglobal/ skills/ tear-sheet/ references/ equity-research.md \\
{} spglobal/ skills/ tear-sheet/ SKILL.md & spglobal/ skills/ tear-sheet/ references/ ib-ma.md \\
{} spglobal/ skills/ tear-sheet/ SKILL.md & spglobal/ skills/ tear-sheet/ references/ sales-bd.md \\
\bottomrule
\end{longtable}
\endgroup

\newpage

\section{Translated Prompt Corpus}

This section translates every markdown prompt under the partner-built directory into LaTeX-native structure, preserving hierarchy and dependency flow.

\subsection{Directory: lseg}

\subsection{File: lseg/CONNECTORS.md}\label{file:lseg-CONNECTORS.md}

\paragraph{Connectors}\mbox{}\par

This plugin connects to the \textbf{LFA MCP Server}, which provides financial analytics tools from LSEG (London Stock Exchange Group). All tools are served by a single MCP server — no additional connectors are needed.


\subparagraph{How Commands Reference Tools}\mbox{}\par

Commands in this plugin reference MCP tools by their exact tool name (e.g., \emph{bond\_\allowbreak{}price}, \emph{interest\_\allowbreak{}rate\_\allowbreak{}curve}). The tools are organized into categories for clarity:


\subparagraph{Tool Categories}\mbox{}\par

\begingroup
\small
\setlength{\tabcolsep}{2pt}
\begin{longtable}{>{\RaggedRight\arraybackslash\hspace{0pt}}p{0.240\textwidth}>{\RaggedRight\arraybackslash\hspace{0pt}}p{0.240\textwidth}>{\RaggedRight\arraybackslash\hspace{0pt}}p{0.240\textwidth}>{\RaggedRight\arraybackslash\hspace{0pt}}p{0.240\textwidth}}
\toprule
{} \textbf{Category} & \textbf{Placeholder} & \textbf{Tools} & \textbf{Description} \\
\midrule
{} Bond Pricing & \emph{\textasciitilde{}\textasciitilde{}bond-pricing} & \emph{bond\_\allowbreak{} price}, \emph{bond\_\allowbreak{} future\_\allowbreak{} price} & Price bonds and bond futures with full analytics \\
{} FX Pricing & \emph{\textasciitilde{}\textasciitilde{}fx-pricing} & \emph{fx\_\allowbreak{} spot\_\allowbreak{} price}, \emph{fx\_\allowbreak{} forward\_\allowbreak{} price} & FX spot and forward rate pricing \\
{} Interest Rate Curves & \emph{\textasciitilde{}\textasciitilde{}ir-curves} & \emph{interest\_\allowbreak{} rate\_\allowbreak{} curve}, \emph{inflation\_\allowbreak{} curve} & Government yield curves and inflation breakevens \\
{} Credit Curves & \emph{\textasciitilde{}\textasciitilde{}credit-curves} & \emph{credit\_\allowbreak{} curve} & Credit spread curves by issuer type \\
{} FX Curves & \emph{\textasciitilde{}\textasciitilde{}fx-curves} & \emph{fx\_\allowbreak{} forward\_\allowbreak{} curve} & FX forward point curves \\
{} Options & \emph{\textasciitilde{}\textasciitilde{}options} & \emph{option\_\allowbreak{} value}, \emph{option\_\allowbreak{} template\_\allowbreak{} list} & Option valuation with Greeks \\
{} Swaps & \emph{\textasciitilde{}\textasciitilde{}swaps} & \emph{ir\_\allowbreak{} swap} & Interest rate swap pricing \\
{} Volatility Surfaces & \emph{\textasciitilde{}\textasciitilde{}volatility} & \emph{fx\_\allowbreak{} vol\_\allowbreak{} surface}, \emph{equity\_\allowbreak{} vol\_\allowbreak{} surface} & FX and equity implied vol surfaces \\
{} Quantitative Analytics & \emph{\textasciitilde{}\textasciitilde{}qa} & \emph{qa\_\allowbreak{} ibes\_\allowbreak{} consensus}, \emph{qa\_\allowbreak{} company\_\allowbreak{} fundamentals}, \emph{qa\_\allowbreak{} historical\_\allowbreak{} equity\_\allowbreak{} price}, \emph{qa\_\allowbreak{} macroeconomic} & Analyst estimates, fundamentals, prices, macro data \\
{} Time Series & \emph{\textasciitilde{}\textasciitilde{}time-series} & \emph{tscc\_\allowbreak{} historical\_\allowbreak{} pricing\_\allowbreak{} summaries} & Historical pricing summaries (interday/ intraday) \\
{} Fixed Income Analytics & \emph{\textasciitilde{}\textasciitilde{}yieldbook} & \emph{yieldbook\_\allowbreak{} bond\_\allowbreak{} reference}, \emph{yieldbook\_\allowbreak{} cashflow}, \emph{yieldbook\_\allowbreak{} scenario}, \emph{fixed\_\allowbreak{} income\_\allowbreak{} risk\_\allowbreak{} analytics} & Bond reference data, cashflows, scenarios, OAS/ duration \\
\bottomrule
\end{longtable}
\endgroup

\subparagraph{Complete Tool Reference}\mbox{}\par

\subparagraph{Bond Domain}\mbox{}\par
\begin{itemize}[leftmargin=*,itemsep=2pt,topsep=2pt]
\item {} \textbf{@@TOKEN0@@} — Calculate bond pricing, valuation, and analytics. Accepts ISIN, RIC, CUSIP, or AssetId. Returns yield, duration, convexity, DV01, accrued interest. Supports what-if scenarios via price/yield overrides.
\item {} \textbf{@@TOKEN0@@} — Calculate bond future pricing and analytics. Returns fair value, cheapest-to-deliver identification, delivery basket, conversion factors, and contract DV01.
\end{itemize}

\subparagraph{FX Domain}\mbox{}\par
\begin{itemize}[leftmargin=*,itemsep=2pt,topsep=2pt]
\item {} \textbf{@@TOKEN0@@} — FX spot rate pricing for ISO currency pairs. Returns mid/bid/ask rates.
\item {} \textbf{@@TOKEN0@@} — FX forward rate pricing at specific tenors or dates. Returns forward points, outright rates, and carry.
\end{itemize}

\subparagraph{Curves Domain}\mbox{}\par
\begin{itemize}[leftmargin=*,itemsep=2pt,topsep=2pt]
\item {} \textbf{@@TOKEN0@@} — Government yield curves. Two-phase: list available curves, then calculate curve points. Returns par/zero rates, discount factors, forward rates.
\item {} \textbf{@@TOKEN0@@} — Credit spread curves. Search by country and issuer type (Corporate, Sovereign, Agency, etc.), then calculate spread term structure.
\item {} \textbf{@@TOKEN0@@} — Inflation breakeven curves. Search by country/currency, then calculate breakeven rates and real yields.
\item {} \textbf{@@TOKEN0@@} — FX forward point curves. List curves, then calculate forward points across all standard tenors.
\end{itemize}

\subparagraph{Swaps Domain}\mbox{}\par
\begin{itemize}[leftmargin=*,itemsep=2pt,topsep=2pt]
\item {} \textbf{@@TOKEN0@@} — Interest rate swap pricing. Two-phase: list templates by currency/index, then price swaps at specified tenors. Returns par rates, DV01, NPV.
\end{itemize}

\subparagraph{Options Domain}\mbox{}\par
\begin{itemize}[leftmargin=*,itemsep=2pt,topsep=2pt]
\item {} \textbf{@@TOKEN0@@} — Option valuation supporting vanilla, barrier, binary, and Asian options. Returns premium, full Greeks (delta, gamma, vega, theta, rho), and risk metrics.
\item {} \textbf{@@TOKEN0@@} — List available option templates for pricing.
\end{itemize}

\subparagraph{Volatility Domain}\mbox{}\par
\begin{itemize}[leftmargin=*,itemsep=2pt,topsep=2pt]
\item {} \textbf{@@TOKEN0@@} — FX volatility surface generation using SABR model. Returns vol surface across tenors and delta strikes.
\item {} \textbf{@@TOKEN0@@} — Equity implied volatility surface. Supports equities/indices via RIC and futures via RICROOT.
\end{itemize}

\subparagraph{Quantitative Analytics Domain}\mbox{}\par
\begin{itemize}[leftmargin=*,itemsep=2pt,topsep=2pt]
\item {} \textbf{@@TOKEN0@@} — IBES analyst consensus estimates (EPS, revenue, EBITDA, DPS). Forward-looking estimates with analyst count, dispersion, and high/low ranges.
\item {} \textbf{@@TOKEN0@@} — Reported company financials (income statement, balance sheet metrics). Historical fiscal year data.
\item {} \textbf{@@TOKEN0@@} — Historical equity prices with OHLCV, total returns, and beta.
\item {} \textbf{@@TOKEN0@@} — Macroeconomic indicators database. Search by mnemonic or description, retrieve latest values or time series.
\end{itemize}

\subparagraph{Time Series Domain}\mbox{}\par
\begin{itemize}[leftmargin=*,itemsep=2pt,topsep=2pt]
\item {} \textbf{@@TOKEN0@@} — Historical pricing summaries for any RIC. Supports interday (daily, weekly, monthly) and intraday (1min to 1hr) intervals.
\end{itemize}

\subparagraph{Fixed Income Analytics (YieldBook) Domain}\mbox{}\par
\begin{itemize}[leftmargin=*,itemsep=2pt,topsep=2pt]
\item {} \textbf{@@TOKEN0@@} — Bond reference data: security type, sector, ratings, coupon, maturity, issuer.
\item {} \textbf{@@TOKEN0@@} — Bond cashflow projections: future coupon and principal payment schedules.
\item {} \textbf{@@TOKEN0@@} — Bond scenario analysis: price/yield under parallel rate shifts.
\item {} \textbf{@@TOKEN0@@} — Bond risk analytics: OAS, effective duration, key rate durations, convexity.
\end{itemize}

\vspace{0.8em}\hrule\vspace{0.8em}

\subsection{File: lseg/README.md}\label{file:lseg-README.md}

\paragraph{LSEG Financial Analytics Plugin}\mbox{}\par

Price bonds, analyze yield curves, evaluate FX carry trades, value options, and build macro dashboards using LSEG financial data and analytics.


\subparagraph{What This Plugin Does}\mbox{}\par

This plugin packages LSEG's financial analytics MCP tools into 8 high-level workflows that stitch together multiple tool calls for common financial analysis tasks. Instead of calling individual tools one at a time, each command orchestrates 4-5 tools into a cohesive analysis.


\subparagraph{Commands}\mbox{}\par

\begingroup
\small
\setlength{\tabcolsep}{2pt}
\begin{longtable}{>{\RaggedRight\arraybackslash\hspace{0pt}}p{0.480\textwidth}>{\RaggedRight\arraybackslash\hspace{0pt}}p{0.480\textwidth}}
\toprule
{} \textbf{Command} & \textbf{Description} \\
\midrule
{} \emph{/ analyze-bond-rv} & Analyze bond relative value with spread decomposition and scenario stress testing \\
{} \emph{/ analyze-fx-carry} & Evaluate FX carry trade opportunities with spot, forwards, vol surface, and historical context \\
{} \emph{/ research-equity} & Generate equity research snapshot with consensus estimates, fundamentals, and price performance \\
{} \emph{/ analyze-swap-curve} & Analyze the swap curve with government and inflation overlays for curve trade ideas \\
{} \emph{/ analyze-option-vol} & Analyze option volatility with vol surface, Greeks, and implied vs realized comparison \\
{} \emph{/ review-fi-portfolio} & Review a fixed income portfolio with pricing, cashflows, and scenario analysis \\
{} \emph{/ macro-rates} & Build a macro and rates dashboard with economic indicators, yield curves, and swap spreads \\
{} \emph{/ analyze-bond-basis} & Analyze bond futures basis with CTD identification and implied repo rate \\
\bottomrule
\end{longtable}
\endgroup

\subparagraph{Skills}\mbox{}\par

Each command is backed by a corresponding skill that provides deep domain knowledge:


\begingroup
\small
\setlength{\tabcolsep}{2pt}
\begin{longtable}{>{\RaggedRight\arraybackslash\hspace{0pt}}p{0.480\textwidth}>{\RaggedRight\arraybackslash\hspace{0pt}}p{0.480\textwidth}}
\toprule
{} \textbf{Skill} & \textbf{Domain Knowledge} \\
\midrule
{} \emph{bond-relative-value} & Spread frameworks, G-spread/ Z-spread/ OAS, rich-cheap analysis \\
{} \emph{fx-carry-trade} & Carry mechanics, carry-to-vol ratios, G10 and EM carry dynamics \\
{} \emph{equity-research} & IBES consensus interpretation, fundamental analysis, valuation metrics \\
{} \emph{swap-curve-strategy} & Swap curve construction, curve trades, real rate analysis \\
{} \emph{option-vol-analysis} & Vol surface interpretation, SABR model, Greeks, implied vs realized vol \\
{} \emph{fixed-income-portfolio} & Portfolio analytics, key rate duration, cashflow analysis, scenario testing \\
{} \emph{macro-rates-monitor} & Macro indicators, yield curve shapes, real rates, financial conditions \\
{} \emph{bond-futures-basis} & CTD mechanics, basis calculation, implied repo, delivery options \\
\bottomrule
\end{longtable}
\endgroup

\subparagraph{Integrations}\mbox{}\par

This plugin connects to the \textbf{LFA MCP Server} which provides access to LSEG financial data and analytics across these domains:


\begin{itemize}[leftmargin=*,itemsep=2pt,topsep=2pt]
\item {} \textbf{Bond Pricing} — Bond and bond future valuation
\item {} \textbf{FX Pricing} — Spot and forward rates
\item {} \textbf{Curves} — Interest rate, credit, inflation, and FX forward curves
\item {} \textbf{Swaps} — Interest rate swap pricing
\item {} \textbf{Options} — Option valuation with full Greeks
\item {} \textbf{Volatility} — FX and equity implied volatility surfaces
\item {} \textbf{Quantitative Analytics} — Analyst estimates, company fundamentals, equity prices, macro data
\item {} \textbf{Time Series} — Historical pricing summaries
\item {} \textbf{YieldBook} — Fixed income reference data, cashflows, scenarios, and risk analytics
\end{itemize}

See \href{CONNECTORS.md}{CONNECTORS.md} for the complete tool reference.


\subparagraph{Installation}\mbox{}\par

\textbf{Structured Template}
\begin{itemize}[leftmargin=*,itemsep=2pt,topsep=2pt]
\item {} claude plugins add LSEG
\end{itemize}

\subparagraph{Requirements}\mbox{}\par

\begin{itemize}[leftmargin=*,itemsep=2pt,topsep=2pt]
\item {} Access to the LSEG MCP Server with valid credentials
\item {} LSEG data entitlements for the relevant product offerings
\end{itemize}

\vspace{0.8em}\hrule\vspace{0.8em}

\subsection{Command: analyze-bond-basis.md}\label{file:lseg-commands-analyze-bond-basis.md}

\paragraph{File Metadata}\mbox{}\par
\begingroup
\small
\setlength{\tabcolsep}{2pt}
\begin{longtable}{>{\RaggedRight\arraybackslash\hspace{0pt}}p{0.480\textwidth}>{\RaggedRight\arraybackslash\hspace{0pt}}p{0.480\textwidth}}
\toprule
{} \textbf{Field} & \textbf{Value} \\
\midrule
{} description & Analyze the bond futures basis with CTD identification, implied repo rate, and basis trade assessment \\
{} argument-hint & <bond future RIC e.g. FGBLc1> \\
\bottomrule
\end{longtable}
\endgroup


\paragraph{Analyze Bond Futures Basis}\mbox{}\par

\begin{quote}
``This command uses LSEG bond future pricing, bond pricing, yield curves, and historical data tools. See \href{../CONNECTORS.md}{CONNECTORS.md} for available tools.''
\end{quote}

Analyze the bond futures basis by pricing the future, identifying the cheapest-to-deliver bond, computing gross and net basis, and assessing basis trade opportunities.


See the \textbf{bond-futures-basis} skill for domain knowledge on basis mechanics and trading strategies.


\subparagraph{Workflow}\mbox{}\par

\subparagraph{1. Gather Input}\mbox{}\par

Ask the user for:

\begin{itemize}[leftmargin=*,itemsep=2pt,topsep=2pt]
\item {} Bond future RIC (required) — e.g., FGBLc1 (Euro Bund), TYc1 (US 10Y Note), FFIc1 (UK Gilt)
\item {} Market data date (optional, defaults to today)
\end{itemize}

\subparagraph{2. Price the Bond Future}\mbox{}\par

Call \emph{bond\_\allowbreak{}future\_\allowbreak{}price} with the future RIC.


Extract: fair price, CTD bond identifier, delivery basket with conversion factors, contract DV01, delivery dates.


\subparagraph{3. Price the CTD Bond}\mbox{}\par

Call \emph{bond\_\allowbreak{}price} for the CTD identifier from Step 2.


Extract: clean/dirty price, yield, duration, DV01, accrued interest, coupon.


Compute: gross basis, invoice price, carry, net basis.


\subparagraph{4. Compute Implied Repo Rate}\mbox{}\par

Call \emph{interest\_\allowbreak{}rate\_\allowbreak{}curve} (list then calculate) for the future's currency. Use short-end rate as repo proxy.


Compute implied repo rate and compare to market repo.


\subparagraph{5. Track Historical Basis}\mbox{}\par

Call \emph{tscc\_\allowbreak{}historical\_\allowbreak{}pricing\_\allowbreak{}summaries} for both the future and CTD bond with \emph{tenor: ``3M''}, \emph{interval: ``P1D''}.


Assess: basis trend, volatility, and historical range.


\subparagraph{6. Sovereign Credit Context}\mbox{}\par

Call \emph{credit\_\allowbreak{}curve} for the relevant sovereign (e.g., ``DE'' for Bund, ``US'' for Treasury).


\subparagraph{7. Synthesize the Report}\mbox{}\par

Present: future summary table, CTD bond analytics, basis calculation table (gross/net basis, implied repo vs market repo), historical context, and trade recommendation (long basis / short basis / neutral).


\subparagraph{Output Format}\mbox{}\par

Lead with the basis trade assessment (long/short/neutral) and implied repo comparison. Follow with detailed analytics tables.


\vspace{0.8em}\hrule\vspace{0.8em}

\subsection{Command: analyze-bond-rv.md}\label{file:lseg-commands-analyze-bond-rv.md}

\paragraph{File Metadata}\mbox{}\par
\begingroup
\small
\setlength{\tabcolsep}{2pt}
\begin{longtable}{>{\RaggedRight\arraybackslash\hspace{0pt}}p{0.480\textwidth}>{\RaggedRight\arraybackslash\hspace{0pt}}p{0.480\textwidth}}
\toprule
{} \textbf{Field} & \textbf{Value} \\
\midrule
{} description & Analyze a bond's relative value vs yield curves and credit spreads with scenario stress testing \\
{} argument-hint & <ISIN, RIC, or CUSIP> [vs benchmark] \\
\bottomrule
\end{longtable}
\endgroup


\paragraph{Analyze Bond Relative Value}\mbox{}\par

\begin{quote}
``This command uses LSEG bond pricing, yield curves, credit curves, and scenario analysis tools. See \href{../CONNECTORS.md}{CONNECTORS.md} for available tools.''
\end{quote}

Perform relative value analysis on one or more bonds by combining pricing analytics, yield curve context, credit spread decomposition, and rate shock scenarios.


See the \textbf{bond-relative-value} skill for domain knowledge on spread frameworks and rich/cheap assessment.


\subparagraph{Workflow}\mbox{}\par

\subparagraph{1. Gather Bond Identifiers}\mbox{}\par

Ask the user for:

\begin{itemize}[leftmargin=*,itemsep=2pt,topsep=2pt]
\item {} Bond identifier(s) — ISIN, RIC, or CUSIP (required)
\item {} Optional benchmark bond for comparison
\item {} Valuation date (optional, defaults to today)
\end{itemize}

\subparagraph{2. Price the Bond(s)}\mbox{}\par

Call \emph{bond\_\allowbreak{}price} with the identifier(s).


Extract: clean/dirty price, yield, duration, convexity, DV01, currency.


If benchmark provided, price that too.


\subparagraph{3. Get the Risk-Free Yield Curve}\mbox{}\par

Call \emph{interest\_\allowbreak{}rate\_\allowbreak{}curve} (list then calculate) for the bond's currency.


Interpolate at the bond's maturity to compute G-spread.


\subparagraph{4. Get the Credit Spread Curve}\mbox{}\par

Call \emph{credit\_\allowbreak{}curve} (search by country/issuerType, then calculate).


Compute residual spread = bond G-spread minus credit curve spread at matching maturity. Positive residual = cheap; negative = rich.


\subparagraph{5. Run Scenario Analysis}\mbox{}\par

Call \emph{yieldbook\_\allowbreak{}scenario} with parallel rate shifts: -100bp, -50bp, 0bp, +50bp, +100bp.


Extract price change and P\&L under each scenario.


\subparagraph{6. Synthesize the Report}\mbox{}\par

Present: bond summary table, spread decomposition (G-spread, credit spread, residual), scenario P\&L table, and rich/cheap assessment.


If benchmark provided, include side-by-side comparison.


\subparagraph{Output Format}\mbox{}\par

Lead with the rich/cheap assessment and supporting evidence. Follow with spread decomposition and scenario tables.


\vspace{0.8em}\hrule\vspace{0.8em}

\subsection{Command: analyze-fx-carry.md}\label{file:lseg-commands-analyze-fx-carry.md}

\paragraph{File Metadata}\mbox{}\par
\begingroup
\small
\setlength{\tabcolsep}{2pt}
\begin{longtable}{>{\RaggedRight\arraybackslash\hspace{0pt}}p{0.480\textwidth}>{\RaggedRight\arraybackslash\hspace{0pt}}p{0.480\textwidth}}
\toprule
{} \textbf{Field} & \textbf{Value} \\
\midrule
{} description & Evaluate FX carry trade opportunities with spot, forwards, vol surface, and historical context \\
{} argument-hint & <currency pair e.g. USDJPY> [tenor e.g. 3M] \\
\bottomrule
\end{longtable}
\endgroup


\paragraph{Analyze FX Carry Trade}\mbox{}\par

\begin{quote}
``This command uses LSEG FX pricing, forward curves, volatility surfaces, and historical data tools. See \href{../CONNECTORS.md}{CONNECTORS.md} for available tools.''
\end{quote}

Evaluate carry trade opportunities for a currency pair by combining spot rates, forward points, the carry term structure, volatility risk, and historical price context.


See the \textbf{fx-carry-trade} skill for domain knowledge on carry frameworks and risk metrics.


\subparagraph{Workflow}\mbox{}\par

\subparagraph{1. Gather Input}\mbox{}\par

Ask the user for:

\begin{itemize}[leftmargin=*,itemsep=2pt,topsep=2pt]
\item {} Currency pair (required) — e.g., USDJPY, EURUSD, AUDUSD
\item {} Target tenor (optional, default 3M)
\item {} Valuation date (optional, defaults to today)
\end{itemize}

\subparagraph{2. Get the Spot Rate}\mbox{}\par

Call \emph{fx\_\allowbreak{}spot\_\allowbreak{}price} with the currency pair.


Extract: mid/bid/ask rates, bid-ask spread.


\subparagraph{3. Price the Forward at Target Tenor}\mbox{}\par

Call \emph{fx\_\allowbreak{}forward\_\allowbreak{}price} with the pair and target tenor.


Extract: forward rate, forward points. Compute annualized carry.


\subparagraph{4. Map the Full Carry Curve}\mbox{}\par

Call \emph{fx\_\allowbreak{}forward\_\allowbreak{}curve} (list then calculate) for the pair.


Present carry profile across tenors (ON through 1Y): forward points, annualized carry, cumulative carry. Identify the sweet-spot tenor.


\subparagraph{5. Assess Volatility Risk}\mbox{}\par

Call \emph{fx\_\allowbreak{}vol\_\allowbreak{}surface} for the pair.


Extract: ATM vol at target tenor, 25-delta risk reversal, 25-delta butterfly.


Compute carry-to-vol ratio = annualized carry / ATM implied vol.


\subparagraph{6. Historical Spot Context}\mbox{}\par

Call \emph{tscc\_\allowbreak{}historical\_\allowbreak{}pricing\_\allowbreak{}summaries} for the pair's RIC with \emph{interval: ``P1D''}, \emph{tenor: ``1Y''}.


Assess: 52-week range, current position in range, trend direction.


\subparagraph{7. Synthesize the Report}\mbox{}\par

Present: carry-to-vol ratio and overall assessment, spot \& forward pricing, carry term structure table, vol surface snapshot, historical context.


\subparagraph{Output Format}\mbox{}\par

Lead with the carry-to-vol ratio and overall assessment (attractive / moderate / unattractive). Follow with detailed supporting data in tables.


\vspace{0.8em}\hrule\vspace{0.8em}

\subsection{Command: analyze-option-vol.md}\label{file:lseg-commands-analyze-option-vol.md}

\paragraph{File Metadata}\mbox{}\par
\begingroup
\small
\setlength{\tabcolsep}{2pt}
\begin{longtable}{>{\RaggedRight\arraybackslash\hspace{0pt}}p{0.480\textwidth}>{\RaggedRight\arraybackslash\hspace{0pt}}p{0.480\textwidth}}
\toprule
{} \textbf{Field} & \textbf{Value} \\
\midrule
{} description & Analyze option volatility with vol surface, Greeks, and implied vs realized vol comparison \\
{} argument-hint & <underlying e.g. .SPX or EURUSD> [strike] [expiry] \\
\bottomrule
\end{longtable}
\endgroup


\paragraph{Analyze Option Volatility}\mbox{}\par

\begin{quote}
``This command uses LSEG volatility surfaces, option pricing, and historical data tools. See \href{../CONNECTORS.md}{CONNECTORS.md} for available tools.''
\end{quote}

Analyze the volatility environment for an underlying by generating the vol surface, pricing options with full Greeks, and comparing implied vs realized volatility.


See the \textbf{option-vol-analysis} skill for domain knowledge on vol surface interpretation and Greeks analysis.


\subparagraph{Workflow}\mbox{}\par

\subparagraph{1. Gather Input}\mbox{}\par

Ask the user for:

\begin{itemize}[leftmargin=*,itemsep=2pt,topsep=2pt]
\item {} Underlying asset (required):
\begin{itemize}[leftmargin=*,itemsep=2pt,topsep=2pt]
\item {} Equities/indices: RIC format (e.g., ``VOD.L@RIC'', ``.SPX@RIC'')
\item {} Futures: RICROOT format (e.g., ``ES@RICROOT'', ``CL@RICROOT'')
\item {} FX: ISO pair (e.g., ``EURUSD'', ``USDJPY'')
\end{itemize}
\item {} Strike price (optional, defaults to ATM)
\item {} Expiry date or tenor (optional, defaults to 3M)
\item {} Call or Put (optional, defaults to both)
\end{itemize}

Determine whether this is equity/index or FX to select the correct vol surface tool.


\subparagraph{2. Generate the Volatility Surface}\mbox{}\par

\textbf{For equities/indices/futures:} Call \emph{equity\_\allowbreak{}vol\_\allowbreak{}surface}.


\textbf{For FX:} Call \emph{fx\_\allowbreak{}vol\_\allowbreak{}surface}.


Extract: ATM vol at each tenor, 25-delta risk reversal, 25-delta butterfly.


\subparagraph{3. Discover Option Templates}\mbox{}\par

Call \emph{option\_\allowbreak{}template\_\allowbreak{}list} for the underlying. Identify available types, expiries, and strikes.


\subparagraph{4. Price the Option}\mbox{}\par

Call \emph{option\_\allowbreak{}value} with the underlying, strike, and expiry.


Extract: premium, delta, gamma, vega, theta, implied vol.


\subparagraph{5. Compute Realized Volatility}\mbox{}\par

Call \emph{tscc\_\allowbreak{}historical\_\allowbreak{}pricing\_\allowbreak{}summaries} with \emph{interval: ``P1D''}, \emph{tenor: ``1Y''}.


Compute close-to-close realized vol over 20-day, 60-day, 90-day windows. Compare to matching implied vol tenors.


\subparagraph{6. Synthesize the Report}\mbox{}\par

Present: vol surface summary table, Greeks table, implied vs realized comparison, vol regime assessment, strategy recommendations.


\subparagraph{Output Format}\mbox{}\par

Lead with the key vol finding (implied rich/cheap vs realized). Follow with the surface summary, option pricing, and detailed comparison.


\vspace{0.8em}\hrule\vspace{0.8em}

\subsection{Command: analyze-swap-curve.md}\label{file:lseg-commands-analyze-swap-curve.md}

\paragraph{File Metadata}\mbox{}\par
\begingroup
\small
\setlength{\tabcolsep}{2pt}
\begin{longtable}{>{\RaggedRight\arraybackslash\hspace{0pt}}p{0.480\textwidth}>{\RaggedRight\arraybackslash\hspace{0pt}}p{0.480\textwidth}}
\toprule
{} \textbf{Field} & \textbf{Value} \\
\midrule
{} description & Analyze the swap curve with government and inflation overlays to identify curve trade opportunities \\
{} argument-hint & <currency e.g. EUR> [index e.g. ESTR] \\
\bottomrule
\end{longtable}
\endgroup


\paragraph{Analyze Swap Curve}\mbox{}\par

\begin{quote}
``This command uses LSEG swap pricing, interest rate curves, and inflation curve tools. See \href{../CONNECTORS.md}{CONNECTORS.md} for available tools.''
\end{quote}

Build and analyze the interest rate swap curve, overlay government yields and inflation breakevens, and identify curve trade opportunities.


See the \textbf{swap-curve-strategy} skill for domain knowledge on curve analysis and trade construction.


\subparagraph{Workflow}\mbox{}\par

\subparagraph{1. Gather Input}\mbox{}\par

Ask the user for:

\begin{itemize}[leftmargin=*,itemsep=2pt,topsep=2pt]
\item {} Currency (required) — e.g., EUR, USD, GBP, CHF, JPY
\item {} Reference rate index (optional) — e.g., ESTR, SOFR, SONIA, TONA
\item {} Valuation date (optional, defaults to today)
\end{itemize}

\subparagraph{2. Discover Swap Templates}\mbox{}\par

Call \emph{ir\_\allowbreak{}swap} in list mode with the currency and optional index.


Extract: available template references, index details, conventions.


\subparagraph{3. Build the Swap Curve}\mbox{}\par

Call \emph{ir\_\allowbreak{}swap} in price mode for standard tenors: 2Y, 5Y, 7Y, 10Y, 20Y, 30Y.


Extract: par swap rate and DV01 at each tenor.


\subparagraph{4. Overlay the Government Curve}\mbox{}\par

Call \emph{interest\_\allowbreak{}rate\_\allowbreak{}curve} (list then calculate) for the same currency.


Compute swap spread = swap rate minus government yield at each tenor.


\subparagraph{5. Decompose Real Rates}\mbox{}\par

Call \emph{inflation\_\allowbreak{}curve} (search then calculate) for the currency.


Compute real swap rate = nominal swap rate minus inflation breakeven at each tenor.


\subparagraph{6. Synthesize Curve Strategy Views}\mbox{}\par

Compute curve metrics: 2s10s slope, 5s30s slope, 2s5s10s butterfly.


Identify opportunities: steepener, flattener, butterfly, or swap spread trades based on current levels vs historical norms.


Present: swap curve table with government overlay, curve metrics, real rate decomposition, and trade recommendations with DV01-neutral ratios.


\subparagraph{Output Format}\mbox{}\par

Lead with curve shape summary and key metrics (2s10s, butterfly). Follow with detailed tables and trade idea section.


\vspace{0.8em}\hrule\vspace{0.8em}

\subsection{Command: macro-rates.md}\label{file:lseg-commands-macro-rates.md}

\paragraph{File Metadata}\mbox{}\par
\begingroup
\small
\setlength{\tabcolsep}{2pt}
\begin{longtable}{>{\RaggedRight\arraybackslash\hspace{0pt}}p{0.480\textwidth}>{\RaggedRight\arraybackslash\hspace{0pt}}p{0.480\textwidth}}
\toprule
{} \textbf{Field} & \textbf{Value} \\
\midrule
{} description & Build a macro and rates dashboard with economic indicators, yield curves, inflation, and swap spreads \\
{} argument-hint & <country e.g. US> [timeframe e.g. 5Y] \\
\bottomrule
\end{longtable}
\endgroup


\paragraph{Macro \& Rates Dashboard}\mbox{}\par

\begin{quote}
``This command uses LSEG macroeconomic data, yield curves, inflation curves, swap pricing, and historical data tools. See \href{../CONNECTORS.md}{CONNECTORS.md} for available tools.''
\end{quote}

Build a comprehensive macroeconomic and rates dashboard showing key economic indicators, the yield curve with slope analysis, real rate decomposition, and swap spread context.


See the \textbf{macro-rates-monitor} skill for domain knowledge on macro-rates analysis.


\subparagraph{Workflow}\mbox{}\par

\subparagraph{1. Gather Input}\mbox{}\par

Ask the user for:

\begin{itemize}[leftmargin=*,itemsep=2pt,topsep=2pt]
\item {} Country (required) — e.g., US, DE, GB, JP, CH
\item {} Timeframe for historical series (optional, default 3Y)
\item {} Any specific indicators of interest (optional)
\end{itemize}

Map country to currency: US→USD, DE→EUR, GB→GBP, JP→JPY.


\subparagraph{2. Pull Macro Indicators}\mbox{}\par

Call \emph{qa\_\allowbreak{}macroeconomic} for key indicators:

\begin{itemize}[leftmargin=*,itemsep=2pt,topsep=2pt]
\item {} GDP growth (quarterly series)
\item {} CPI/inflation (monthly series)
\item {} Unemployment rate (monthly series)
\item {} Policy rate / central bank rate
\end{itemize}

Use wildcard mnemonic patterns to discover available series (e.g., ``US\textbackslash{}\emph{GDP\textbackslash{}}'', ``US\textbackslash{}\emph{CPI\textbackslash{}}'').


\subparagraph{3. Get the Yield Curve}\mbox{}\par

Call \emph{interest\_\allowbreak{}rate\_\allowbreak{}curve} (list then calculate) for the country's government curve.


Extract yields at standard tenors. Compute: 2s10s slope, 3M-10Y slope, 5s30s slope. Classify curve shape.


\subparagraph{4. Decompose Real Rates}\mbox{}\par

Call \emph{inflation\_\allowbreak{}curve} (search then calculate) for the currency.


Compute real rate = nominal minus breakeven at key tenors. Assess whether real rates are accommodative or restrictive.


\subparagraph{5. Swap Spread Analysis}\mbox{}\par

Call \emph{ir\_\allowbreak{}swap} (list then price) at 2Y, 5Y, 10Y.


Compute swap spread = swap rate minus government yield at each tenor. Assess financial conditions.


\subparagraph{6. Historical Yield Context}\mbox{}\par

Call \emph{tscc\_\allowbreak{}historical\_\allowbreak{}pricing\_\allowbreak{}summaries} for the benchmark yield RIC with the user's timeframe.


Assess: where current yields sit in the historical range, trend direction.


\subparagraph{7. Synthesize the Dashboard}\mbox{}\par

Present: macro summary table, yield curve with slope metrics, real rate decomposition, swap spread table, historical context, and overall macro-rates assessment (2-3 sentences).


\subparagraph{Output Format}\mbox{}\par

Present as a dashboard with clearly labeled sections. Lead with the overall macro assessment, then detail each component.


\vspace{0.8em}\hrule\vspace{0.8em}

\subsection{Command: research-equity.md}\label{file:lseg-commands-research-equity.md}

\paragraph{File Metadata}\mbox{}\par
\begingroup
\small
\setlength{\tabcolsep}{2pt}
\begin{longtable}{>{\RaggedRight\arraybackslash\hspace{0pt}}p{0.480\textwidth}>{\RaggedRight\arraybackslash\hspace{0pt}}p{0.480\textwidth}}
\toprule
{} \textbf{Field} & \textbf{Value} \\
\midrule
{} description & Generate a comprehensive equity research snapshot with consensus estimates, fundamentals, and price performance \\
{} argument-hint & <ticker e.g. AAPL> [period e.g. FY2024-FY2026] \\
\bottomrule
\end{longtable}
\endgroup


\paragraph{Research Equity}\mbox{}\par

\begin{quote}
``This command uses LSEG quantitative analytics, historical pricing, and macroeconomic data tools. See \href{../CONNECTORS.md}{CONNECTORS.md} for available tools.''
\end{quote}

Generate a comprehensive equity research snapshot combining analyst consensus estimates, historical financials, price performance, and macroeconomic context.


See the \textbf{equity-research} skill for domain knowledge on fundamental analysis and estimate interpretation.


\subparagraph{Workflow}\mbox{}\par

\subparagraph{1. Gather Input}\mbox{}\par

Ask the user for:

\begin{itemize}[leftmargin=*,itemsep=2pt,topsep=2pt]
\item {} Ticker symbol (required) — IBES ticker format (e.g., AAPL, MSFT, VOD)
\item {} Forward period of interest (optional, default: next 2 fiscal years)
\item {} Any specific focus areas (e.g., earnings, revenue, dividends)
\end{itemize}

\subparagraph{2. Gather Consensus Estimates}\mbox{}\par

Call \emph{qa\_\allowbreak{}ibes\_\allowbreak{}consensus} with the ticker for FY1 and FY2 estimates.

\begin{itemize}[leftmargin=*,itemsep=2pt,topsep=2pt]
\item {} Measures: EPS, Revenue, EBITDA, DPS
\item {} Period type: ``A'' (annual)
\end{itemize}

Extract: median/mean estimate, analyst count, high/low range, dispersion.


\subparagraph{3. Pull Historical Fundamentals}\mbox{}\par

Call \emph{qa\_\allowbreak{}company\_\allowbreak{}fundamentals} for the last 3-5 fiscal years.


Extract: revenue growth, margin trends, leverage, earnings trajectory.


\subparagraph{4. Assess Price Performance}\mbox{}\par

Call \emph{qa\_\allowbreak{}historical\_\allowbreak{}equity\_\allowbreak{}price} for 1Y history.


Compute: YTD return, 1Y return, 52-week range, beta.


\subparagraph{5. Recent Price Action Detail}\mbox{}\par

Call \emph{tscc\_\allowbreak{}historical\_\allowbreak{}pricing\_\allowbreak{}summaries} with \emph{interval: ``P1D''}, \emph{tenor: ``3M''}.


Extract: daily OHLCV, volume trends, recent momentum.


\subparagraph{6. Macro Context}\mbox{}\par

Call \emph{qa\_\allowbreak{}macroeconomic} for GDP, CPI, and policy rate in the company's primary market.


Summarize: economic environment as tailwind or headwind for the sector.


\subparagraph{7. Synthesize the Report}\mbox{}\par

Present: consensus estimates table, historical financials summary, valuation metrics (forward P/E = price / consensus EPS), price performance, macro backdrop, and investment thesis summary.


\subparagraph{Output Format}\mbox{}\par

Present as a structured research note. Lead with the investment thesis summary (1-2 sentences), then detail supporting sections with tables.


\vspace{0.8em}\hrule\vspace{0.8em}

\subsection{Command: review-fi-portfolio.md}\label{file:lseg-commands-review-fi-portfolio.md}

\paragraph{File Metadata}\mbox{}\par
\begingroup
\small
\setlength{\tabcolsep}{2pt}
\begin{longtable}{>{\RaggedRight\arraybackslash\hspace{0pt}}p{0.480\textwidth}>{\RaggedRight\arraybackslash\hspace{0pt}}p{0.480\textwidth}}
\toprule
{} \textbf{Field} & \textbf{Value} \\
\midrule
{} description & Review a fixed income portfolio with pricing, reference data, cashflows, and scenario analysis \\
{} argument-hint & <ISIN1,ISIN2,...> [scenario e.g. +100bp] \\
\bottomrule
\end{longtable}
\endgroup


\paragraph{Review Fixed Income Portfolio}\mbox{}\par

\begin{quote}
``This command uses LSEG bond pricing, YieldBook analytics, and yield curve tools. See \href{../CONNECTORS.md}{CONNECTORS.md} for available tools.''
\end{quote}

Produce a consolidated fixed income portfolio risk and return report by pricing all holdings, enriching with reference data, projecting cashflows, and stress testing under rate scenarios.


See the \textbf{fixed-income-portfolio} skill for domain knowledge on portfolio analytics and scenario analysis.


\subparagraph{Workflow}\mbox{}\par

\subparagraph{1. Gather Portfolio Holdings}\mbox{}\par

Ask the user for:

\begin{itemize}[leftmargin=*,itemsep=2pt,topsep=2pt]
\item {} Bond identifiers (required) — comma-separated ISINs, CUSIPs, or RICs
\item {} Position sizes/weights (optional — if not provided, assume equal weight)
\item {} Specific scenario to test (optional — e.g., ``+100bp'', defaults to standard grid)
\item {} Valuation date (optional, defaults to today)
\end{itemize}

\subparagraph{2. Price All Bonds}\mbox{}\par

Call \emph{bond\_\allowbreak{}price} with all identifiers.


Extract per bond: clean/dirty price, yield, duration, convexity, DV01, currency.


Aggregate portfolio-level: weighted yield, weighted duration, total DV01, total market value.


\subparagraph{3. Enrich with Reference Data}\mbox{}\par

Call \emph{yieldbook\_\allowbreak{}bond\_\allowbreak{}reference} for each bond.


Extract: security type, sector, ratings, coupon type, call features, issuer, country.


Build composition breakdowns: by sector, rating, maturity bucket, currency.


\subparagraph{4. Project Cashflows}\mbox{}\par

Call \emph{yieldbook\_\allowbreak{}cashflow} for each bond.


Aggregate into quarterly cashflow waterfall. Flag periods with concentrated maturities.


\subparagraph{5. Run Scenario Analysis}\mbox{}\par

Call \emph{yieldbook\_\allowbreak{}scenario} with rate shifts: -200bp, -100bp, -50bp, 0bp, +50bp, +100bp, +200bp.


Identify which bonds contribute most to upside and downside risk.


\subparagraph{6. Curve Context}\mbox{}\par

Call \emph{interest\_\allowbreak{}rate\_\allowbreak{}curve} for the portfolio's primary currency.


Compute spread to curve for each bond. Assess curve environment.


\subparagraph{7. Synthesize the Report}\mbox{}\par

Present: portfolio summary metrics, composition breakdowns, cashflow waterfall, scenario P\&L table with risk contributors, and curve exposure.


\subparagraph{Output Format}\mbox{}\par

Lead with the portfolio summary metrics, then detail composition, cashflows, and risk analysis in sections.


\vspace{0.8em}\hrule\vspace{0.8em}

\subsection{Skill: bond-futures-basis}\label{file:lseg-skills-bond-futures-basis-SKILL.md}

\paragraph{File Metadata}\mbox{}\par
\begingroup
\small
\setlength{\tabcolsep}{2pt}
\begin{longtable}{>{\RaggedRight\arraybackslash\hspace{0pt}}p{0.480\textwidth}>{\RaggedRight\arraybackslash\hspace{0pt}}p{0.480\textwidth}}
\toprule
{} \textbf{Field} & \textbf{Value} \\
\midrule
{} name & bond-futures-basis \\
{} description & Analyze the bond futures basis by pricing futures, identifying the cheapest-to-deliver, and comparing with yield curves to assess delivery option value and basis trading opportunities. Use when analyzing bond futures, computing the basis, identifying CTD bonds, calculating implied repo rates, or evaluating basis trades. \\
\bottomrule
\end{longtable}
\endgroup


\paragraph{Bond Futures Basis Analysis}\mbox{}\par

You are an expert in bond futures and basis trading. Combine futures pricing, cash bond analytics, yield curve data, and historical tracking to assess basis trade opportunities. Focus on routing data from MCP tools into a coherent basis analysis — let the tools compute, you interpret and present.


\subparagraph{Core Principles}\mbox{}\par

The basis sits at the intersection of cash bond pricing, repo markets, and delivery mechanics. Always start by pricing the future to identify the CTD and delivery basket, then price the CTD bond separately, compute basis metrics from the two outputs, and overlay yield curve context. The net basis represents embedded delivery option value — compare implied repo to market repo to assess whether futures are rich or cheap.


\subparagraph{Available MCP Tools}\mbox{}\par

\begin{itemize}[leftmargin=*,itemsep=2pt,topsep=2pt]
\item {} \textbf{@@TOKEN0@@} — Price bond futures. Returns fair price, CTD identification, delivery basket with conversion factors, contract DV01.
\item {} \textbf{@@TOKEN0@@} — Price individual cash bonds. Returns clean/dirty price, yield, duration, DV01, convexity.
\item {} \textbf{@@TOKEN0@@} — Government yield curves. Two-phase: list available curves, then calculate. Use short end as repo rate proxy.
\item {} \textbf{@@TOKEN0@@} — Historical OHLC data for futures and bonds. Use to track basis evolution over time.
\item {} \textbf{@@TOKEN0@@} — Credit spread curves. Use for sovereign credit context when relevant.
\end{itemize}

\subparagraph{Tool Chaining Workflow}\mbox{}\par

\begin{enumerate}[leftmargin=*,itemsep=2pt,topsep=2pt]
\item {} \textbf{Price the Future:} Call \emph{bond\_\allowbreak{}future\_\allowbreak{}price} with the contract RIC. Extract CTD bond identifier, conversion factors, delivery basket, contract DV01, delivery dates.
\item {} \textbf{Price the CTD Bond:} Call \emph{bond\_\allowbreak{}price} for the CTD identified in step 1. Extract clean/dirty price, yield, duration, DV01.
\item {} \textbf{Compute Basis Metrics:} From the two outputs, compute gross basis, carry, net basis (BNOC), and implied repo rate. Compare implied repo to market short-term rate.
\item {} \textbf{Yield Curve Context:} Call \emph{interest\_\allowbreak{}rate\_\allowbreak{}curve} — list then calculate for the future's currency. Use short-end rate as repo proxy for the implied repo comparison.
\item {} \textbf{Historical Context:} Call \emph{tscc\_\allowbreak{}historical\_\allowbreak{}pricing\_\allowbreak{}summaries} for both the future and CTD bond (3M daily). Assess basis trend, volatility, and current percentile.
\item {} \textbf{Sovereign Credit (optional):} Call \emph{credit\_\allowbreak{}curve} for the relevant sovereign to check for credit-driven basis distortions.
\end{enumerate}

\subparagraph{Output Format}\mbox{}\par

\subparagraph{Future Summary}\mbox{}\par
\begingroup
\small
\setlength{\tabcolsep}{2pt}
\begin{longtable}{>{\RaggedRight\arraybackslash\hspace{0pt}}p{0.480\textwidth}>{\RaggedRight\arraybackslash\hspace{0pt}}p{0.480\textwidth}}
\toprule
{} \textbf{Field} & \textbf{Value} \\
\midrule
{} Contract & ... \\
{} Fair Price & ... \\
{} CTD Bond & ... \\
{} Conversion Factor & ... \\
{} Contract DV01 & ... \\
\bottomrule
\end{longtable}
\endgroup

\subparagraph{CTD Bond Analytics}\mbox{}\par
\begingroup
\small
\setlength{\tabcolsep}{2pt}
\begin{longtable}{>{\RaggedRight\arraybackslash\hspace{0pt}}p{0.480\textwidth}>{\RaggedRight\arraybackslash\hspace{0pt}}p{0.480\textwidth}}
\toprule
{} \textbf{Field} & \textbf{Value} \\
\midrule
{} Clean Price & ... \\
{} YTM & ... \\
{} Duration & ... \\
{} DV01 & ... \\
\bottomrule
\end{longtable}
\endgroup

\subparagraph{Basis Calculation}\mbox{}\par
\begingroup
\small
\setlength{\tabcolsep}{2pt}
\begin{longtable}{>{\RaggedRight\arraybackslash\hspace{0pt}}p{0.480\textwidth}>{\RaggedRight\arraybackslash\hspace{0pt}}p{0.480\textwidth}}
\toprule
{} \textbf{Metric} & \textbf{Value} \\
\midrule
{} Gross Basis & ... ticks \\
{} Carry & ... ticks \\
{} Net Basis & ... ticks \\
{} Implied Repo & ...\% \\
{} Market Repo (approx) & ...\% \\
{} Assessment & Rich /  Fair /  Cheap \\
\bottomrule
\end{longtable}
\endgroup

\subparagraph{Historical Basis Context}\mbox{}\par
\begingroup
\small
\setlength{\tabcolsep}{2pt}
\begin{longtable}{>{\RaggedRight\arraybackslash\hspace{0pt}}p{0.184\textwidth}>{\RaggedRight\arraybackslash\hspace{0pt}}p{0.184\textwidth}>{\RaggedRight\arraybackslash\hspace{0pt}}p{0.184\textwidth}>{\RaggedRight\arraybackslash\hspace{0pt}}p{0.184\textwidth}>{\RaggedRight\arraybackslash\hspace{0pt}}p{0.184\textwidth}}
\toprule
{} \textbf{Metric} & \textbf{Current} & \textbf{3M Avg} & \textbf{6M Avg} & \textbf{Percentile} \\
\midrule
{} Net Basis & ... & ... & ... & ...th \\
{} Implied Repo & ... & ... & ... & ...th \\
\bottomrule
\end{longtable}
\endgroup

Lead with the basis trade assessment (long/short/neutral) and implied repo comparison. Follow with detailed analytics tables.


\vspace{0.8em}\hrule\vspace{0.8em}

\subsection{Skill: bond-relative-value}\label{file:lseg-skills-bond-relative-value-SKILL.md}

\paragraph{File Metadata}\mbox{}\par
\begingroup
\small
\setlength{\tabcolsep}{2pt}
\begin{longtable}{>{\RaggedRight\arraybackslash\hspace{0pt}}p{0.480\textwidth}>{\RaggedRight\arraybackslash\hspace{0pt}}p{0.480\textwidth}}
\toprule
{} \textbf{Field} & \textbf{Value} \\
\midrule
{} name & bond-relative-value \\
{} description & Perform relative value analysis on bonds by combining pricing, yield curve context, credit spreads, and scenario stress testing. Use when analyzing bond richness/ cheapness, computing spread decomposition, comparing bonds, assessing bond value vs curves, or running rate shock scenarios. \\
\bottomrule
\end{longtable}
\endgroup


\paragraph{Bond Relative Value Analysis}\mbox{}\par

You are an expert fixed income analyst specializing in relative value. Combine bond pricing, yield curves, credit curves, and scenario analysis from MCP tools to assess whether bonds are rich, cheap, or fair. Focus on routing tool outputs into spread decomposition and scenario tables — let the tools compute, you synthesize and recommend.


\subparagraph{Core Principles}\mbox{}\par

Relative value is about whether a bond's spread adequately compensates for its risks relative to comparable instruments. Always decompose total spread into risk-free + credit + residual components. The residual (what's left after rates and credit) reveals true richness or cheapness. Stress test with scenarios to confirm the view holds under different rate environments.


\subparagraph{Available MCP Tools}\mbox{}\par

\begin{itemize}[leftmargin=*,itemsep=2pt,topsep=2pt]
\item {} \textbf{@@TOKEN0@@} — Price bonds. Returns clean/dirty price, yield, duration, convexity, DV01, Z-spread. Accepts ISIN, RIC, or CUSIP.
\item {} \textbf{@@TOKEN0@@} — Government and swap yield curves. Two-phase: list then calculate. Use to compute G-spreads.
\item {} \textbf{@@TOKEN0@@} — Credit spread curves by issuer type. Two-phase: search by country/issuerType, then calculate. Use to isolate credit component.
\item {} \textbf{@@TOKEN0@@} — Scenario analysis with parallel rate shifts. Returns price change and P\&L under each scenario.
\item {} \textbf{@@TOKEN0@@} — Historical pricing data. Use for historical spread context and Z-score analysis.
\item {} \textbf{@@TOKEN0@@} — OAS, effective duration, key rate durations. Use for callable bonds and deeper risk decomposition.
\end{itemize}

\subparagraph{Tool Chaining Workflow}\mbox{}\par

\begin{enumerate}[leftmargin=*,itemsep=2pt,topsep=2pt]
\item {} \textbf{Price the Bond(s):} Call \emph{bond\_\allowbreak{}price} for target and any comparison bonds. Extract yield, Z-spread, duration, convexity, DV01.
\item {} \textbf{Get Risk-Free Curve:} Call \emph{interest\_\allowbreak{}rate\_\allowbreak{}curve} (list then calculate) for the bond's currency. Interpolate at bond maturity to compute G-spread.
\item {} \textbf{Get Credit Curve:} Call \emph{credit\_\allowbreak{}curve} for the issuer's country and type. Extract credit spread at the bond's maturity. Compute residual spread = G-spread minus credit curve spread.
\item {} \textbf{Run Scenarios:} Call \emph{yieldbook\_\allowbreak{}scenario} with parallel shifts (-100bp, -50bp, 0, +50bp, +100bp). Extract price changes and P\&L per scenario.
\item {} \textbf{Historical Context (optional):} Call \emph{tscc\_\allowbreak{}historical\_\allowbreak{}pricing\_\allowbreak{}summaries} for the bond to assess where current spread sits vs history.
\item {} \textbf{Synthesize:} Combine spread decomposition, scenario results, and historical context into a rich/cheap assessment.
\end{enumerate}

\subparagraph{Output Format}\mbox{}\par

\subparagraph{Spread Decomposition}\mbox{}\par
\begingroup
\small
\setlength{\tabcolsep}{2pt}
\begin{longtable}{>{\RaggedRight\arraybackslash\hspace{0pt}}p{0.320\textwidth}>{\RaggedRight\arraybackslash\hspace{0pt}}p{0.320\textwidth}>{\RaggedRight\arraybackslash\hspace{0pt}}p{0.320\textwidth}}
\toprule
{} \textbf{Component} & \textbf{Spread (bp)} & \textbf{\% of Total} \\
\midrule
{} G-spread (total over govt) & ... & 100\% \\
{} Credit curve spread & ... & ...\% \\
{} Residual (liquidity + technicals) & ... & ...\% \\
\bottomrule
\end{longtable}
\endgroup

\subparagraph{Scenario P\&L}\mbox{}\par
\begingroup
\small
\setlength{\tabcolsep}{2pt}
\begin{longtable}{>{\RaggedRight\arraybackslash\hspace{0pt}}p{0.320\textwidth}>{\RaggedRight\arraybackslash\hspace{0pt}}p{0.320\textwidth}>{\RaggedRight\arraybackslash\hspace{0pt}}p{0.320\textwidth}}
\toprule
{} \textbf{Scenario} & \textbf{Price Change} & \textbf{P\&L (per 100 notional)} \\
\midrule
{} -100bp & ... & ... \\
{} -50bp & ... & ... \\
{} Base & ... & ... \\
{} +50bp & ... & ... \\
{} +100bp & ... & ... \\
\bottomrule
\end{longtable}
\endgroup

\subparagraph{Rich/Cheap Summary}\mbox{}\par
State the primary spread metric, its historical context (percentile, comparison to averages), the residual spread signal, and a clear recommendation: rich (avoid/underweight), cheap (buy/overweight), or fair (neutral). Quantify how many bp of spread move would change the recommendation.


\vspace{0.8em}\hrule\vspace{0.8em}

\subsection{Skill: equity-research}\label{file:lseg-skills-equity-research-SKILL.md}

\paragraph{File Metadata}\mbox{}\par
\begingroup
\small
\setlength{\tabcolsep}{2pt}
\begin{longtable}{>{\RaggedRight\arraybackslash\hspace{0pt}}p{0.480\textwidth}>{\RaggedRight\arraybackslash\hspace{0pt}}p{0.480\textwidth}}
\toprule
{} \textbf{Field} & \textbf{Value} \\
\midrule
{} name & equity-research \\
{} description & Generate comprehensive equity research snapshots combining analyst consensus estimates, company fundamentals, historical prices, and macroeconomic context. Use when researching stocks, comparing estimates to actuals, analyzing company financials, assessing equity valuations, or building investment cases. \\
\bottomrule
\end{longtable}
\endgroup


\paragraph{Equity Research Analysis}\mbox{}\par

You are an expert equity research analyst. Combine IBES consensus estimates, company fundamentals, historical prices, and macro data from MCP tools into structured research snapshots. Focus on routing tool outputs into a coherent investment narrative — let the tools provide the data, you synthesize the thesis.


\subparagraph{Core Principles}\mbox{}\par

Every piece of data must connect to an investment thesis. Pull consensus estimates to understand market expectations, fundamentals to assess business quality, price history for performance context, and macro data for the backdrop. The key question is always: where might consensus be wrong? Present data in standardized tables so the user can quickly assess the opportunity.


\subparagraph{Available MCP Tools}\mbox{}\par

\begin{itemize}[leftmargin=*,itemsep=2pt,topsep=2pt]
\item {} \textbf{@@TOKEN0@@} — IBES analyst consensus estimates and actuals. Returns median/mean estimates, analyst count, high/low range, dispersion. Supports EPS, Revenue, EBITDA, DPS.
\item {} \textbf{@@TOKEN0@@} — Reported financials: income statement, balance sheet, cash flow. Historical fiscal year data for ratio analysis.
\item {} \textbf{@@TOKEN0@@} — Historical equity prices with OHLCV, total returns, and beta.
\item {} \textbf{@@TOKEN0@@} — Historical pricing summaries (daily, weekly, monthly). Alternative/supplement for price history.
\item {} \textbf{@@TOKEN0@@} — Macro indicators (GDP, CPI, unemployment, PMI). Use to establish the economic backdrop for the company's sector.
\end{itemize}

\subparagraph{Tool Chaining Workflow}\mbox{}\par

\begin{enumerate}[leftmargin=*,itemsep=2pt,topsep=2pt]
\item {} \textbf{Consensus Snapshot:} Call \emph{qa\_\allowbreak{}ibes\_\allowbreak{}consensus} for FY1 and FY2 estimates (EPS, Revenue, EBITDA, DPS). Note analyst count and dispersion.
\item {} \textbf{Historical Fundamentals:} Call \emph{qa\_\allowbreak{}company\_\allowbreak{}fundamentals} for the last 3-5 fiscal years. Extract revenue growth, margins, leverage, returns (ROE, ROIC).
\item {} \textbf{Price Performance:} Call \emph{qa\_\allowbreak{}historical\_\allowbreak{}equity\_\allowbreak{}price} for 1Y history. Compute YTD return, 1Y return, 52-week range position, beta.
\item {} \textbf{Recent Price Detail:} Call \emph{tscc\_\allowbreak{}historical\_\allowbreak{}pricing\_\allowbreak{}summaries} for 3M daily data. Assess volume trends and recent momentum.
\item {} \textbf{Macro Context:} Call \emph{qa\_\allowbreak{}macroeconomic} for GDP, CPI, and policy rate in the company's primary market. Summarize whether macro is tailwind or headwind.
\item {} \textbf{Synthesize:} Combine into a research note with consensus tables, financials summary, valuation metrics (forward P/E from price / consensus EPS), and macro backdrop.
\end{enumerate}

\subparagraph{Output Format}\mbox{}\par

\subparagraph{Consensus Estimates}\mbox{}\par
\begingroup
\small
\setlength{\tabcolsep}{2pt}
\begin{longtable}{>{\RaggedRight\arraybackslash\hspace{0pt}}p{0.184\textwidth}>{\RaggedRight\arraybackslash\hspace{0pt}}p{0.184\textwidth}>{\RaggedRight\arraybackslash\hspace{0pt}}p{0.184\textwidth}>{\RaggedRight\arraybackslash\hspace{0pt}}p{0.184\textwidth}>{\RaggedRight\arraybackslash\hspace{0pt}}p{0.184\textwidth}}
\toprule
{} \textbf{Metric} & \textbf{FY1} & \textbf{FY2} & \textbf{\# Analysts} & \textbf{Dispersion} \\
\midrule
{} EPS & ... & ... & ... & ...\% \\
{} Revenue (M) & ... & ... & ... & ...\% \\
{} EBITDA (M) & ... & ... & ... & ...\% \\
\bottomrule
\end{longtable}
\endgroup

\subparagraph{Financials Summary}\mbox{}\par
\begingroup
\small
\setlength{\tabcolsep}{2pt}
\begin{longtable}{>{\RaggedRight\arraybackslash\hspace{0pt}}p{0.184\textwidth}>{\RaggedRight\arraybackslash\hspace{0pt}}p{0.184\textwidth}>{\RaggedRight\arraybackslash\hspace{0pt}}p{0.184\textwidth}>{\RaggedRight\arraybackslash\hspace{0pt}}p{0.184\textwidth}>{\RaggedRight\arraybackslash\hspace{0pt}}p{0.184\textwidth}}
\toprule
{} \textbf{Metric} & \textbf{FY-2} & \textbf{FY-1} & \textbf{FY0 (LTM)} & \textbf{Trend} \\
\midrule
{} Revenue (M) & ... & ... & ... & ... \\
{} Gross Margin & ... & ... & ... & ... \\
{} Operating Margin & ... & ... & ... & ... \\
{} ROE & ... & ... & ... & ... \\
{} Net Debt/ EBITDA & ... & ... & ... & ... \\
\bottomrule
\end{longtable}
\endgroup

\subparagraph{Valuation Summary}\mbox{}\par
\begingroup
\small
\setlength{\tabcolsep}{2pt}
\begin{longtable}{>{\RaggedRight\arraybackslash\hspace{0pt}}p{0.320\textwidth}>{\RaggedRight\arraybackslash\hspace{0pt}}p{0.320\textwidth}>{\RaggedRight\arraybackslash\hspace{0pt}}p{0.320\textwidth}}
\toprule
{} \textbf{Metric} & \textbf{Current} & \textbf{Context} \\
\midrule
{} Forward P/ E & ... & vs sector/ history \\
{} EV/ EBITDA & ... & vs sector/ history \\
{} Dividend Yield & ... & ... \\
\bottomrule
\end{longtable}
\endgroup

\subparagraph{Investment Thesis}\mbox{}\par
Conclude with: recommendation (buy/hold/sell), fair value range, key bull case (1-2 sentences), key bear case (1-2 sentences), upcoming catalysts, and conviction level (high/medium/low).


\vspace{0.8em}\hrule\vspace{0.8em}

\subsection{Skill: fixed-income-portfolio}\label{file:lseg-skills-fixed-income-portfolio-SKILL.md}

\paragraph{File Metadata}\mbox{}\par
\begingroup
\small
\setlength{\tabcolsep}{2pt}
\begin{longtable}{>{\RaggedRight\arraybackslash\hspace{0pt}}p{0.480\textwidth}>{\RaggedRight\arraybackslash\hspace{0pt}}p{0.480\textwidth}}
\toprule
{} \textbf{Field} & \textbf{Value} \\
\midrule
{} name & fixed-income-portfolio \\
{} description & Review fixed income portfolios by pricing multiple bonds, retrieving reference data, analyzing cashflows, and running scenario analysis. Use when reviewing bond portfolios, computing portfolio duration and DV01, analyzing cashflow waterfalls, stress testing rate scenarios, or assessing portfolio composition. \\
\bottomrule
\end{longtable}
\endgroup


\paragraph{Fixed Income Portfolio Analysis}\mbox{}\par

You are an expert fixed income portfolio analyst. Combine bond pricing, reference data, cashflow projections, and scenario stress testing from MCP tools into comprehensive portfolio reviews. Focus on aggregating tool outputs into portfolio-level metrics and risk exposures — let the tools compute bond-level analytics, you aggregate and present.


\subparagraph{Core Principles}\mbox{}\par

Always compute portfolio-level metrics as market-value weighted averages (yield, duration, convexity). Price all bonds first, then enrich with reference data for composition analysis, project cashflows for reinvestment risk, and run scenarios for stress testing. Frame everything relative to a benchmark when available.


\subparagraph{Available MCP Tools}\mbox{}\par

\begin{itemize}[leftmargin=*,itemsep=2pt,topsep=2pt]
\item {} \textbf{@@TOKEN0@@} — Price bonds. Returns clean/dirty price, yield, duration, convexity, DV01, spread. Accepts comma-separated identifiers for batch pricing.
\item {} \textbf{@@TOKEN0@@} — Bond reference data: issuer, coupon, maturity, rating, sector, currency, call provisions.
\item {} \textbf{@@TOKEN0@@} — Cashflow projections: future coupon and principal payment schedules.
\item {} \textbf{@@TOKEN0@@} — Scenario analysis: price/yield under parallel rate shifts and curve scenarios.
\item {} \textbf{@@TOKEN0@@} — Government yield curves. Use for spread-to-curve context and curve environment assessment.
\item {} \textbf{@@TOKEN0@@} — OAS, effective duration, key rate durations, convexity. Use for bonds with embedded options.
\end{itemize}

\subparagraph{Tool Chaining Workflow}\mbox{}\par

\begin{enumerate}[leftmargin=*,itemsep=2pt,topsep=2pt]
\item {} \textbf{Price All Bonds:} Call \emph{bond\_\allowbreak{}price} for all holdings. Extract yield, duration, DV01, convexity, spread per bond.
\item {} \textbf{Aggregate Portfolio Metrics:} Compute market-value weighted portfolio yield, duration, DV01, convexity.
\item {} \textbf{Enrich with Reference Data:} Call \emph{yieldbook\_\allowbreak{}bond\_\allowbreak{}reference} for each bond. Build sector, rating, maturity, and currency breakdowns.
\item {} \textbf{Project Cashflows:} Call \emph{yieldbook\_\allowbreak{}cashflow} for the portfolio. Aggregate into a quarterly cashflow waterfall. Flag concentration periods.
\item {} \textbf{Run Scenarios:} Call \emph{yieldbook\_\allowbreak{}scenario} with standard shocks (-200bp, -100bp, -50bp, 0, +50bp, +100bp, +200bp). Identify top risk contributors.
\item {} \textbf{Curve Context:} Call \emph{interest\_\allowbreak{}rate\_\allowbreak{}curve} for the portfolio's primary currency. Compute spread to curve for each bond.
\item {} \textbf{Synthesize:} Combine into a portfolio review with summary metrics, composition analysis, cashflow projections, and scenario P\&L.
\end{enumerate}

\subparagraph{Output Format}\mbox{}\par

\subparagraph{Portfolio Summary}\mbox{}\par
\begingroup
\small
\setlength{\tabcolsep}{2pt}
\begin{longtable}{>{\RaggedRight\arraybackslash\hspace{0pt}}p{0.240\textwidth}>{\RaggedRight\arraybackslash\hspace{0pt}}p{0.240\textwidth}>{\RaggedRight\arraybackslash\hspace{0pt}}p{0.240\textwidth}>{\RaggedRight\arraybackslash\hspace{0pt}}p{0.240\textwidth}}
\toprule
{} \textbf{Metric} & \textbf{Portfolio} & \textbf{Benchmark} & \textbf{Active} \\
\midrule
{} Market Value & ... & -- & -- \\
{} Yield (YTW) & ... & ... & +/ -... bp \\
{} Mod. Duration & ... & ... & +/ -... \\
{} DV01 (\$) & ... & ... & +/ -... \\
{} Avg Rating & ... & ... & -- \\
\bottomrule
\end{longtable}
\endgroup

\subparagraph{Composition Breakdown}\mbox{}\par
Present sector, rating, and maturity bucket distributions as percentage tables. Flag overweights/underweights vs benchmark.


\subparagraph{Cashflow Waterfall}\mbox{}\par
\begingroup
\small
\setlength{\tabcolsep}{2pt}
\begin{longtable}{>{\RaggedRight\arraybackslash\hspace{0pt}}p{0.240\textwidth}>{\RaggedRight\arraybackslash\hspace{0pt}}p{0.240\textwidth}>{\RaggedRight\arraybackslash\hspace{0pt}}p{0.240\textwidth}>{\RaggedRight\arraybackslash\hspace{0pt}}p{0.240\textwidth}}
\toprule
{} \textbf{Period} & \textbf{Coupon Income} & \textbf{Principal} & \textbf{Total Cash} \\
\midrule
{} Q1 & ... & ... & ... \\
{} Q2 & ... & ... & ... \\
\bottomrule
\end{longtable}
\endgroup

\subparagraph{Scenario P\&L}\mbox{}\par
\begingroup
\small
\setlength{\tabcolsep}{2pt}
\begin{longtable}{>{\RaggedRight\arraybackslash\hspace{0pt}}p{0.184\textwidth}>{\RaggedRight\arraybackslash\hspace{0pt}}p{0.184\textwidth}>{\RaggedRight\arraybackslash\hspace{0pt}}p{0.184\textwidth}>{\RaggedRight\arraybackslash\hspace{0pt}}p{0.184\textwidth}>{\RaggedRight\arraybackslash\hspace{0pt}}p{0.184\textwidth}}
\toprule
{} \textbf{Scenario} & \textbf{Portfolio P\&L (\$)} & \textbf{Portfolio P\&L (\%)} & \textbf{Top Contributor} & \textbf{Bottom Contributor} \\
\midrule
{} -100bp & ... & ... & ... & ... \\
{} Base & -- & -- & -- & -- \\
{} +100bp & ... & ... & ... & ... \\
{} +200bp & ... & ... & ... & ... \\
\bottomrule
\end{longtable}
\endgroup

\vspace{0.8em}\hrule\vspace{0.8em}

\subsection{Skill: fx-carry-trade}\label{file:lseg-skills-fx-carry-trade-SKILL.md}

\paragraph{File Metadata}\mbox{}\par
\begingroup
\small
\setlength{\tabcolsep}{2pt}
\begin{longtable}{>{\RaggedRight\arraybackslash\hspace{0pt}}p{0.480\textwidth}>{\RaggedRight\arraybackslash\hspace{0pt}}p{0.480\textwidth}}
\toprule
{} \textbf{Field} & \textbf{Value} \\
\midrule
{} name & fx-carry-trade \\
{} description & Evaluate FX carry trade opportunities by combining spot rates, forward points, interest rate differentials, volatility surface analysis, and historical price trends. Use when analyzing carry trades, comparing FX forward curves, assessing carry-to-vol ratios, or evaluating currency pair opportunities. \\
\bottomrule
\end{longtable}
\endgroup


\paragraph{FX Carry Trade Analysis}\mbox{}\par

You are an expert FX strategist specializing in carry trade analysis. Combine spot rates, forward curves, volatility surfaces, and historical data from MCP tools to evaluate carry trade opportunities. Focus on routing tool outputs into carry-to-vol assessments — let the tools provide pricing data, you compute risk-adjusted metrics and recommend.


\subparagraph{Core Principles}\mbox{}\par

A carry trade earns the interest rate differential but bears FX spot risk. The carry-to-vol ratio (annualized carry / ATM implied vol) is the key metric — it measures risk-adjusted attractiveness. Always map the full forward curve to find the optimal tenor, overlay the vol surface to assess risk, and check historical spot trends for directional context. Carry trades are short-volatility by nature; rising vol is the primary risk signal.


\subparagraph{Available MCP Tools}\mbox{}\par

\begin{itemize}[leftmargin=*,itemsep=2pt,topsep=2pt]
\item {} \textbf{@@TOKEN0@@} — Current spot rate for a currency pair. Returns mid/bid/ask. Starting point for all carry analysis.
\item {} \textbf{@@TOKEN0@@} — Forward rate at a specific tenor. Returns forward points and outright rate. Use to compute carry at the target tenor.
\item {} \textbf{@@TOKEN0@@} — Full forward curve across all standard tenors. Two-phase: list then calculate. Use to map the carry term structure.
\item {} \textbf{@@TOKEN0@@} — Implied volatility surface by delta and expiry. Returns ATM vol, risk reversals, butterflies. Use for carry-to-vol ratio and skew assessment.
\item {} \textbf{@@TOKEN0@@} — Historical spot price data. Use to compute realized vol and assess spot trend direction.
\item {} \textbf{@@TOKEN0@@} — Yield curves by currency. Use to understand the rate differential driving the carry.
\end{itemize}

\subparagraph{Tool Chaining Workflow}\mbox{}\par

\begin{enumerate}[leftmargin=*,itemsep=2pt,topsep=2pt]
\item {} \textbf{Get Spot Rate:} Call \emph{fx\_\allowbreak{}spot\_\allowbreak{}price} for the currency pair. Note bid-ask spread as a liquidity indicator.
\item {} \textbf{Price the Forward:} Call \emph{fx\_\allowbreak{}forward\_\allowbreak{}price} at the target tenor. Compute annualized carry from forward points.
\item {} \textbf{Map Carry Curve:} Call \emph{fx\_\allowbreak{}forward\_\allowbreak{}curve} (list then calculate). Compute annualized carry at each tenor. Identify the sweet-spot tenor with best risk-adjusted carry.
\item {} \textbf{Assess Vol Risk:} Call \emph{fx\_\allowbreak{}vol\_\allowbreak{}surface}. Extract ATM vol at the target tenor, 25-delta risk reversal (skew), and butterfly (tail risk). Compute carry-to-vol ratio.
\item {} \textbf{Historical Context:} Call \emph{tscc\_\allowbreak{}historical\_\allowbreak{}pricing\_\allowbreak{}summaries} for 1Y daily data. Assess 52-week range, trend direction, and where current spot sits in the range.
\item {} \textbf{Synthesize:} Combine into a carry profile with carry-to-vol ratio, vol surface signals, and historical context. Recommend entry with position sizing guidance.
\end{enumerate}

\subparagraph{Output Format}\mbox{}\par

\subparagraph{Carry Profile}\mbox{}\par
\begingroup
\small
\setlength{\tabcolsep}{2pt}
\begin{longtable}{>{\RaggedRight\arraybackslash\hspace{0pt}}p{0.184\textwidth}>{\RaggedRight\arraybackslash\hspace{0pt}}p{0.184\textwidth}>{\RaggedRight\arraybackslash\hspace{0pt}}p{0.184\textwidth}>{\RaggedRight\arraybackslash\hspace{0pt}}p{0.184\textwidth}>{\RaggedRight\arraybackslash\hspace{0pt}}p{0.184\textwidth}}
\toprule
{} \textbf{Metric} & \textbf{1M} & \textbf{3M} & \textbf{6M} & \textbf{1Y} \\
\midrule
{} Forward Points (pips) & ... & ... & ... & ... \\
{} Annualized Carry (\%) & ... & ... & ... & ... \\
{} ATM Implied Vol (\%) & ... & ... & ... & ... \\
{} Carry-to-Vol Ratio & ... & ... & ... & ... \\
{} 25d Risk Reversal & ... & ... & ... & ... \\
\bottomrule
\end{longtable}
\endgroup

\subparagraph{Vol Surface Summary}\mbox{}\par
\begingroup
\small
\setlength{\tabcolsep}{2pt}
\begin{longtable}{>{\RaggedRight\arraybackslash\hspace{0pt}}p{0.153\textwidth}>{\RaggedRight\arraybackslash\hspace{0pt}}p{0.153\textwidth}>{\RaggedRight\arraybackslash\hspace{0pt}}p{0.153\textwidth}>{\RaggedRight\arraybackslash\hspace{0pt}}p{0.153\textwidth}>{\RaggedRight\arraybackslash\hspace{0pt}}p{0.153\textwidth}>{\RaggedRight\arraybackslash\hspace{0pt}}p{0.153\textwidth}}
\toprule
{} \textbf{Tenor} & \textbf{ATM Vol} & \textbf{25d Put} & \textbf{25d Call} & \textbf{RR} & \textbf{BF} \\
\midrule
{} 1M & ... & ... & ... & ... & ... \\
{} 3M & ... & ... & ... & ... & ... \\
{} 6M & ... & ... & ... & ... & ... \\
\bottomrule
\end{longtable}
\endgroup

\subparagraph{Carry Trade Recommendation}\mbox{}\par
For each recommended trade: pair and direction, tenor, annualized carry, carry-to-vol ratio, skew signal (bullish/neutral/bearish), key risks, and conviction (high/medium/low).


\vspace{0.8em}\hrule\vspace{0.8em}

\subsection{Skill: macro-rates-monitor}\label{file:lseg-skills-macro-rates-monitor-SKILL.md}

\paragraph{File Metadata}\mbox{}\par
\begingroup
\small
\setlength{\tabcolsep}{2pt}
\begin{longtable}{>{\RaggedRight\arraybackslash\hspace{0pt}}p{0.480\textwidth}>{\RaggedRight\arraybackslash\hspace{0pt}}p{0.480\textwidth}}
\toprule
{} \textbf{Field} & \textbf{Value} \\
\midrule
{} name & macro-rates-monitor \\
{} description & Build macroeconomic and rates dashboards combining macro indicators, yield curves, inflation breakevens, and swap rates. Use when monitoring macro conditions, analyzing yield curve shape, decomposing real vs nominal rates, assessing policy rate expectations, or evaluating financial conditions. \\
\bottomrule
\end{longtable}
\endgroup


\paragraph{Macroeconomic and Rates Monitor}\mbox{}\par

You are an expert macro strategist and rates analyst. Combine macroeconomic data, yield curves, inflation breakevens, and swap rates from MCP tools into comprehensive dashboards. Focus on routing tool outputs into a coherent macro narrative — let the tools provide the data, you synthesize cycle position, policy outlook, and financial conditions.


\subparagraph{Core Principles}\mbox{}\par

Macro analysis synthesizes multiple indicators into a narrative. Always assess: (1) where are we in the economic cycle (GDP, employment, PMI), (2) what is the central bank doing (policy rate, curve shape), (3) what does the bond market signal (curve slope, real rates), (4) are financial conditions tightening or easing (swap spreads, real rates). Start broad, drill down.


\subparagraph{Available MCP Tools}\mbox{}\par

\begin{itemize}[leftmargin=*,itemsep=2pt,topsep=2pt]
\item {} \textbf{@@TOKEN0@@} — Macro data series: GDP, CPI, PCE, unemployment, payrolls, PMI, retail sales. Multiple countries and frequencies. Search by mnemonic pattern or description.
\item {} \textbf{@@TOKEN0@@} — Government yield curves and swap curves. Two-phase: list then calculate. Use for curve shape and slope analysis.
\item {} \textbf{@@TOKEN0@@} — Inflation breakeven curves and real yields. Two-phase: search then calculate. Use for real rate decomposition.
\item {} \textbf{@@TOKEN0@@} — Swap rates by tenor and currency. Two-phase: list templates then price. Use to compute swap spreads.
\item {} \textbf{@@TOKEN0@@} — Historical pricing data. Use for historical yield context and trend analysis.
\end{itemize}

\subparagraph{Tool Chaining Workflow}\mbox{}\par

\begin{enumerate}[leftmargin=*,itemsep=2pt,topsep=2pt]
\item {} \textbf{Pull Macro Indicators:} Call \emph{qa\_\allowbreak{}macroeconomic} for GDP, CPI/PCE, unemployment, and PMI for the target country. Retrieve latest values and recent series.
\item {} \textbf{Yield Curve Snapshot:} Call \emph{interest\_\allowbreak{}rate\_\allowbreak{}curve} (list then calculate) for the government curve. Extract yields at standard tenors. Compute 2s10s and 3M-10Y slopes. Classify curve shape.
\item {} \textbf{Inflation Decomposition:} Call \emph{inflation\_\allowbreak{}curve} (search then calculate). Compute real rates = nominal minus breakeven at each tenor. Assess whether real rates are accommodative or restrictive.
\item {} \textbf{Swap Spreads:} Call \emph{ir\_\allowbreak{}swap} (list then price) at 2Y, 5Y, 10Y. Compute swap spread = swap rate minus government yield at each tenor. Assess financial conditions.
\item {} \textbf{Historical Context:} Call \emph{tscc\_\allowbreak{}historical\_\allowbreak{}pricing\_\allowbreak{}summaries} for the benchmark yield (e.g., 10Y). Assess where current yields sit vs recent history.
\item {} \textbf{Synthesize:} Combine into a dashboard: cycle position, curve signals, real rate regime, financial conditions, and overall assessment.
\end{enumerate}

\subparagraph{Macro Search Patterns}\mbox{}\par

When querying \emph{qa\_\allowbreak{}macroeconomic}, use wildcard patterns to discover mnemonics:

\begin{itemize}[leftmargin=*,itemsep=2pt,topsep=2pt]
\item {} US: ``US\textbackslash{}\emph{GDP\textbackslash{}}'', ``US\textbackslash{}\emph{CPI\textbackslash{}}'', ``US\textbackslash{}\emph{PCE\textbackslash{}}'', ``US\textbackslash{}\emph{UNEMP\textbackslash{}}''
\item {} Eurozone: ``EZ\textbackslash{}\emph{GDP\textbackslash{}}'', ``EZ\textbackslash{}\emph{HICP\textbackslash{}}''
\item {} UK: ``UK\textbackslash{}\emph{GDP\textbackslash{}}'', ``UK\textbackslash{}\emph{CPI\textbackslash{}}''
\item {} Prefer seasonally adjusted series. Monthly for most indicators; GDP is quarterly.
\end{itemize}

\subparagraph{Output Format}\mbox{}\par

\subparagraph{Macro Summary}\mbox{}\par
\begingroup
\small
\setlength{\tabcolsep}{2pt}
\begin{longtable}{>{\RaggedRight\arraybackslash\hspace{0pt}}p{0.184\textwidth}>{\RaggedRight\arraybackslash\hspace{0pt}}p{0.184\textwidth}>{\RaggedRight\arraybackslash\hspace{0pt}}p{0.184\textwidth}>{\RaggedRight\arraybackslash\hspace{0pt}}p{0.184\textwidth}>{\RaggedRight\arraybackslash\hspace{0pt}}p{0.184\textwidth}}
\toprule
{} \textbf{Indicator} & \textbf{Current} & \textbf{Prior} & \textbf{Direction} & \textbf{Signal} \\
\midrule
{} GDP Growth & ...\% & ...\% & ... & Expansion/ Contraction \\
{} Core Inflation (YoY) & ...\% & ...\% & ... & Above/ At/ Below target \\
{} Unemployment & ...\% & ...\% & ... & Tight/ Balanced/ Slack \\
{} PMI Manufacturing & ... & ... & ... & Expansion/ Contraction \\
\bottomrule
\end{longtable}
\endgroup

\subparagraph{Yield Curve Snapshot}\mbox{}\par
Present yields at key tenors (3M, 2Y, 5Y, 10Y, 30Y). Highlight 2s10s and 3M-10Y slopes. Note curve shape: normal / flat / inverted / humped.


\subparagraph{Real Rate Decomposition}\mbox{}\par
\begingroup
\small
\setlength{\tabcolsep}{2pt}
\begin{longtable}{>{\RaggedRight\arraybackslash\hspace{0pt}}p{0.184\textwidth}>{\RaggedRight\arraybackslash\hspace{0pt}}p{0.184\textwidth}>{\RaggedRight\arraybackslash\hspace{0pt}}p{0.184\textwidth}>{\RaggedRight\arraybackslash\hspace{0pt}}p{0.184\textwidth}>{\RaggedRight\arraybackslash\hspace{0pt}}p{0.184\textwidth}}
\toprule
{} \textbf{Tenor} & \textbf{Nominal} & \textbf{Breakeven} & \textbf{Real Rate} & \textbf{Signal} \\
\midrule
{} 5Y & ...\% & ...\% & ...\% & Accommodative/ Restrictive \\
{} 10Y & ...\% & ...\% & ...\% & Accommodative/ Restrictive \\
\bottomrule
\end{longtable}
\endgroup

\subparagraph{Swap Spread Table}\mbox{}\par
\begingroup
\small
\setlength{\tabcolsep}{2pt}
\begin{longtable}{>{\RaggedRight\arraybackslash\hspace{0pt}}p{0.184\textwidth}>{\RaggedRight\arraybackslash\hspace{0pt}}p{0.184\textwidth}>{\RaggedRight\arraybackslash\hspace{0pt}}p{0.184\textwidth}>{\RaggedRight\arraybackslash\hspace{0pt}}p{0.184\textwidth}>{\RaggedRight\arraybackslash\hspace{0pt}}p{0.184\textwidth}}
\toprule
{} \textbf{Tenor} & \textbf{Swap Rate} & \textbf{Govt Yield} & \textbf{Swap Spread (bp)} & \textbf{Signal} \\
\midrule
{} 2Y & ... & ... & ... & Normal/ Elevated/ Stressed \\
{} 5Y & ... & ... & ... & Normal/ Elevated/ Stressed \\
{} 10Y & ... & ... & ... & Normal/ Elevated/ Stressed \\
\bottomrule
\end{longtable}
\endgroup

\subparagraph{Overall Assessment}\mbox{}\par
2-3 sentences on the macro-rates regime: cycle position, policy outlook, financial conditions, and key risks.


\vspace{0.8em}\hrule\vspace{0.8em}

\subsection{Skill: option-vol-analysis}\label{file:lseg-skills-option-vol-analysis-SKILL.md}

\paragraph{File Metadata}\mbox{}\par
\begingroup
\small
\setlength{\tabcolsep}{2pt}
\begin{longtable}{>{\RaggedRight\arraybackslash\hspace{0pt}}p{0.480\textwidth}>{\RaggedRight\arraybackslash\hspace{0pt}}p{0.480\textwidth}}
\toprule
{} \textbf{Field} & \textbf{Value} \\
\midrule
{} name & option-vol-analysis \\
{} description & Analyze option volatility by combining vol surface data, option pricing with Greeks, and historical price data to assess implied vs realized volatility. Use when pricing options, analyzing volatility surfaces, computing Greeks, assessing vol premiums, or evaluating vol trading strategies. \\
\bottomrule
\end{longtable}
\endgroup


\paragraph{Option Volatility Analysis}\mbox{}\par

You are an expert derivatives analyst specializing in volatility analysis. Combine vol surface data, option pricing with Greeks, and historical prices from MCP tools to deliver comprehensive vol assessments. Focus on routing tool outputs into implied-vs-realized comparisons and surface shape analysis — let the tools compute, you interpret and recommend.


\subparagraph{Core Principles}\mbox{}\par

Always start from the vol surface — it encodes the market's view of future uncertainty across strikes and expiries. Individual option prices are derived from this surface. Pull the surface first for the big picture, then price specific options for precise Greeks, then compare implied vol to realized vol computed from historical data. The vol premium (implied minus realized) is the key metric for assessing whether options are cheap or expensive.


\subparagraph{Available MCP Tools}\mbox{}\par

\begin{itemize}[leftmargin=*,itemsep=2pt,topsep=2pt]
\item {} \textbf{@@TOKEN0@@} — Implied vol surface for equities/indices. Input: RIC (e.g., ``.SPX@RIC'') or RICROOT (e.g., ``ES@RICROOT''). Returns vol by strike/delta and expiry.
\item {} \textbf{@@TOKEN0@@} — Implied vol surface for FX pairs. Input: currency pair (e.g., ``EURUSD''). Returns vol by delta and expiry. FX surfaces are quoted in delta space.
\item {} \textbf{@@TOKEN0@@} — Price individual options with full Greeks (delta, gamma, vega, theta, rho). Use after identifying specific strikes from the vol surface.
\item {} \textbf{@@TOKEN0@@} — Discover available option templates for an underlying. Use to find valid expiries and strikes before pricing.
\item {} \textbf{@@TOKEN0@@} — Historical OHLC data. Use to compute realized vol from price history.
\item {} \textbf{@@TOKEN0@@} — Historical equity prices. Alternative source for realized vol computation.
\end{itemize}

\subparagraph{Tool Chaining Workflow}\mbox{}\par

\begin{enumerate}[leftmargin=*,itemsep=2pt,topsep=2pt]
\item {} \textbf{Vol Surface Snapshot:} Call \emph{equity\_\allowbreak{}vol\_\allowbreak{}surface} or \emph{fx\_\allowbreak{}vol\_\allowbreak{}surface} (based on asset type). Extract ATM vol term structure, 25-delta risk reversals (skew), and butterflies (smile curvature).
\item {} \textbf{Template Discovery:} Call \emph{option\_\allowbreak{}template\_\allowbreak{}list} to find available option types, expiries, and strikes for the underlying.
\item {} \textbf{Option Pricing:} Call \emph{option\_\allowbreak{}value} for specific options of interest. Extract premium, delta, gamma, vega, theta, implied vol.
\item {} \textbf{Historical Data:} Call \emph{tscc\_\allowbreak{}historical\_\allowbreak{}pricing\_\allowbreak{}summaries} or \emph{qa\_\allowbreak{}historical\_\allowbreak{}equity\_\allowbreak{}price} for 1Y daily history.
\item {} \textbf{Realized Vol Computation:} From historical prices, compute close-to-close realized vol over 20-day, 60-day, and 90-day windows. Compare to matching implied vol tenors.
\item {} \textbf{Synthesize:} Combine surface shape, Greeks, and implied-vs-realized comparison into a vol assessment with strategy recommendations.
\end{enumerate}

\subparagraph{Output Format}\mbox{}\par

\subparagraph{Vol Surface Summary}\mbox{}\par
\begingroup
\small
\setlength{\tabcolsep}{2pt}
\begin{longtable}{>{\RaggedRight\arraybackslash\hspace{0pt}}p{0.240\textwidth}>{\RaggedRight\arraybackslash\hspace{0pt}}p{0.240\textwidth}>{\RaggedRight\arraybackslash\hspace{0pt}}p{0.240\textwidth}>{\RaggedRight\arraybackslash\hspace{0pt}}p{0.240\textwidth}}
\toprule
{} \textbf{Tenor} & \textbf{ATM Vol} & \textbf{25d RR} & \textbf{25d BF} \\
\midrule
{} 1M & ... & ... & ... \\
{} 3M & ... & ... & ... \\
{} 6M & ... & ... & ... \\
{} 1Y & ... & ... & ... \\
\bottomrule
\end{longtable}
\endgroup

\subparagraph{Greeks Table}\mbox{}\par
\begingroup
\small
\setlength{\tabcolsep}{2pt}
\begin{longtable}{>{\RaggedRight\arraybackslash\hspace{0pt}}p{0.320\textwidth}>{\RaggedRight\arraybackslash\hspace{0pt}}p{0.320\textwidth}>{\RaggedRight\arraybackslash\hspace{0pt}}p{0.320\textwidth}}
\toprule
{} \textbf{Greek} & \textbf{Call} & \textbf{Put} \\
\midrule
{} Premium & ... & ... \\
{} Delta & ... & ... \\
{} Gamma & ... & ... \\
{} Vega & ... & ... \\
{} Theta & ... & ... \\
{} Implied Vol & ... & ... \\
\bottomrule
\end{longtable}
\endgroup

\subparagraph{Implied vs Realized Comparison}\mbox{}\par
\begingroup
\small
\setlength{\tabcolsep}{2pt}
\begin{longtable}{>{\RaggedRight\arraybackslash\hspace{0pt}}p{0.184\textwidth}>{\RaggedRight\arraybackslash\hspace{0pt}}p{0.184\textwidth}>{\RaggedRight\arraybackslash\hspace{0pt}}p{0.184\textwidth}>{\RaggedRight\arraybackslash\hspace{0pt}}p{0.184\textwidth}>{\RaggedRight\arraybackslash\hspace{0pt}}p{0.184\textwidth}}
\toprule
{} \textbf{Window} & \textbf{Realized Vol} & \textbf{Implied Vol (matching tenor)} & \textbf{Premium (IV - RV)} & \textbf{Signal} \\
\midrule
{} 20d & ... & 1M ATM & ... & Rich/ Cheap \\
{} 60d & ... & 3M ATM & ... & Rich/ Cheap \\
{} 90d & ... & 6M ATM & ... & Rich/ Cheap \\
\bottomrule
\end{longtable}
\endgroup

\subparagraph{Assessment}\mbox{}\par
State the vol regime (low/normal/elevated/crisis), whether implied is rich or cheap vs realized, surface shape signals (skew direction, term structure shape), and recommended strategies with key Greeks and rationale.


\vspace{0.8em}\hrule\vspace{0.8em}

\subsection{Skill: swap-curve-strategy}\label{file:lseg-skills-swap-curve-strategy-SKILL.md}

\paragraph{File Metadata}\mbox{}\par
\begingroup
\small
\setlength{\tabcolsep}{2pt}
\begin{longtable}{>{\RaggedRight\arraybackslash\hspace{0pt}}p{0.480\textwidth}>{\RaggedRight\arraybackslash\hspace{0pt}}p{0.480\textwidth}}
\toprule
{} \textbf{Field} & \textbf{Value} \\
\midrule
{} name & swap-curve-strategy \\
{} description & Analyze the interest rate swap curve by pricing swaps at multiple tenors, overlaying government and inflation curves, and identifying curve trade opportunities. Use when analyzing swap curves, computing swap spreads, decomposing real rates, identifying steepener/ flattener/ butterfly trades, or comparing swap rates across currencies. \\
\bottomrule
\end{longtable}
\endgroup


\paragraph{Swap Curve Strategy Analysis}\mbox{}\par

You are an expert rates strategist specializing in swap curve analysis. Combine swap pricing, government yield curves, and inflation curves from MCP tools to analyze curve shape, compute swap spreads, decompose real rates, and identify curve trade opportunities. Focus on routing tool outputs into curve metrics and trade recommendations — let the tools price, you analyze the shape and recommend.


\subparagraph{Core Principles}\mbox{}\par

The swap curve prices the market's expectation of future short-term rates, credit conditions, and funding costs. Always build the full swap curve first, overlay the government curve to compute swap spreads, then add inflation breakevens for real rate decomposition. Curve metrics (2s10s slope, 5s30s slope, butterfly) and their historical context drive trade ideas. For trade recommendations, always include DV01-neutral sizing and carry/roll-down estimates.


\subparagraph{Available MCP Tools}\mbox{}\par

\begin{itemize}[leftmargin=*,itemsep=2pt,topsep=2pt]
\item {} \textbf{@@TOKEN0@@} — Swap pricing. Two-phase: list templates (by currency/index) then price at specific tenors. Returns par swap rate, DV01, NPV.
\item {} \textbf{@@TOKEN0@@} — Government yield curves. Two-phase: list then calculate. Use for swap spread computation and curve shape context.
\item {} \textbf{@@TOKEN0@@} — Inflation breakeven curves. Two-phase: search then calculate. Use for real rate decomposition.
\item {} \textbf{@@TOKEN0@@} — Historical pricing data. Use for historical curve slope context and trend analysis.
\item {} \textbf{@@TOKEN0@@} — Macro data. Use to establish economic context for curve analysis and assess consistency with curve signals.
\end{itemize}

\subparagraph{Tool Chaining Workflow}\mbox{}\par

\begin{enumerate}[leftmargin=*,itemsep=2pt,topsep=2pt]
\item {} \textbf{Discover Swap Templates:} Call \emph{ir\_\allowbreak{}swap} in list mode for the target currency. Identify available indices and tenors.
\item {} \textbf{Build Swap Curve:} Call \emph{ir\_\allowbreak{}swap} in price mode for standard tenors (2Y, 5Y, 7Y, 10Y, 20Y, 30Y). Extract par swap rate and DV01 at each point.
\item {} \textbf{Overlay Government Curve:} Call \emph{interest\_\allowbreak{}rate\_\allowbreak{}curve} (list then calculate) for the same currency. Compute swap spread = swap rate minus government yield at each tenor.
\item {} \textbf{Inflation Decomposition:} Call \emph{inflation\_\allowbreak{}curve} (search then calculate). Compute real rate = nominal swap rate minus inflation breakeven at each tenor.
\item {} \textbf{Compute Curve Metrics:} From the swap curve: 2s10s slope, 5s30s slope, 2s5s10s butterfly. Note curve shape classification.
\item {} \textbf{Synthesize:} Combine into a complete analysis with swap curve table, swap spreads, real rate decomposition, curve metrics, and trade recommendations with DV01-neutral sizing.
\end{enumerate}

\subparagraph{Output Format}\mbox{}\par

\subparagraph{Swap Curve Table}\mbox{}\par
\begingroup
\small
\setlength{\tabcolsep}{2pt}
\begin{longtable}{>{\RaggedRight\arraybackslash\hspace{0pt}}p{0.131\textwidth}>{\RaggedRight\arraybackslash\hspace{0pt}}p{0.131\textwidth}>{\RaggedRight\arraybackslash\hspace{0pt}}p{0.131\textwidth}>{\RaggedRight\arraybackslash\hspace{0pt}}p{0.131\textwidth}>{\RaggedRight\arraybackslash\hspace{0pt}}p{0.131\textwidth}>{\RaggedRight\arraybackslash\hspace{0pt}}p{0.131\textwidth}>{\RaggedRight\arraybackslash\hspace{0pt}}p{0.131\textwidth}}
\toprule
{} \textbf{Tenor} & \textbf{Swap Rate (\%)} & \textbf{Govt Yield (\%)} & \textbf{Swap Spread (bp)} & \textbf{DV01} & \textbf{Inflation BE (\%)} & \textbf{Real Rate (\%)} \\
\midrule
{} 2Y & ... & ... & ... & ... & ... & ... \\
{} 5Y & ... & ... & ... & ... & ... & ... \\
{} 10Y & ... & ... & ... & ... & ... & ... \\
{} 30Y & ... & ... & ... & ... & ... & ... \\
\bottomrule
\end{longtable}
\endgroup

\subparagraph{Curve Metrics}\mbox{}\par
\begingroup
\small
\setlength{\tabcolsep}{2pt}
\begin{longtable}{>{\RaggedRight\arraybackslash\hspace{0pt}}p{0.480\textwidth}>{\RaggedRight\arraybackslash\hspace{0pt}}p{0.480\textwidth}}
\toprule
{} \textbf{Metric} & \textbf{Current} \\
\midrule
{} 2s10s slope (bp) & ... \\
{} 5s30s slope (bp) & ... \\
{} 2s5s10s butterfly (bp) & ... \\
{} Curve shape & Normal /  Flat /  Inverted /  Humped \\
\bottomrule
\end{longtable}
\endgroup

\subparagraph{Real Rate Decomposition}\mbox{}\par
\begingroup
\small
\setlength{\tabcolsep}{2pt}
\begin{longtable}{>{\RaggedRight\arraybackslash\hspace{0pt}}p{0.184\textwidth}>{\RaggedRight\arraybackslash\hspace{0pt}}p{0.184\textwidth}>{\RaggedRight\arraybackslash\hspace{0pt}}p{0.184\textwidth}>{\RaggedRight\arraybackslash\hspace{0pt}}p{0.184\textwidth}>{\RaggedRight\arraybackslash\hspace{0pt}}p{0.184\textwidth}}
\toprule
{} \textbf{Tenor} & \textbf{Nominal Swap} & \textbf{Inflation BE} & \textbf{Real Rate} & \textbf{Signal} \\
\midrule
{} 2Y & ...\% & ...\% & ...\% & Accommodative/ Restrictive \\
{} 5Y & ...\% & ...\% & ...\% & Accommodative/ Restrictive \\
{} 10Y & ...\% & ...\% & ...\% & Accommodative/ Restrictive \\
\bottomrule
\end{longtable}
\endgroup

\subparagraph{Curve Trade Recommendation}\mbox{}\par
For each trade: structure (e.g., 2s10s steepener), legs, DV01-neutral notionals, estimated 3M carry, estimated 3M roll-down, breakeven curve move, target, stop-loss, and thesis (1-2 sentences).


\vspace{0.8em}\hrule\vspace{0.8em}

\subsection{Directory: spglobal}

\subsection{File: spglobal/README.md}\label{file:spglobal-README.md}

\paragraph{S\&P Global Plugin}\mbox{}\par

This plugin brings S\&P Global's financial data and analytics directly into your AI workflow via a set of pre-built skills. It is designed for financial professionals who want AI-assisted research, analysis, and document generation grounded in authoritative S\&P Global data.


The skills are built on open standards (MCP) and are designed to work across AI platforms and agent frameworks. While the plugin follows the Claude Cowork standard, all skills and the underlying data layer are platform-agnostic. If you want to use these skills in another environment, please do.


We recognize that every organization has unique needs. These skills are starting points to get you started on your finished products. We encourage you to adapt the prompts, outputs, and workflows to fit your firm's specific processes, templates, and data needs.


The skills in this plugin are provided as-is. Generated outputs and data are not guaranteed to be correct. \textbf{Always verify outputs generated by an LLM for correctness.}


\subparagraph{Skills Contained}\mbox{}\par

\subparagraph{Tearsheets}\mbox{}\par
\textbf{Requires}: \href{https://www.marketplace.spglobal.com/en/solutions/kensho-llm-ready-api-\%28a156fe9f-5564-4f60-a624-95d8645dc98f\%29}{S\&P Global LLM-ready API} subscription


Generates a formatted, one-to-two page company tearsheet as a Word document, populated with live data from S\&P Capital IQ. Supports four audience types, each optimized for a different use case:

\begin{itemize}[leftmargin=*,itemsep=2pt,topsep=2pt]
\item {} Equity Research: Investment thesis snapshot for buy-side/sell-side analysts
\item {} Investment Banking / M\&A: Company profile in transaction context
\item {} Corporate Development: Acquisition target profile for internal strategic teams
\item {} Sales / Business Development: Client meeting prep for commercial teams
\end{itemize}

\textbf{Example prompt}: ``Generate a business development tearsheet for Palantir.''


\subparagraph{Industry Transaction Summaries}\mbox{}\par
\textbf{Requires}: \href{https://www.marketplace.spglobal.com/en/solutions/kensho-llm-ready-api-\%28a156fe9f-5564-4f60-a624-95d8645dc98f\%29}{S\&P Global LLM-ready API} subscription


Summarizes recent M\&A and deal activity within a sector or for a specific company, drawing on S\&P Capital IQ transaction data. Useful for market mapping, pitch preparation, and competitive intelligence.


\textbf{Example prompt}: ``Summarize recent transactions in the data infrastructure space”


\subparagraph{Earnings Previews}\mbox{}\par
\textbf{Requires}: \href{https://www.marketplace.spglobal.com/en/solutions/kensho-llm-ready-api-\%28a156fe9f-5564-4f60-a624-95d8645dc98f\%29}{S\&P Global LLM-ready API} subscription


Generates a structured earnings preview for an upcoming report, including consensus estimates, recent guidance, analyst sentiment, and key things to watch — all sourced from S\&P Capital IQ.


\textbf{Example prompt}: ``Give me an earnings preview for Salesforce.''


\subparagraph{How to Use}\mbox{}\par
The plugin and skills require access to S\&P Global data to work with, either \href{https://www.spglobal.com/market-intelligence/en/solutions/products/sp-capital-iq-pro}{Capital IQ Pro} or \href{https://www.marketplace.spglobal.com/en/solutions/kensho-llm-ready-api-\%28a156fe9f-5564-4f60-a624-95d8645dc98f\%29}{S\&P Global LLM-ready API} subscriptions.


The LLM-ready API is easy to integrate in Claude or other applications via its MCP server. Follow \href{docs.kensho.com/llmreadyapi/mcp/third-party/claude}{these steps} to set it up.


\subparagraph{In Cowork}\mbox{}\par
You'll need a paid Claude plan (Pro, Max, Team, or Enterprise) and the Claude Desktop app for macOS or Windows.


\begin{enumerate}[leftmargin=*,itemsep=2pt,topsep=2pt]
\item {} Open Claude Desktop and navigate to the \textbf{Cowork} tab
\item {} Click \textbf{Customize with Plugins}
\item {} In Browse Plugins, select \textbf{Personal}
\item {} Click the \textbf{plus sign “+”} to add plug in
\item {} Authenticate with your S\&P Global credentials when prompted
\end{enumerate}

Once installed, skills activate automatically when relevant – just describe what you need in natural language. You can also invoke specific skills explicitly by typing / in the chat to see available commands.


To customize the plugin for your firm's workflows, templates, or terminology, click \textbf{Customize} while viewing the installed plugin. We encourage this; the defaults are a starting point, not a prescription.


\subparagraph{In Claude Desktop (individual skills)}\mbox{}\par
To install individual skills in Claude Desktop without the full plugin:


\begin{enumerate}[leftmargin=*,itemsep=2pt,topsep=2pt]
\item {} Open \textbf{Settings}
\item {} Navigate to \textbf{Capabilities → Skills}
\item {} Click \textbf{Add}
\item {} Upload the skill file(s) from this repository
\end{enumerate}

Skills are available immediately in your Claude Desktop sessions once uploaded. You can install one or several depending on your needs.


\subparagraph{In Claude Code (individual skills)}\mbox{}\par

Follow the instructions on the \href{https://code.claude.com/docs/en/discover-plugins\#add-from-github}{Claude Code documentation}.


\subparagraph{Other Platforms}\mbox{}\par
The skills in this repository are markdown files. Any AI platform that supports custom instructions, system prompts, or knowledge file uploads can use them – the mechanism varies by platform but the principle is the same: load the skill content so the model has it as persistent context.


\textbf{ChatGPT}: Paste skill content into Custom Instructions (Settings → Customize ChatGPT), upload as a knowledge file within a Project, or add to a Custom GPT's configuration. Custom instructions apply globally across sessions; Project-level files scope context to a specific workflow.


\textbf{Microsoft Copilot}: Paste skill content into a custom prompt or system instruction depending on your Copilot configuration (M365 Copilot, Copilot Studio, etc.). Enterprise deployments via Copilot Studio allow uploading knowledge sources directly.


\textbf{Other platforms}: If your platform supports a system prompt or persistent instruction layer, paste the skill markdown there. If it supports file-based knowledge retrieval, upload the skill file. The skills are plain markdown, no special formatting or tooling required.


\subparagraph{What Comes Next}\mbox{}\par
We are actively building additional skills and plugins across a range of financial workflows. We'd love to hear what would be most useful to you! For general questions, feedback, or partnership inquiries, contact \href{mailto:commercial@kensho.com}{commercial@kensho.com} or open an issue in this repository.


\paragraph{License}\mbox{}\par

Licensed under the Apache 2.0 License. Unless required by applicable law or agreed to in writing, software distributed under the License is distributed on an ``AS IS'' BASIS, WITHOUT WARRANTIES OR CONDITIONS OF ANY KIND, either express or implied. See the License for the specific language governing permissions and limitations under the License.


Copyright 2026-present Kensho Technologies, LLC. The present date is determined by the timestamp of the most recent commit in the repository.


\vspace{0.8em}\hrule\vspace{0.8em}

\subsection{Skill: earnings-preview-beta}\label{file:spglobal-skills-earnings-preview-beta-SKILL.md}

\paragraph{File Metadata}\mbox{}\par
\begingroup
\small
\setlength{\tabcolsep}{2pt}
\begin{longtable}{>{\RaggedRight\arraybackslash\hspace{0pt}}p{0.480\textwidth}>{\RaggedRight\arraybackslash\hspace{0pt}}p{0.480\textwidth}}
\toprule
{} \textbf{Field} & \textbf{Value} \\
\midrule
{} name & earnings-preview-single \\
{} description & Generate a concise 4-5 page equity research earnings preview for a single company. Analyzes the most recent earnings transcript, competitor landscape, valuation, and recent news to produce a professional HTML report. \\
\bottomrule
\end{longtable}
\endgroup


\paragraph{Single-Company Earnings Preview}\mbox{}\par

Generate a concise, professional equity research earnings preview for a single company. The output is a self-contained HTML file targeting 4-5 printed pages. The report is dense with figures and data, with tight narrative that gets straight to the point.


\textbf{Data Sources (ZERO EXCEPTIONS):} The ONLY permitted data sources are \textbf{Kensho Grounding MCP} (\emph{search}) and \textbf{S\&P Global MCP} (\emph{kfinance}). Absolutely NO other tools, data sources, or web access of any kind. Specifically:

\begin{itemize}[leftmargin=*,itemsep=2pt,topsep=2pt]
\item {} Do NOT use \emph{WebSearch}, \emph{WebFetch}, \emph{web\_\allowbreak{}search}, \emph{brave\_\allowbreak{}search}, \emph{google\_\allowbreak{}search}, or ANY generic web/internet search tool — even if Kensho is slow, returns no results, or is temporarily unavailable.
\item {} Do NOT use any browser, URL fetch, or web scraping tool.
\item {} If Kensho Grounding returns no results for a query, try rephrasing the query or note ``data not available'' in the report. \textbf{NEVER fall back to web search as an alternative.}
\item {} Every piece of information in the report must be traceable to either a \emph{kfinance} MCP function call or a Kensho \emph{search} call. If it cannot be sourced to one of these two, it must not appear in the report.
\end{itemize}

\textbf{Critical Rule:} You MUST complete ALL research and data collection (Phases 1-5) BEFORE writing any part of the report.


\textbf{Intermediate File Rule:} All raw data from MCP tool calls MUST be written to files in \emph{/tmp/earnings-preview/} \textbf{immediately after each tool call returns} — before moving to the next call. This protects data from context window compression. Do NOT hold data only in memory. At the start of Phase 1, run \emph{mkdir -p /tmp/earnings-preview} to create the directory. \textbf{Before generating the HTML report (Phase 7), you MUST read ALL intermediate files back into context using @@TOKEN2@@ commands. The files — not your memory of earlier conversation — are the single source of truth for every number, quote, and source URL in the report. If you skip reading the files, the report WILL contain errors.}


\textbf{Fiscal Quarter Rule:} NEVER infer the fiscal quarter from the calendar report date. Many companies have non-standard fiscal years (e.g., Walmart's FY ends Jan 31, so a Feb 2026 report covers Q4 FY2026, not Q4 2025 or Q1 2026). Always use the fiscal quarter and fiscal year exactly as stated in the earnings call name returned by \emph{get\_\allowbreak{}next\_\allowbreak{}earnings\_\allowbreak{}from\_\allowbreak{}identifiers} or \emph{get\_\allowbreak{}earnings\_\allowbreak{}from\_\allowbreak{}identifiers} (e.g., ``Walmart Q4 FY2026 Earnings Call'' means the quarter is Q4 FY2026). Use that verbatim in the report title, headers, tables, and all references. If the call name is ambiguous, cross-reference with \emph{get\_\allowbreak{}financial\_\allowbreak{}line\_\allowbreak{}item\_\allowbreak{}from\_\allowbreak{}identifiers} period labels.


\textbf{Length Rule:} The report must be concise. Target 4-5 pages when printed. Do NOT write long multi-paragraph narratives. Use tight, punchy bullet points. Every sentence must earn its place. If you can say it in fewer words, do so.


\textbf{Verbatim Quote Rule:} When quoting management in \emph{<blockquote>} tags, the text MUST be copied \textbf{exactly} from the transcript — word for word, including filler words and sentence fragments. Do NOT paraphrase, rearrange, combine sentences from different parts of the transcript, or ``clean up'' quotes. If you cannot find the exact phrase in the transcript, do NOT present it as a direct quote. Instead, paraphrase in your own narrative voice without blockquote formatting (e.g., ``Management noted that data center demand remains significant''). Every blockquote must be a verbatim, copy-paste excerpt that can be verified against the transcript.


\textbf{Calculation Integrity Rule:} For any multi-step calculation (implied quarterly figures from annual guidance, LTM P/E, y/y growth rates, segment y/y changes), write out each step explicitly and verify intermediate results before using them in the next step. If you state A + B + C = X, verify X is arithmetically correct before using X in a subsequent formula. If the appendix shows a sum that does not equal its stated components, the report is wrong. When in doubt, recompute from raw data rather than reusing a previously calculated intermediate.


\textbf{Ratio Nomenclature Rule:} All valuation ratios must be explicitly labeled as \textbf{LTM} (Last Twelve Months) or \textbf{NTM} (Next Twelve Months). Never use ``trailing'' or ``forward'' — always use LTM or NTM. LTM ratios use the sum of the most recent 4 reported quarters. NTM ratios use the \textbf{sum of the next 4 quarterly consensus mean EPS estimates} from \emph{get\_\allowbreak{}consensus\_\allowbreak{}estimates\_\allowbreak{}from\_\allowbreak{}identifiers} — NOT a single annual figure. Both LTM and NTM P/E must be computed and displayed in the competitor comparison table.


\textbf{Hyperlink Rule (STRICTLY ENFORCED):} Every claim in the report — numeric AND non-numeric — MUST be wrapped in an \emph{<a href=``\#ref-N'' class=``data-ref''>} hyperlink pointing to the corresponding entry in the Appendix. \textbf{This is not optional. Every single number in the report must be a clickable link.} This includes: revenue figures, EPS, margins, growth rates, market caps, P/E ratios, stock returns, price targets, segment revenue, and any other financial metric. It also includes qualitative claims from transcripts or Kensho searches. If you state it as fact, it must link to a source. Assign each unique claim a sequential reference ID (\emph{ref-1}, \emph{ref-2}, etc.). The hyperlink style is subtle — navy color, no underline, dotted underline on hover. \textbf{Do NOT write any number in the report body without wrapping it in an @@TOKEN3@@ tag.} Example: write \emph{<a href=``\#ref-1'' class=``data-ref''>\$152.3B</a>}, NEVER write \emph{\$152.3B} as plain text.


\vspace{0.5em}\hrule\vspace{0.5em}

\subparagraph{Phase 1: Company Profile \& Setup}\mbox{}\par

\begin{enumerate}[leftmargin=*,itemsep=2pt,topsep=2pt]
\item {} Parse the single company ticker from \emph{\$ARGUMENTS} (strip whitespace).
\item {} Run \emph{mkdir -p /tmp/earnings-preview} to create the working directory.
\item {} Call \emph{get\_\allowbreak{}latest()} to establish current reporting period context.
\item {} Call \emph{get\_\allowbreak{}info\_\allowbreak{}from\_\allowbreak{}identifiers} — record market cap, industry.
\item {} Call \emph{get\_\allowbreak{}company\_\allowbreak{}summary\_\allowbreak{}from\_\allowbreak{}identifiers} — record business description.
\item {} Call \emph{get\_\allowbreak{}next\_\allowbreak{}earnings\_\allowbreak{}from\_\allowbreak{}identifiers} — record upcoming earnings date and fiscal quarter name.
\end{enumerate}

\textbf{Immediately write} \emph{/tmp/earnings-preview/company-info.txt}:

\textbf{Structured Template}
\begin{itemize}[leftmargin=*,itemsep=2pt,topsep=2pt]
\item {} TICKER: [ticker]
\item {} COMPANY: [full name]
\item {} INDUSTRY: [industry]
\item {} MARKET\_\allowbreak{}CAP: [value] (as of [date])
\item {} NEXT\_\allowbreak{}EARNINGS\_\allowbreak{}DATE: [date]
\item {} NEXT\_\allowbreak{}EARNINGS\_\allowbreak{}QUARTER: [Q\# FY\#\#\#\# exactly as returned by API]
\item {} BUSINESS\_\allowbreak{}DESCRIPTION: [2-3 sentence summary]
\end{itemize}

\vspace{0.5em}\hrule\vspace{0.5em}

\subparagraph{Phase 2: Earnings Transcript Analysis (MANDATORY — COMPLETE BEFORE WRITING)}\mbox{}\par

\begin{enumerate}[leftmargin=*,itemsep=2pt,topsep=2pt]
\item {} Call \emph{get\_\allowbreak{}latest\_\allowbreak{}earnings\_\allowbreak{}from\_\allowbreak{}identifiers} to get the most recent completed earnings call \emph{key\_\allowbreak{}dev\_\allowbreak{}id}.
\item {} Call \emph{get\_\allowbreak{}transcript\_\allowbreak{}from\_\allowbreak{}key\_\allowbreak{}dev\_\allowbreak{}id} for that transcript.
\item {} \textbf{Immediately write} \emph{/tmp/earnings-preview/transcript-extracts.txt} with the following sections. Write this file WHILE you still have the transcript in context — do not wait:
\end{enumerate}

\textbf{Structured Template}
\begin{itemize}[leftmargin=*,itemsep=2pt,topsep=2pt]
\item {} TRANSCRIPT\_\allowbreak{}SOURCE: [Call Name, e.g., ``Q3 2025 Earnings Call'']
\item {} KEY\_\allowbreak{}DEV\_\allowbreak{}ID: [key\_\allowbreak{}dev\_\allowbreak{}id]
\item {} CALL\_\allowbreak{}DATE: [date]
\item {} FISCAL\_\allowbreak{}QUARTER: [Q\# FY\#\#\#\#]
\item {} === VERBATIM QUOTES (copy-paste exactly — do NOT paraphrase) ===
\item {} QUOTE\_\allowbreak{}1: ``[exact text from transcript]''
\item {} SPEAKER\_\allowbreak{}1: [Name], [Title]
\item {} CONTEXT\_\allowbreak{}1: [1 sentence on where this appeared — prepared remarks or Q\&A]
\item {} QUOTE\_\allowbreak{}2: ``[exact text from transcript]''
\item {} SPEAKER\_\allowbreak{}2: [Name], [Title]
\item {} CONTEXT\_\allowbreak{}2: [context]
\item {} QUOTE\_\allowbreak{}3: ``[exact text from transcript]''
\item {} SPEAKER\_\allowbreak{}3: [Name], [Title]
\item {} CONTEXT\_\allowbreak{}3: [context]
\item {} QUOTE\_\allowbreak{}4: ``[exact text from transcript]''
\item {} SPEAKER\_\allowbreak{}4: [Name], [Title]
\item {} CONTEXT\_\allowbreak{}4: [context]
\item {} === GUIDANCE (quantitative only) ===
\item {} - [metric]: [range or point estimate as stated by management]
\item {} - [metric]: [range or point estimate]
\item {} === KEY DRIVERS ===
\item {} - [driver 1 with supporting data point]
\item {} - [driver 2 with supporting data point]
\item {} - [driver 3 with supporting data point]
\item {} === HEADWINDS \& RISKS ===
\item {} - [risk 1 with quantification if available]
\item {} - [risk 2]
\item {} === ANALYST Q\&A THEMES ===
\item {} - [theme 1: what analysts pushed on]
\item {} - [theme 2]
\item {} - [theme 3]
\item {} === SYNTHESIS: THEMES TO WATCH NEXT QUARTER ===
\item {} - [theme 1]
\item {} - [theme 2]
\item {} - [theme 3]
\end{itemize}

\vspace{0.5em}\hrule\vspace{0.5em}

\subparagraph{Phase 3: Competitor Analysis}\mbox{}\par

\begin{enumerate}[leftmargin=*,itemsep=2pt,topsep=2pt]
\item {} Call \emph{get\_\allowbreak{}competitors\_\allowbreak{}from\_\allowbreak{}identifiers} with \emph{competitor\_\allowbreak{}source=``all''}.
\item {} Select \textbf{top 5-7 most relevant public competitors}.
\item {} For the company AND all selected competitors, gather:
\begin{itemize}[leftmargin=*,itemsep=2pt,topsep=2pt]
\item {} \emph{get\_\allowbreak{}prices\_\allowbreak{}from\_\allowbreak{}identifiers} with \emph{periodicity=``day''}, last 12 months
\item {} \emph{get\_\allowbreak{}financial\_\allowbreak{}line\_\allowbreak{}item\_\allowbreak{}from\_\allowbreak{}identifiers} for \emph{diluted\_\allowbreak{}eps}, \emph{period\_\allowbreak{}type=``quarterly''}, \emph{num\_\allowbreak{}periods=8}
\item {} \emph{get\_\allowbreak{}capitalization\_\allowbreak{}from\_\allowbreak{}identifiers} with \emph{capitalization=``market\_\allowbreak{}cap''} (latest)
\item {} \emph{get\_\allowbreak{}consensus\_\allowbreak{}estimates\_\allowbreak{}from\_\allowbreak{}identifiers} with \emph{period\_\allowbreak{}type=``quarterly''}, \emph{num\_\allowbreak{}periods\_\allowbreak{}forward=4} — this returns consensus mean EPS estimates for the next 4 quarters, which are summed to compute NTM EPS
\end{itemize}
\end{enumerate}

\textbf{After each tool call returns, immediately append the raw data to the appropriate intermediate file:}


\textbf{Write} \emph{/tmp/earnings-preview/prices.csv} — one row per (ticker, date, close). Include the \emph{source} column with the exact MCP function call. Write the subject company's prices first, then each competitor's as you fetch them:

\textbf{Structured Template}
\begin{itemize}[leftmargin=*,itemsep=2pt,topsep=2pt]
\item {} ticker,date,close,source
\item {} D,2025-02-19,55.67,get\_\allowbreak{}prices\_\allowbreak{}from\_\allowbreak{}identifiers(identifier='D',periodicity='day')
\item {} D,2025-02-20,55.82,get\_\allowbreak{}prices\_\allowbreak{}from\_\allowbreak{}identifiers(identifier='D',periodicity='day')
\item {} ...
\item {} DUK,2025-02-19,111.79,get\_\allowbreak{}prices\_\allowbreak{}from\_\allowbreak{}identifiers(identifier='DUK',periodicity='day')
\item {} ...
\end{itemize}
Note: the \emph{source} value is the same for all rows from a single call — write it on every row so it's always available.


\textbf{Write} \emph{/tmp/earnings-preview/peer-eps.csv} — one row per (ticker, period, eps). Write immediately after each \emph{diluted\_\allowbreak{}eps} call:

\textbf{Structured Template}
\begin{itemize}[leftmargin=*,itemsep=2pt,topsep=2pt]
\item {} ticker,period,diluted\_\allowbreak{}eps,source
\item {} D,Q4 2024,1.09,get\_\allowbreak{}financial\_\allowbreak{}line\_\allowbreak{}item\_\allowbreak{}from\_\allowbreak{}identifiers(identifier='D',line\_\allowbreak{}item='diluted\_\allowbreak{}eps',period\_\allowbreak{}type='quarterly')
\item {} D,Q1 2025,-0.11,get\_\allowbreak{}financial\_\allowbreak{}line\_\allowbreak{}item\_\allowbreak{}from\_\allowbreak{}identifiers(identifier='D',line\_\allowbreak{}item='diluted\_\allowbreak{}eps',period\_\allowbreak{}type='quarterly')
\item {} ...
\item {} DUK,Q4 2024,1.52,get\_\allowbreak{}financial\_\allowbreak{}line\_\allowbreak{}item\_\allowbreak{}from\_\allowbreak{}identifiers(identifier='DUK',line\_\allowbreak{}item='diluted\_\allowbreak{}eps',period\_\allowbreak{}type='quarterly')
\item {} ...
\end{itemize}

\textbf{Write} \emph{/tmp/earnings-preview/peer-market-caps.csv} — one row per ticker. Write immediately after each \emph{market\_\allowbreak{}cap} call:

\textbf{Structured Template}
\begin{itemize}[leftmargin=*,itemsep=2pt,topsep=2pt]
\item {} ticker,market\_\allowbreak{}cap,retrieval\_\allowbreak{}date,source
\item {} D,55900000000,2026-02-19,get\_\allowbreak{}capitalization\_\allowbreak{}from\_\allowbreak{}identifiers(identifier='D',capitalization='market\_\allowbreak{}cap')
\item {} DUK,98300000000,2026-02-19,get\_\allowbreak{}capitalization\_\allowbreak{}from\_\allowbreak{}identifiers(identifier='DUK',capitalization='market\_\allowbreak{}cap')
\item {} ...
\end{itemize}

\textbf{Write} \emph{/tmp/earnings-preview/consensus-eps.csv} — one row per (ticker, period, consensus mean EPS). Write immediately after each \emph{get\_\allowbreak{}consensus\_\allowbreak{}estimates\_\allowbreak{}from\_\allowbreak{}identifiers} call:

\textbf{Structured Template}
\begin{itemize}[leftmargin=*,itemsep=2pt,topsep=2pt]
\item {} ticker,period,consensus\_\allowbreak{}mean\_\allowbreak{}eps,num\_\allowbreak{}estimates,source
\item {} D,Q4 2025,0.88,12,get\_\allowbreak{}consensus\_\allowbreak{}estimates\_\allowbreak{}from\_\allowbreak{}identifiers(identifier='D',period\_\allowbreak{}type='quarterly',num\_\allowbreak{}periods\_\allowbreak{}forward=4)
\item {} D,Q1 2026,0.72,10,get\_\allowbreak{}consensus\_\allowbreak{}estimates\_\allowbreak{}from\_\allowbreak{}identifiers(identifier='D',period\_\allowbreak{}type='quarterly',num\_\allowbreak{}periods\_\allowbreak{}forward=4)
\item {} D,Q2 2026,0.91,9,get\_\allowbreak{}consensus\_\allowbreak{}estimates\_\allowbreak{}from\_\allowbreak{}identifiers(identifier='D',period\_\allowbreak{}type='quarterly',num\_\allowbreak{}periods\_\allowbreak{}forward=4)
\item {} D,Q3 2026,1.05,8,get\_\allowbreak{}consensus\_\allowbreak{}estimates\_\allowbreak{}from\_\allowbreak{}identifiers(identifier='D',period\_\allowbreak{}type='quarterly',num\_\allowbreak{}periods\_\allowbreak{}forward=4)
\item {} DUK,Q4 2025,1.48,14,get\_\allowbreak{}consensus\_\allowbreak{}estimates\_\allowbreak{}from\_\allowbreak{}identifiers(identifier='DUK',period\_\allowbreak{}type='quarterly',num\_\allowbreak{}periods\_\allowbreak{}forward=4)
\item {} ...
\end{itemize}

\begin{enumerate}[leftmargin=*,itemsep=2pt,topsep=2pt]
\item {} \textbf{Do NOT calculate P/E or returns yet.} The raw data is now on disk. Calculations happen in Phase 6 (Verification), reading from these files.
\end{enumerate}

\textbf{Date Consistency Rule (stock returns):} When computing comparative stock returns (YTD \%, 1-yr \%, 30d \%, 90d \%), ALL tickers MUST use the \textbf{exact same start and end dates}. After writing all price data to \emph{prices.csv}, identify the first trading date that appears in ALL tickers' data and use that as the common base date. Do NOT use different base dates for different tickers (e.g., the subject from Feb 19 and peers from Feb 28). If a ticker's data starts later than others, use the first overlapping date for ALL calculations. State the common base date in the appendix for every return calculation.


\textbf{P/E Currency Rule (LTM P/E):} When computing LTM P/E for each company, use that company's \textbf{most recent 4 reported quarters} from \emph{peer-eps.csv} — not a fixed calendar window applied to all. If a peer has already reported Q4 2025 while the subject company has only reported through Q3 2025, the peer's LTM EPS should include Q4 2025. Check the latest reported period for each company and use the 4 most recent periods per company. Note in the appendix which 4 quarters were used for each P/E calculation.


\textbf{Market Cap Date-Stamp:} When reporting market cap, use the \emph{retrieval\_\allowbreak{}date} from \emph{peer-market-caps.csv}. If it differs from the report date, note this in the appendix.


\vspace{0.5em}\hrule\vspace{0.5em}

\subparagraph{Phase 4: News, Estimates \& Sector Intelligence (via Kensho Grounding)}\mbox{}\par

Run these \emph{search} queries for \textbf{each} category below. Do NOT skip any.


\textbf{CRITICAL — Capture Source URLs:} Every Kensho \emph{search} result includes a \textbf{source URL} for the underlying article, report, or data page. You MUST record the URL alongside each finding.


\textbf{After EACH search call, immediately append the results to} \emph{/tmp/earnings-preview/kensho-findings.txt} using the format below. Do NOT wait until all searches are done — write after each one:


\textbf{Structured Template}
\begin{itemize}[leftmargin=*,itemsep=2pt,topsep=2pt]
\item {} === SEARCH: ``[query used]'' ===
\item {} DATE\_\allowbreak{}RUN: [today's date]
\item {} CATEGORY: [estimates|analyst\_\allowbreak{}ratings|risks|news|sector]
\item {} FINDING\_\allowbreak{}1: [key finding or excerpt]
\item {} URL\_\allowbreak{}1: [source URL from search result]
\item {} SOURCE\_\allowbreak{}1: [publication name, date if available]
\item {} FINDING\_\allowbreak{}2: [key finding or excerpt]
\item {} URL\_\allowbreak{}2: [source URL]
\item {} SOURCE\_\allowbreak{}2: [publication name, date]
\item {} [...continue for all relevant results from this search...]
\end{itemize}

\textbf{Earnings estimates \& analyst sentiment:}

\begin{enumerate}[leftmargin=*,itemsep=2pt,topsep=2pt]
\item {} \emph{search} for ``[TICKER] earnings estimates consensus EPS revenue upcoming quarter''
\begin{itemize}[leftmargin=*,itemsep=2pt,topsep=2pt]
\item {} Record: consensus EPS, consensus revenue, estimate revision direction over last 90 days.
\item {} \textbf{Append to kensho-findings.txt immediately.}
\end{itemize}
\item {} \emph{search} for ``[TICKER] analyst ratings price target upgrades downgrades''
\begin{itemize}[leftmargin=*,itemsep=2pt,topsep=2pt]
\item {} Record: recent upgrades/downgrades, price target range, bull/bear thesis summaries.
\item {} \textbf{Append to kensho-findings.txt immediately.}
\end{itemize}
\item {} \emph{search} for ``[TICKER] risks bear case concerns investors''
\begin{itemize}[leftmargin=*,itemsep=2pt,topsep=2pt]
\item {} Record: key debates, bear arguments, swing factors for the upcoming print.
\item {} \textbf{Append to kensho-findings.txt immediately.}
\end{itemize}
\end{enumerate}

\textbf{Recent news (MANDATORY — do not skip):}

\begin{enumerate}[leftmargin=*,itemsep=2pt,topsep=2pt]
\item {} \emph{search} for ``[TICKER] [company name] recent news developments''
\begin{itemize}[leftmargin=*,itemsep=2pt,topsep=2pt]
\item {} Record: material news from the last 60 days — M\&A, product launches, executive changes, regulatory actions, partnerships, legal developments, tariffs, or any event that could affect the upcoming earnings print or forward guidance.
\item {} For each item, note the date, headline, potential earnings impact.
\item {} \textbf{Append to kensho-findings.txt immediately.}
\end{itemize}
\end{enumerate}

\textbf{Sector context:}

\begin{enumerate}[leftmargin=*,itemsep=2pt,topsep=2pt]
\item {} \emph{search} for ``[company industry/sector] sector outlook trends''
\begin{itemize}[leftmargin=*,itemsep=2pt,topsep=2pt]
\item {} Record: sector-level tailwinds/headwinds, macro data, competitive dynamics.
\item {} \textbf{Append to kensho-findings.txt immediately.}
\end{itemize}
\end{enumerate}

\vspace{0.5em}\hrule\vspace{0.5em}

\subparagraph{Phase 5: Financial Data Collection}\mbox{}\par

\textbf{Quarterly financials (last 8 quarters):} \emph{get\_\allowbreak{}financial\_\allowbreak{}line\_\allowbreak{}item\_\allowbreak{}from\_\allowbreak{}identifiers} with \emph{period\_\allowbreak{}type=``quarterly''}, \emph{num\_\allowbreak{}periods=8} for: \emph{revenue}, \emph{gross\_\allowbreak{}profit}, \emph{operating\_\allowbreak{}income}, \emph{ebitda}, \emph{net\_\allowbreak{}income}, \emph{diluted\_\allowbreak{}eps}


\textbf{After each line item call returns, immediately append to} \emph{/tmp/earnings-preview/financials.csv}. Write the raw values exactly as returned — do NOT round or convert yet. Include the \emph{source} column with the exact MCP function call and parameters:

\textbf{Structured Template}
\begin{itemize}[leftmargin=*,itemsep=2pt,topsep=2pt]
\item {} ticker,period,line\_\allowbreak{}item,value,source
\item {} D,Q4 2024,revenue,3941000000,get\_\allowbreak{}financial\_\allowbreak{}line\_\allowbreak{}item\_\allowbreak{}from\_\allowbreak{}identifiers(identifier='D',line\_\allowbreak{}item='revenue',period\_\allowbreak{}type='quarterly')
\item {} D,Q1 2025,revenue,3400000000,get\_\allowbreak{}financial\_\allowbreak{}line\_\allowbreak{}item\_\allowbreak{}from\_\allowbreak{}identifiers(identifier='D',line\_\allowbreak{}item='revenue',period\_\allowbreak{}type='quarterly')
\item {} D,Q2 2025,revenue,4076000000,get\_\allowbreak{}financial\_\allowbreak{}line\_\allowbreak{}item\_\allowbreak{}from\_\allowbreak{}identifiers(identifier='D',line\_\allowbreak{}item='revenue',period\_\allowbreak{}type='quarterly')
\item {} D,Q3 2025,revenue,3810000000,get\_\allowbreak{}financial\_\allowbreak{}line\_\allowbreak{}item\_\allowbreak{}from\_\allowbreak{}identifiers(identifier='D',line\_\allowbreak{}item='revenue',period\_\allowbreak{}type='quarterly')
\item {} D,Q4 2024,diluted\_\allowbreak{}eps,1.09,get\_\allowbreak{}financial\_\allowbreak{}line\_\allowbreak{}item\_\allowbreak{}from\_\allowbreak{}identifiers(identifier='D',line\_\allowbreak{}item='diluted\_\allowbreak{}eps',period\_\allowbreak{}type='quarterly')
\item {} D,Q1 2025,diluted\_\allowbreak{}eps,-0.11,get\_\allowbreak{}financial\_\allowbreak{}line\_\allowbreak{}item\_\allowbreak{}from\_\allowbreak{}identifiers(identifier='D',line\_\allowbreak{}item='diluted\_\allowbreak{}eps',period\_\allowbreak{}type='quarterly')
\item {} ...
\end{itemize}

\textbf{Do NOT calculate margins or growth rates yet.} Write raw data only. Calculations happen in Phase 6.


\textbf{Segment data:}

\begin{itemize}[leftmargin=*,itemsep=2pt,topsep=2pt]
\item {} \emph{get\_\allowbreak{}segments\_\allowbreak{}from\_\allowbreak{}identifiers} with \emph{segment\_\allowbreak{}type=``business''}, \emph{period\_\allowbreak{}type=``quarterly''}, \emph{num\_\allowbreak{}periods=8}
\item {} You need 8 quarters (not 4) so you have the year-ago quarter for y/y comparisons. To calculate y/y for Q3 2025, you need Q3 2024 — which is the 5th quarter back. \textbf{If the prior-year quarter's segment data is not available in the API response, do NOT estimate or fabricate it. State ``y/y not available'' in the report.}
\end{itemize}

\textbf{Immediately write} \emph{/tmp/earnings-preview/segments.csv}:

\textbf{Structured Template}
\begin{itemize}[leftmargin=*,itemsep=2pt,topsep=2pt]
\item {} ticker,period,segment\_\allowbreak{}name,revenue,source
\item {} D,Q3 2024,Dominion Energy Virginia,2762000000,get\_\allowbreak{}segments\_\allowbreak{}from\_\allowbreak{}identifiers(identifier='D',segment\_\allowbreak{}type='business',period\_\allowbreak{}type='quarterly')
\item {} D,Q3 2024,Dominion Energy South Carolina,848000000,get\_\allowbreak{}segments\_\allowbreak{}from\_\allowbreak{}identifiers(identifier='D',segment\_\allowbreak{}type='business',period\_\allowbreak{}type='quarterly')
\item {} D,Q3 2024,Contracted Energy,260000000,get\_\allowbreak{}segments\_\allowbreak{}from\_\allowbreak{}identifiers(identifier='D',segment\_\allowbreak{}type='business',period\_\allowbreak{}type='quarterly')
\item {} D,Q3 2025,Dominion Energy Virginia,3311000000,get\_\allowbreak{}segments\_\allowbreak{}from\_\allowbreak{}identifiers(identifier='D',segment\_\allowbreak{}type='business',period\_\allowbreak{}type='quarterly')
\item {} D,Q3 2025,Dominion Energy South Carolina,945000000,get\_\allowbreak{}segments\_\allowbreak{}from\_\allowbreak{}identifiers(identifier='D',segment\_\allowbreak{}type='business',period\_\allowbreak{}type='quarterly')
\item {} D,Q3 2025,Contracted Energy,297000000,get\_\allowbreak{}segments\_\allowbreak{}from\_\allowbreak{}identifiers(identifier='D',segment\_\allowbreak{}type='business',period\_\allowbreak{}type='quarterly')
\item {} ...
\end{itemize}

\textbf{Earnings history (for stock chart annotations):}

\begin{itemize}[leftmargin=*,itemsep=2pt,topsep=2pt]
\item {} \emph{get\_\allowbreak{}earnings\_\allowbreak{}from\_\allowbreak{}identifiers} — collect past earnings dates within the 12-month price window.
\item {} \textbf{Immediately write} \emph{/tmp/earnings-preview/earnings-dates.csv}:
\end{itemize}
\textbf{Structured Template}
\begin{itemize}[leftmargin=*,itemsep=2pt,topsep=2pt]
\item {} ticker,earnings\_\allowbreak{}date,call\_\allowbreak{}name,source
\item {} D,2025-05-02,Q1 2025 Earnings Call,get\_\allowbreak{}earnings\_\allowbreak{}from\_\allowbreak{}identifiers(identifier='D')
\item {} D,2025-08-01,Q2 2025 Earnings Call,get\_\allowbreak{}earnings\_\allowbreak{}from\_\allowbreak{}identifiers(identifier='D')
\item {} D,2025-10-31,Q3 2025 Earnings Call,get\_\allowbreak{}earnings\_\allowbreak{}from\_\allowbreak{}identifiers(identifier='D')
\item {} ...
\end{itemize}

\vspace{0.5em}\hrule\vspace{0.5em}

\subparagraph{Phase 6: Verification \& Calculations (MANDATORY — DO NOT SKIP)}\mbox{}\par

Before generating the report, read back ALL intermediate files and perform calculations from the clean data. This phase ensures data integrity by working from files rather than compressed conversation context.


\begin{enumerate}[leftmargin=*,itemsep=2pt,topsep=2pt]
\item {} \textbf{Read all intermediate files} using bash \emph{cat} commands:
\begin{itemize}[leftmargin=*,itemsep=2pt,topsep=2pt]
\item {} \emph{cat /tmp/earnings-preview/company-info.txt}
\item {} \emph{cat /tmp/earnings-preview/transcript-extracts.txt}
\item {} \emph{cat /tmp/earnings-preview/financials.csv}
\item {} \emph{cat /tmp/earnings-preview/segments.csv}
\item {} \emph{cat /tmp/earnings-preview/prices.csv}
\item {} \emph{cat /tmp/earnings-preview/peer-eps.csv}
\item {} \emph{cat /tmp/earnings-preview/peer-market-caps.csv}
\item {} \emph{cat /tmp/earnings-preview/consensus-eps.csv}
\item {} \emph{cat /tmp/earnings-preview/kensho-findings.txt}
\item {} \emph{cat /tmp/earnings-preview/earnings-dates.csv}
\end{itemize}
\item {} \textbf{Calculate derived metrics} from the raw data now in context:
\begin{itemize}[leftmargin=*,itemsep=2pt,topsep=2pt]
\item {} Gross margin \% = gross\_\allowbreak{}profit / revenue (per quarter)
\item {} Operating margin \% = operating\_\allowbreak{}income / revenue (per quarter)
\item {} Revenue y/y growth \% = (current Q revenue - year-ago Q revenue) / year-ago Q revenue
\item {} EPS y/y growth \% = same logic; use ``n.m.'' if base is negative
\item {} Segment y/y growth \% = match segment by name to year-ago Q; if missing, note ``y/y not available''
\item {} LTM P/E per company = latest price / sum of most recent 4 quarterly EPS (check which 4 quarters are available per ticker using \emph{peer-eps.csv})
\item {} NTM P/E per company = latest price / NTM EPS, where \textbf{NTM EPS = sum of the next 4 quarterly consensus mean EPS estimates} from \emph{consensus-eps.csv}. Add all 4 quarters' consensus\_\allowbreak{}mean\_\allowbreak{}eps values for each ticker. If fewer than 4 forward quarters are available for a peer, mark NTM P/E as ``n/a''. Note in the appendix which 4 quarters were summed.
\item {} Stock returns (YTD, 1-yr, 30d, 90d) = find the \textbf{common first date across all tickers} in \emph{prices.csv}, then compute returns from that date
\end{itemize}
\item {} \textbf{Cross-check}:
\begin{itemize}[leftmargin=*,itemsep=2pt,topsep=2pt]
\item {} Verify every segment y/y has the actual prior-year row in \emph{segments.csv}. If not, mark ``y/y not available.''
\item {} Verify all stock return base dates are identical across tickers.
\item {} Verify any multi-step calculation by re-summing components (e.g., LTM EPS sum matches the 4 quarterly values).
\item {} Verify all verbatim quotes in \emph{transcript-extracts.txt} are exact copy-pastes (not paraphrases).
\end{itemize}
\item {} \textbf{Write} \emph{/tmp/earnings-preview/calculations.csv} with all derived values:
\end{enumerate}
\textbf{Structured Template}
\begin{itemize}[leftmargin=*,itemsep=2pt,topsep=2pt]
\item {} ticker,metric,value,formula,components
\item {} D,gross\_\allowbreak{}margin\_\allowbreak{}Q3\_\allowbreak{}2025,32.5\%,gross\_\allowbreak{}profit/revenue,``gross\_\allowbreak{}profit=1238100000,revenue=3810000000''
\item {} D,revenue\_\allowbreak{}yoy\_\allowbreak{}Q3\_\allowbreak{}2025,+9.3\%,(Q3\_\allowbreak{}2025-Q3\_\allowbreak{}2024)/Q3\_\allowbreak{}2024,``Q3\_\allowbreak{}2025=3810000000,Q3\_\allowbreak{}2024=3486000000''
\item {} D,ltm\_\allowbreak{}pe,24.2x,price/ltm\_\allowbreak{}eps,``price=65.46,ltm\_\allowbreak{}eps=2.70,quarters=Q4\_\allowbreak{}2024+Q1\_\allowbreak{}2025+Q2\_\allowbreak{}2025+Q3\_\allowbreak{}2025''
\item {} D,ntm\_\allowbreak{}pe,18.5x,price/ntm\_\allowbreak{}eps,``price=65.46,ntm\_\allowbreak{}eps=3.56,quarters=Q4\_\allowbreak{}2025(0.88)+Q1\_\allowbreak{}2026(0.72)+Q2\_\allowbreak{}2026(0.91)+Q3\_\allowbreak{}2026(1.05),source=get\_\allowbreak{}consensus\_\allowbreak{}estimates\_\allowbreak{}from\_\allowbreak{}identifiers''
\item {} D,yoy\_\allowbreak{}return,+17.6\%,(end-start)/start,``end=65.46,start=55.67,base\_\allowbreak{}date=2025-02-19''
\item {} DUK,yoy\_\allowbreak{}return,+13.0\%,(end-start)/start,``end=126.32,start=111.79,base\_\allowbreak{}date=2025-02-19''
\item {} ...
\end{itemize}

This file becomes the single source of truth for all numbers in the report.


\vspace{0.5em}\hrule\vspace{0.5em}

\subparagraph{Phase 7: Generate the HTML Report}\mbox{}\par

\textbf{STOP — BEFORE WRITING ANY HTML, YOU MUST READ ALL INTERMEDIATE FILES. THIS IS A BLOCKING PREREQUISITE.}


This is not optional. You MUST run each \emph{cat} command below as a \textbf{separate bash tool call} (not combined into one). This ensures each file's contents are individually loaded and visible in the conversation. Do NOT combine them into a single command. Do NOT skip any file.


Run these commands \textbf{one at a time, each as its own bash call}:


\begin{enumerate}[leftmargin=*,itemsep=2pt,topsep=2pt]
\item {} \emph{cat /tmp/earnings-preview/company-info.txt}
\item {} \emph{cat /tmp/earnings-preview/transcript-extracts.txt}
\item {} \emph{cat /tmp/earnings-preview/financials.csv}
\item {} \emph{cat /tmp/earnings-preview/segments.csv}
\item {} \emph{cat /tmp/earnings-preview/prices.csv}
\item {} \emph{cat /tmp/earnings-preview/peer-eps.csv}
\item {} \emph{cat /tmp/earnings-preview/peer-market-caps.csv}
\item {} \emph{cat /tmp/earnings-preview/consensus-eps.csv}
\item {} \emph{cat /tmp/earnings-preview/kensho-findings.txt}
\item {} \emph{cat /tmp/earnings-preview/earnings-dates.csv}
\item {} \emph{cat /tmp/earnings-preview/calculations.csv}
\end{enumerate}

\textbf{After reading ALL files, you MUST print a summary message to the user} that lists every file and its status. Use exactly this format:


\textbf{Structured Template}
\begin{itemize}[leftmargin=*,itemsep=2pt,topsep=2pt]
\item {} --- DATA FILE VERIFICATION ---
\item {} 1. company-info.txt ✓ loaded ([N] lines)
\item {} 2. transcript-extracts.txt ✓ loaded ([N] lines)
\item {} 3. financials.csv ✓ loaded ([N] rows)
\item {} 4. segments.csv ✓ loaded ([N] rows)
\item {} 5. prices.csv ✓ loaded ([N] rows)
\item {} 6. peer-eps.csv ✓ loaded ([N] rows)
\item {} 7. peer-market-caps.csv ✓ loaded ([N] rows)
\item {} 8. consensus-eps.csv ✓ loaded ([N] rows)
\item {} 9. kensho-findings.txt ✓ loaded ([N] lines)
\item {} 10. earnings-dates.csv ✓ loaded ([N] rows)
\item {} 11. calculations.csv ✓ loaded ([N] rows)
\item {} All intermediate data files loaded successfully.
\item {} Generating report using file data as the single source of truth.
\end{itemize}

If any file is missing or empty, STOP and tell the user which file failed. Do NOT proceed to generate the report with missing data.


\textbf{Every number, quote, source URL, and MCP function call reference in the HTML report must come from these files — not from your memory of earlier conversation turns.} The files are the single source of truth. Earlier conversation context may have been compressed or summarized and WILL contain errors if relied upon. If a data point is not in the files, it should not appear in the report.


See \href{report-template.md}{report-template.md} for the complete HTML template, CSS, and Chart.js configuration.


\textbf{MANDATORY — Use Template Helper Functions for Charts:} The report-template.md provides pre-built, debugged Chart.js helper functions. You MUST use these exact functions to create charts. Do NOT write custom inline Chart.js code. The helpers are:

\begin{itemize}[leftmargin=*,itemsep=2pt,topsep=2pt]
\item {} \emph{createRevEpsChart(canvasId, labels, revenueData, epsData, revLabel)} — for Figure 1
\item {} \emph{createMarginChart(canvasId, labels, grossMargins, opMargins)} — for Figure 2
\item {} \emph{createRevGrowthChart(canvasId, labels, growthData)} — for Figure 3
\item {} \emph{createAnnotatedPriceChart(canvasId, labels, prices, earningsDates, ticker)} — for Figure 5
\item {} \emph{createCompPerfChart(canvasId, labels, datasets)} — for Figure 6
\item {} \emph{createPEChart(canvasId, companies)} — for Figure 7
\end{itemize}

Each chart call MUST be in its own \emph{<script>} tag wrapped in a try-catch block. This ensures a bug in one chart does not prevent other charts from rendering. Example:

\textbf{Structured Template}
\begin{itemize}[leftmargin=*,itemsep=2pt,topsep=2pt]
\item {} <script>
\item {} try \{
\item {} createRevEpsChart('chart-rev-eps', [...], [...], [...], 'Revenue (\$B)');
\item {} \} catch(e) \{ console.error('Figure 1 error:', e); \}
\item {} </script>
\item {} <script>
\item {} try \{
\item {} createMarginChart('chart-margins', [...], [...], [...]);
\item {} \} catch(e) \{ console.error('Figure 2 error:', e); \}
\item {} </script>
\end{itemize}

\subparagraph{Report Structure (4-5 pages total)}\mbox{}\par

The report has two halves: \textbf{narrative} (pages 1-2) and \textbf{figures} (pages 3-5). Keep these tightly integrated.


\vspace{0.5em}\hrule\vspace{0.5em}

\textbf{AI DISCLAIMER (MANDATORY — must appear in 3 places):} You MUST include the following disclaimer text in the report HTML. This is not optional — the report is incomplete without it:


\begin{quote}
``\textbf{``Analysis is AI-generated — please confirm all outputs''}''
\end{quote}

It must appear in exactly these 3 locations:

\begin{enumerate}[leftmargin=*,itemsep=2pt,topsep=2pt]
\item {} \textbf{Header banner} — immediately before the cover header, as a centered yellow banner: \emph{<div class=``ai-disclaimer''>Analysis is AI-generated — please confirm all outputs</div>}
\item {} \textbf{Footer} — inside the page-footer div, as a prominent yellow banner: \emph{<div class=``footer-disclaimer''>Analysis is AI-generated — please confirm all outputs</div>}
\item {} \textbf{Appendix} — as the first line of the appendix section, before the table: \emph{<div class=``ai-disclaimer''>Analysis is AI-generated — please confirm all outputs</div>}
\end{enumerate}

\vspace{0.5em}\hrule\vspace{0.5em}

\textbf{PAGE 1: Cover \& Thesis}


\begin{itemize}[leftmargin=*,itemsep=2pt,topsep=2pt]
\item {} \textbf{AI disclaimer banner} (yellow, centered — see AI DISCLAIMER rule above)
\item {} \textbf{Header}: Company name (TICKER) | Industry | Report date
\item {} \textbf{Title}: Thematic, specific to the quarter (e.g., ``Walmart Inc. (WMT) Q4 FY2026 Earnings Preview: Holiday Harvest — Can Furner's First Print Confirm the \$1T Thesis?'')
\item {} \textbf{Executive thesis} (2-3 short paragraphs max with bullet points):
\begin{itemize}[leftmargin=*,itemsep=2pt,topsep=2pt]
\item {} What we expect from this print in 1-2 sentences
\item {} 4-6 bullet points covering: our EPS estimate vs consensus, guidance expectations, key metrics to watch, what would move the stock, key debates
\item {} Keep it direct and opinionated — take a view, don't hedge everything
\end{itemize}
\item {} \textbf{Key management quotes} from the most recent earnings call woven into the narrative where relevant. Do NOT put these under a separate heading. Integrate them naturally as supporting evidence for your thesis points. Format as indented blockquotes.
\end{itemize}

\vspace{0.5em}\hrule\vspace{0.5em}

\textbf{PAGE 2: Estimates, Themes \& News}


\begin{itemize}[leftmargin=*,itemsep=2pt,topsep=2pt]
\item {} \textbf{Consensus Estimates Table} (single table, labeled as a figure):
\begin{itemize}[leftmargin=*,itemsep=2pt,topsep=2pt]
\item {} Columns: Metric | Consensus | Our Estimate | y/y Change
\item {} Rows: Revenue, EPS, Gross Margin, Operating Income, and 2-3 company-specific KPIs that matter (e.g., comp sales, eComm growth, membership revenue — whatever the Street cares about for THIS company)
\item {} \textbf{Color-coding is strictly mechanical:} If the y/y change value is negative, use \emph{class=``neg''} (red). If positive, use \emph{class=``pos''} (green). If zero or N/A, use \emph{class=``neutral''}. The sign of the number determines the class — do NOT override based on interpretation. A -1.1\% is ALWAYS red, even if the decline is small.
\item {} This is the ONLY guidance/estimates section. Do not repeat estimate data elsewhere.
\end{itemize}
\item {} \textbf{Key Metrics Beyond Headline EPS} (bulleted list, 3-5 items):
\begin{itemize}[leftmargin=*,itemsep=2pt,topsep=2pt]
\item {} The specific metrics that will determine if this is a good or bad quarter beyond the EPS number
\item {} For each: what the metric is, what consensus/management expects, why it matters
\item {} Be specific: ``Walmart Connect ad revenue growth (consensus \textasciitilde{}30\% y/y, 3Q was 33\%)''
\end{itemize}
\item {} \textbf{Themes to Watch} (3-5 bullets):
\begin{itemize}[leftmargin=*,itemsep=2pt,topsep=2pt]
\item {} Forward-looking items for the upcoming report
\item {} What management needs to deliver on, what could surprise, what the bears are focused on
\item {} Each theme: 1-2 sentences max
\end{itemize}
\item {} \textbf{Recent News \& Developments} (3-5 bullets):
\begin{itemize}[leftmargin=*,itemsep=2pt,topsep=2pt]
\item {} Material news from the last 60 days, one line each
\item {} Date + headline + brief impact assessment
\item {} Only include items that could affect the upcoming print or guidance
\end{itemize}
\end{itemize}

\vspace{0.5em}\hrule\vspace{0.5em}

\textbf{PAGES 3-5: Figures (all charts and tables)}


All figures are numbered sequentially. Every figure has a title and source line.


\begin{itemize}[leftmargin=*,itemsep=2pt,topsep=2pt]
\item {} \textbf{Figure 1: Quarterly Revenue \& Diluted EPS} — Bar/line combo, 8 quarters
\item {} \textbf{Figure 2: Margin Trends (Gross \& Operating \%)} — Dual line chart, 8 quarters
\item {} \textbf{Figure 3: Revenue Growth y/y \%} — Bar chart with green/red conditional coloring. \textbf{Only include quarters where both current and year-ago data exist} (typically the most recent 4 quarters from the 8 fetched). Do NOT include quarters where y/y cannot be computed — the chart should have 4 bars, not 8.
\item {} \textbf{Figure 4: Business Segment Revenue} — Table: Segment | Latest Q Rev (\$M) | \% of Total | y/y Change
\item {} \textbf{Figure 5: 1-Year Stock Price with Earnings Dates} — Price line with vertical annotation lines at earnings dates, labeled with quarter and 1-day post-earnings move
\item {} \textbf{Figure 6: Stock Performance vs. Competitors (Indexed to 100)} — Multi-line chart, subject company as thick solid line, competitors as thinner dashed lines
\item {} \textbf{Figure 7: LTM P/E vs. Competitors} — Horizontal bar chart, subject company highlighted in navy
\item {} \textbf{Figure 8: Competitor Comparison Table} — Ticker | Company | Mkt Cap | LTM P/E | NTM P/E | YTD \% | 1-Yr \%
\end{itemize}

\vspace{0.5em}\hrule\vspace{0.5em}

\textbf{APPENDIX: Data Sources \& Calculations (MANDATORY — DO NOT SKIP OR ABBREVIATE)}


The appendix MUST begin with the AI disclaimer banner: \emph{<div class=``ai-disclaimer''>Analysis is AI-generated — please confirm all outputs</div>}


The final page(s) of the report MUST include an Appendix table that documents \textbf{every claim} — numeric and non-numeric — cited in the report. \textbf{Every number that appears in the report body must have a corresponding row in this appendix, and every such number in the report body must be a clickable @@TOKEN0@@ hyperlink that scrolls to its appendix row.} If a number appears in the report without a hyperlink to the appendix, the report is incomplete.


\begin{itemize}[leftmargin=*,itemsep=2pt,topsep=2pt]
\item {} \textbf{Table columns}: Ref \# | Fact | Value | Source \& Derivation
\item {} \textbf{Ref \#}: Sequential ID matching the hyperlink anchors in the report body (\emph{ref-1}, \emph{ref-2}, etc.). Each row has an \emph{id=``ref-N''} attribute so hyperlinks scroll to it.
\item {} \textbf{Fact}: Human-readable label (e.g., ``Q3 FY2026 Revenue'', ``LTM P/E — WMT'', ``Management flagged tariff headwinds'', ``Barclays upgraded to Overweight'')
\item {} \textbf{Value}: The exact figure as displayed in the report (e.g., ``\$152.3B'', ``24.5\%'', ``28.1x''). For non-numeric facts, leave blank or write ``N/A''.
\item {} \textbf{Source \& Derivation}: This is the critical column. \textbf{Every row must have a specific, detailed source — not just a label.} Follow these rules strictly: \textbf{For raw financial data from S\&P Capital IQ (revenue, EPS, gross profit, operating income, net income, EBITDA, prices, market cap, etc.):}
\begin{itemize}[leftmargin=*,itemsep=2pt,topsep=2pt]
\item {} State the MCP function used and its key parameters. Format: \emph{S\&P Capital IQ — @@TOKEN0@@}
\item {} Examples:
\begin{itemize}[leftmargin=*,itemsep=2pt,topsep=2pt]
\item {} \emph{S\&P Capital IQ — get\_\allowbreak{}financial\_\allowbreak{}line\_\allowbreak{}item\_\allowbreak{}from\_\allowbreak{}identifiers(identifier='WMT', line\_\allowbreak{}item='revenue', period\_\allowbreak{}type='quarterly', period='Q3 FY2026')}
\item {} \emph{S\&P Capital IQ — get\_\allowbreak{}financial\_\allowbreak{}line\_\allowbreak{}item\_\allowbreak{}from\_\allowbreak{}identifiers(identifier='WMT', line\_\allowbreak{}item='diluted\_\allowbreak{}eps', period\_\allowbreak{}type='quarterly', period='Q3 FY2026')}
\item {} \emph{S\&P Capital IQ — get\_\allowbreak{}prices\_\allowbreak{}from\_\allowbreak{}identifiers(identifier='WMT', periodicity='day')}
\item {} \emph{S\&P Capital IQ — get\_\allowbreak{}capitalization\_\allowbreak{}from\_\allowbreak{}identifiers(identifier='WMT', capitalization='market\_\allowbreak{}cap')}
\end{itemize}
\item {} \textbf{Do NOT just write ``S\&P Capital IQ'' with no detail.} The reader must know exactly which data point from which tool call produced this number. \textbf{For calculated values (margins, growth rates, P/E, returns, y/y changes):}
\item {} Show the full formula with \textbf{hyperlinked components} — each component must be an \emph{<a href=``\#ref-N''>} link back to the appendix row for that raw data point. This is critical: the reader must be able to click through from the calculated value to each of its inputs.
\item {} Example: \emph{Gross Margin = <a href='\#ref-5'>Gross Profit \$37.2B</a> / <a href='\#ref-1'>Revenue \$152.3B</a> = 24.4\%. Source: S\&P Capital IQ (calculated)}
\item {} Example: \emph{LTM P/E = <a href='\#ref-20'>Price \$172.35</a> / (<a href='\#ref-8'>Q1 EPS \$1.47</a> + <a href='\#ref-9'>Q2 EPS \$1.84</a> + <a href='\#ref-10'>Q3 EPS \$1.53</a> + <a href='\#ref-11'>Q4 EPS \$1.80</a>) = \$172.35 / \$6.64 = 25.9x}
\item {} Example: \emph{Revenue y/y growth = (<a href='\#ref-12'>Q3 FY26 Rev \$165.8B</a> - <a href='\#ref-3'>Q3 FY25 Rev \$160.8B</a>) / <a href='\#ref-3'>Q3 FY25 Rev \$160.8B</a> = +3.1\%}
\item {} \textbf{Every formula component must be a clickable hyperlink.} Do NOT write formulas with plain-text numbers. \textbf{For transcript-sourced claims (quotes, management commentary, guidance):}
\item {} Write the \textbf{verbatim excerpt sentence} from the transcript.
\item {} Reference the transcript by its full name and the \emph{key\_\allowbreak{}dev\_\allowbreak{}id} used to fetch it.
\item {} Format: \emph{``[verbatim quote]'' — [Speaker], [Title]. Source: [Q\# FY\#\#\#\# Earnings Call Transcript] (key\_\allowbreak{}dev\_\allowbreak{}id: [ID])}
\item {} Example: \emph{``We expect comp sales growth of 3-4\% in Q4'' — CEO John Furner. Source: Q3 FY2026 Earnings Call Transcript (key\_\allowbreak{}dev\_\allowbreak{}id: 12345678)} \textbf{For Kensho Grounding search results (news, analyst ratings, consensus estimates):}
\item {} Write the key finding or excerpt from the search result.
\item {} \textbf{MANDATORY: Include the source URL} returned by the Kensho \emph{search} tool as a clickable \emph{<a href=``[URL]'' target=``\_\allowbreak{}blank''>} hyperlink. This is the most important part — readers must be able to click through to the original source.
\item {} Format: \emph{``[finding/excerpt]'' — <a href=``[URL]'' target=``\_\allowbreak{}blank''>[Source Title or Publication]</a>. Query: search(``[query used]'')}
\item {} Example: \emph{``Barclays upgraded WMT to Overweight with \$210 price target on Jan 15, 2026.'' — <a href=``https://www.investing.com/news/barclays-upgrades-wmt'' target=``\_\allowbreak{}blank''>Investing.com, Jan 15 2026</a>. Query: search(``WMT analyst ratings price target upgrades downgrades'')}
\item {} If no URL was returned for a particular result, write ``Source URL not available'' and still include the search query.
\end{itemize}
\end{itemize}

\textbf{Completeness check:} Before finalizing the report, scan every number in the report body. If any number is not wrapped in \emph{<a href=``\#ref-N'' class=``data-ref''>}, fix it. If any appendix row has a Source \& Derivation that is just a bare label like ``S\&P Capital IQ'' with no function call detail, fix it. If any calculated value's formula lacks hyperlinked components, fix it. If any Kensho-sourced claim lacks a source URL, fix it.


Group the appendix rows by section (Financials, Valuation, Estimates \& Consensus, Transcript Claims, News \& Analyst Commentary, Stock Performance) with subheadings. Use smaller font size (10-11px).


\vspace{0.5em}\hrule\vspace{0.5em}

\subparagraph{Phase 8: Output}\mbox{}\par

\begin{enumerate}[leftmargin=*,itemsep=2pt,topsep=2pt]
\item {} Write the complete HTML file to \emph{earnings-preview-[TICKER]-YYYY-MM-DD.html} in the current working directory.
\item {} Open it in the browser: \emph{open earnings-preview-[TICKER]-YYYY-MM-DD.html}
\item {} Tell the user the file has been created and summarize the key findings.
\end{enumerate}

\vspace{0.5em}\hrule\vspace{0.5em}

\subparagraph{Writing Guidelines}\mbox{}\par

\begin{itemize}[leftmargin=*,itemsep=2pt,topsep=2pt]
\item {} \textbf{NO EMOJIS}: Do not use any emojis anywhere in the report. This is a professional research document.
\item {} \textbf{CONCISE}: Target 4-5 printed pages. Every sentence must carry weight. Use bullets, not paragraphs, wherever possible. If a section feels long, cut it.
\item {} \textbf{Be specific with numbers}: ``\$52.4B revenue, up 5.2\% y/y'' not ``strong revenue growth.''
\item {} \textbf{Take a view}: This is an earnings preview, not a summary. State what you expect, what matters, and why. Be opinionated but back it with data.
\item {} \textbf{Management quotes without headers}: Weave 3-4 key management quotes from the most recent call directly into the narrative as blockquotes. Do not create a ``Key Management Quotes'' section heading — let them flow naturally as supporting evidence.
\item {} \textbf{Professional tone}: Sell-side equity research style — analytical, direct, data-driven.
\item {} \textbf{Charts must use real data}: Every chart populated with actual MCP data. Never fabricate.
\item {} \textbf{Competitor context}: Frame valuation relative to peers. A 25x P/E means nothing without knowing peers trade at 20x or 35x.
\item {} \textbf{Hyperlinked claims}: Every factual claim — numeric or qualitative — must be an \emph{<a class=``data-ref''>} tag linking to its appendix entry. Numbers: \emph{<a href=``\#ref-1'' class=``data-ref''>\$152.3B</a>}. Qualitative: \emph{<a href=``\#ref-25'' class=``data-ref''>management flagged tariff headwinds as the primary margin risk</a>}. No fact should appear without a traceable source in the appendix.
\end{itemize}

\vspace{0.8em}\hrule\vspace{0.8em}

\subsection{Skill Resource: spglobal/skills/earnings-preview-beta/report-template.md}\label{file:spglobal-skills-earnings-preview-beta-report-template.md}

\paragraph{HTML Report Template Reference}\mbox{}\par

Use this template as the foundation for the single-company earnings preview HTML report. Customize the data, charts, and narrative content based on the research gathered in Phases 1-5.


\subparagraph{HTML Structure}\mbox{}\par

The report is a single self-contained HTML file with:

\begin{itemize}[leftmargin=*,itemsep=2pt,topsep=2pt]
\item {} Embedded CSS (no external stylesheets)
\item {} Chart.js loaded from CDN for interactive charts
\item {} Print-friendly styles via \emph{@media print}
\item {} Responsive layout that works on screens and in print
\item {} Target: 4-5 printed pages
\end{itemize}

\subparagraph{Complete Template}\mbox{}\par

\textbf{Structured Template}
\begin{itemize}[leftmargin=*,itemsep=2pt,topsep=2pt]
\item {} <!DOCTYPE html>
\item {} <html lang=``en''>
\item {} <head>
\item {} <meta charset=``UTF-8''>
\item {} <meta name=``viewport'' content=``width=device-width, initial-scale=1.0''>
\item {} <title>Earnings Preview — [COMPANY] ([TICKER]) — [DATE]</title>
\item {} <script src=``https://cdn.jsdelivr.net/npm/chart.js@4.4.7/dist/chart.umd.min.js'' integrity=``sha384-vsrfeLOOY6KuIYKDlmVH5UiBmgIdB1oEf7p01YgWHuqmOHfZr374+odEv96n9tNC'' crossorigin=``anonymous''></script>
\item {} <script src=``https://cdn.jsdelivr.net/npm/chartjs-plugin-annotation@3.1.0/dist/chartjs-plugin-annotation.min.js'' integrity=``sha384-3N9GHhCtN3CQef6tNfqgZlv7sQLYIkcChN+uaTZ7xVdzKYp/SjBNPxa92+hM7EAY'' crossorigin=``anonymous''></script>
\item {} <style>
\item {} /\emph{ Reset \& Base }/
\item {} \emph{, }::before, *::after \{ box-sizing: border-box; margin: 0; padding: 0; \}
\item {} html \{ font-size: 15px; \}
\item {} body \{
\item {} font-family: 'Arial Narrow', Arial, sans-serif;
\item {} color: \#1a1a2e;
\item {} background: \#fff;
\item {} line-height: 1.6;
\item {} \}
\item {} /\emph{ Layout }/
\item {} .page \{
\item {} max-width: 1100px;
\item {} margin: 0 auto;
\item {} padding: 40px 48px;
\item {} \}
\item {} .page-break \{
\item {} page-break-before: always;
\item {} border-top: 2px solid \#1a1a4e;
\item {} margin-top: 48px;
\item {} padding-top: 32px;
\item {} \}
\item {} /\emph{ Header / Cover }/
\item {} .cover-header \{
\item {} border-bottom: 3px solid \#1a1a4e;
\item {} padding-bottom: 16px;
\item {} margin-bottom: 24px;
\item {} display: flex;
\item {} justify-content: space-between;
\item {} align-items: flex-end;
\item {} \}
\item {} .cover-header .brand \{
\item {} font-size: 24px;
\item {} font-weight: bold;
\item {} color: \#1a1a4e;
\item {} letter-spacing: 2px;
\item {} text-transform: uppercase;
\item {} \}
\item {} .cover-header .sector \{
\item {} font-size: 13px;
\item {} color: \#555;
\item {} \}
\item {} .cover-header .date \{
\item {} font-size: 14px;
\item {} color: \#333;
\item {} text-align: right;
\item {} \}
\item {} .report-title \{
\item {} font-size: 26px;
\item {} font-weight: bold;
\item {} color: \#1a1a2e;
\item {} margin: 20px 0 16px 0;
\item {} line-height: 1.3;
\item {} \}
\item {} /\emph{ Executive Thesis }/
\item {} .executive-summary \{
\item {} font-size: 14px;
\item {} line-height: 1.65;
\item {} color: \#222;
\item {} margin-bottom: 16px;
\item {} \}
\item {} .executive-summary p \{
\item {} margin-bottom: 10px;
\item {} text-align: justify;
\item {} \}
\item {} .executive-summary ul \{
\item {} margin: 8px 0 10px 20px;
\item {} font-size: 13.5px;
\item {} \}
\item {} .executive-summary ul li \{
\item {} margin-bottom: 5px;
\item {} line-height: 1.5;
\item {} \}
\item {} blockquote \{
\item {} border-left: 3px solid \#b0b8c8;
\item {} padding: 6px 14px;
\item {} margin: 8px 0 8px 12px;
\item {} font-style: italic;
\item {} color: \#444;
\item {} background: \#f9fafb;
\item {} font-size: 12.5px;
\item {} line-height: 1.5;
\item {} \}
\item {} /\emph{ Section Headings }/
\item {} h2.section-title \{
\item {} font-size: 18px;
\item {} font-weight: 700;
\item {} color: \#1a1a4e;
\item {} border-bottom: 2px solid \#1a1a4e;
\item {} padding-bottom: 5px;
\item {} margin: 28px 0 14px 0;
\item {} \}
\item {} h3.subsection-title \{
\item {} font-size: 14px;
\item {} font-weight: 600;
\item {} color: \#1a1a4e;
\item {} margin: 16px 0 8px 0;
\item {} \}
\item {} h4.figure-title \{
\item {} font-size: 12px;
\item {} font-weight: 600;
\item {} color: \#444;
\item {} margin: 14px 0 6px 0;
\item {} \}
\item {} /\emph{ Tables }/
\item {} table \{
\item {} width: 100\%;
\item {} border-collapse: collapse;
\item {} font-size: 12px;
\item {} margin: 10px 0 16px 0;
\item {} \}
\item {} thead th \{
\item {} background: \#1a1a4e;
\item {} color: \#fff;
\item {} padding: 7px 10px;
\item {} text-align: left;
\item {} font-weight: 600;
\item {} font-size: 11px;
\item {} text-transform: uppercase;
\item {} letter-spacing: 0.5px;
\item {} \}
\item {} tbody td \{
\item {} padding: 6px 10px;
\item {} border-bottom: 1px solid \#e0e0e0;
\item {} \}
\item {} tbody tr:nth-child(even) \{
\item {} background: \#f9fafb;
\item {} \}
\item {} tbody tr:hover \{
\item {} background: \#eef0f5;
\item {} \}
\item {} .num \{ text-align: right; font-variant-numeric: tabular-nums; \}
\item {} .pos \{ color: \#0d7a3e; font-weight: 600; \}
\item {} .neg \{ color: \#c0392b; font-weight: 600; \}
\item {} .neutral \{ color: \#555; \}
\item {} .highlight-row \{ background: \#e8eaf6 !important; font-weight: 600; \}
\item {} /\emph{ Chart Containers }/
\item {} .chart-row \{
\item {} display: grid;
\item {} grid-template-columns: 1fr 1fr;
\item {} gap: 20px;
\item {} margin: 12px 0 20px 0;
\item {} \}
\item {} .chart-container \{
\item {} position: relative;
\item {} background: \#fafbfc;
\item {} border: 1px solid \#e8e8e8;
\item {} border-radius: 4px;
\item {} padding: 14px;
\item {} \}
\item {} .chart-container canvas \{
\item {} max-height: 260px;
\item {} \}
\item {} .chart-full \{
\item {} grid-column: 1 / -1;
\item {} \}
\item {} /\emph{ Compact Lists }/
\item {} .key-metrics ul, .themes ul, .news-list ul \{
\item {} margin: 6px 0 6px 18px;
\item {} font-size: 13px;
\item {} line-height: 1.55;
\item {} \}
\item {} .key-metrics li, .themes li, .news-list li \{
\item {} margin-bottom: 5px;
\item {} \}
\item {} /\emph{ Data Reference Links }/
\item {} a.data-ref \{
\item {} color: \#1a1a4e;
\item {} text-decoration: none;
\item {} border-bottom: 1px dotted transparent;
\item {} transition: border-color 0.15s;
\item {} \}
\item {} a.data-ref:hover \{
\item {} border-bottom-color: \#1a1a4e;
\item {} \}
\item {} /\emph{ Appendix }/
\item {} .appendix table \{
\item {} font-size: 10.5px;
\item {} \}
\item {} .appendix thead th \{
\item {} font-size: 10px;
\item {} padding: 5px 8px;
\item {} \}
\item {} .appendix tbody td \{
\item {} padding: 4px 8px;
\item {} font-size: 10.5px;
\item {} vertical-align: top;
\item {} line-height: 1.45;
\item {} \}
\item {} .appendix .ref-id \{
\item {} font-weight: 600;
\item {} color: \#1a1a4e;
\item {} white-space: nowrap;
\item {} \}
\item {} .appendix .source-detail \{
\item {} font-size: 10px;
\item {} color: \#444;
\item {} \}
\item {} .appendix .source-detail .formula \{
\item {} font-family: 'Courier New', monospace;
\item {} font-size: 9.5px;
\item {} color: \#555;
\item {} \}
\item {} .appendix .source-detail .excerpt \{
\item {} font-style: italic;
\item {} color: \#555;
\item {} \}
\item {} .appendix .source-detail .src-label \{
\item {} font-weight: 600;
\item {} color: \#1a1a4e;
\item {} font-size: 9.5px;
\item {} \}
\item {} .appendix .source-detail a.data-ref \{
\item {} font-weight: 600;
\item {} \}
\item {} .appendix a.src-url \{
\item {} color: \#3366cc;
\item {} text-decoration: underline;
\item {} font-size: 10px;
\item {} word-break: break-all;
\item {} \}
\item {} .appendix a.src-url:hover \{
\item {} color: \#1a1a4e;
\item {} \}
\item {} .appendix .transcript-ref \{
\item {} font-weight: 600;
\item {} color: \#1a1a4e;
\item {} font-size: 10px;
\item {} \}
\item {} .appendix-group \{
\item {} font-size: 11px;
\item {} font-weight: 700;
\item {} color: \#1a1a4e;
\item {} background: \#f0f1f5;
\item {} padding: 4px 8px;
\item {} text-transform: uppercase;
\item {} letter-spacing: 0.5px;
\item {} \}
\item {} /\emph{ Source / Footer }/
\item {} .source \{
\item {} font-size: 10px;
\item {} color: \#999;
\item {} margin-top: 3px;
\item {} font-style: italic;
\item {} \}
\item {} .ai-disclaimer \{
\item {} background-color: \#fff3cd;
\item {} border: 1px solid \#ffc107;
\item {} border-radius: 4px;
\item {} padding: 4px 10px;
\item {} font-size: 11px;
\item {} font-weight: 600;
\item {} color: \#664d03;
\item {} text-align: center;
\item {} margin-bottom: 12px;
\item {} \}
\item {} .page-footer \{
\item {} border-top: 2px solid \#1a1a4e;
\item {} padding-top: 10px;
\item {} margin-top: 32px;
\item {} text-align: center;
\item {} \}
\item {} .page-footer .footer-disclaimer \{
\item {} font-size: 11px;
\item {} font-weight: 600;
\item {} color: \#664d03;
\item {} background-color: \#fff3cd;
\item {} border: 1px solid \#ffc107;
\item {} border-radius: 4px;
\item {} padding: 4px 10px;
\item {} display: inline-block;
\item {} margin-bottom: 4px;
\item {} \}
\item {} .page-footer .footer-meta \{
\item {} font-size: 10px;
\item {} color: \#888;
\item {} \}
\item {} /\emph{ Print Styles }/
\item {} @media print \{
\item {} body \{ font-size: 11px; \}
\item {} .page \{ padding: 16px; max-width: none; \}
\item {} .chart-container \{ break-inside: avoid; \}
\item {} table \{ break-inside: avoid; \}
\item {} .page-break \{ margin-top: 0; \}
\item {} .no-print \{ display: none; \}
\item {} \}
\item {} </style>
\item {} </head>
\item {} <body>
\item {} <div class=``page''>
\item {} <!-- -->
\item {} <!-- PAGE 1: COVER \& THESIS -->
\item {} <!-- -->
\item {} <div class=``ai-disclaimer''>Analysis is AI-generated — please confirm all outputs</div>
\item {} <div class=``cover-header''>
\item {} <div>
\item {} <div class=``brand''>Earnings Preview</div>
\item {} <div class=``sector''>[Industry] | [TICKER]</div>
\item {} </div>
\item {} <div class=``date''>[Full Date]</div>
\item {} </div>
\item {} <h1 class=``report-title''>[Company Name] ([TICKER]) [Q\# FY\#\#\#\#] Earnings Preview: [Thematic Subtitle]</h1>
\item {} <div class=``executive-summary''>
\item {} <!-- Executive thesis: 2-3 short paragraphs + bullet points.
\item {} What we expect, our EPS estimate vs consensus, guidance expectations,
\item {} key metrics to watch, what would move the stock, key debates.
\item {} Weave in 3-4 management quotes as blockquotes where they support the thesis.
\item {} Do NOT create a separate ``Key Management Quotes'' section. -->
\item {} <p>[Opening 1-2 sentences: what we expect from this print.]</p>
\item {} <ul>
\item {} <li><strong>EPS:</strong> We estimate <a href=``\#ref-1'' class=``data-ref''>\$X.XX</a> vs consensus <a href=``\#ref-2'' class=``data-ref''>\$X.XX</a>, [rationale]</li>
\item {} <li><strong>Revenue:</strong> We estimate <a href=``\#ref-3'' class=``data-ref''>\$XX.XB</a> vs consensus <a href=``\#ref-4'' class=``data-ref''>\$XX.XB</a>, [rationale]</li>
\item {} <li><strong>Guidance:</strong> [What to expect on forward guidance]</li>
\item {} <li><strong>Key metric:</strong> [Most important sub-headline metric to watch]</li>
\item {} <li><strong>Stock catalyst:</strong> [What would move the stock up/down post-print]</li>
\item {} <li><strong>Key debate:</strong> [What bulls and bears disagree on]</li>
\item {} </ul>
\item {} <blockquote>``[Key management quote supporting a thesis point]'' — [Speaker], [Q\# FY\#\#\#\#] Earnings Call</blockquote>
\item {} <p>[1-2 sentences tying it together — your overall read on the setup.]</p>
\item {} <blockquote>``[Another supporting quote]'' — [Speaker], [Q\# FY\#\#\#\#] Earnings Call</blockquote>
\item {} </div>
\item {} <!-- -->
\item {} <!-- PAGE 2: ESTIMATES, THEMES \& NEWS -->
\item {} <!-- -->
\item {} <div class=``page-break''>
\item {} <!-- Consensus Estimates Table (Figure label inline) -->
\item {} <h2 class=``section-title''>Consensus Estimates — [Q\# FY\#\#\#\#]</h2>
\item {} <h4 class=``figure-title''>[Q\# FY\#\#\#\#] Consensus Estimates</h4>
\item {} <table>
\item {} <thead>
\item {} <tr>
\item {} <th>Metric</th>
\item {} <th class=``num''>Consensus</th>
\item {} <th class=``num''>Our Estimate</th>
\item {} <th class=``num''>y/y Change</th>
\item {} </tr>
\item {} </thead>
\item {} <tbody>
\item {} <tr><td>Revenue</td><td class=``num''><a href=``\#ref-N'' class=``data-ref''>\$[XX.X]B</a></td><td class=``num''><a href=``\#ref-N'' class=``data-ref''>\$[XX.X]B</a></td><td class=``num [pos|neg]''><a href=``\#ref-N'' class=``data-ref''>[+/-X.X\%]</a></td></tr>
\item {} <tr><td>Diluted EPS</td><td class=``num''><a href=``\#ref-N'' class=``data-ref''>\$[X.XX]</a></td><td class=``num''><a href=``\#ref-N'' class=``data-ref''>\$[X.XX]</a></td><td class=``num [pos|neg]''><a href=``\#ref-N'' class=``data-ref''>[+/-X.X\%]</a></td></tr>
\item {} <tr><td>Gross Margin</td><td class=``num''><a href=``\#ref-N'' class=``data-ref''>[XX.X\%]</a></td><td class=``num''><a href=``\#ref-N'' class=``data-ref''>[XX.X\%]</a></td><td class=``num [pos|neg]''><a href=``\#ref-N'' class=``data-ref''>[+/-XXbps]</a></td></tr>
\item {} <tr><td>Operating Income</td><td class=``num''><a href=``\#ref-N'' class=``data-ref''>\$[X.X]B</a></td><td class=``num''><a href=``\#ref-N'' class=``data-ref''>\$[X.X]B</a></td><td class=``num [pos|neg]''><a href=``\#ref-N'' class=``data-ref''>[+/-X.X\%]</a></td></tr>
\item {} <!-- Add 2-3 company-specific KPIs below (e.g., comp sales, eComm growth, membership revenue) -->
\item {} <tr><td>[Company KPI 1]</td><td class=``num''>[Value]</td><td class=``num''>[Value]</td><td class=``num [pos|neg]''>[Change]</td></tr>
\item {} <tr><td>[Company KPI 2]</td><td class=``num''>[Value]</td><td class=``num''>[Value]</td><td class=``num [pos|neg]''>[Change]</td></tr>
\item {} </tbody>
\item {} </table>
\item {} <div class=``source''>Source: Kensho, S\&P Capital IQ</div>
\item {} <!-- Key Metrics Beyond Headline EPS -->
\item {} <h3 class=``subsection-title''>Key Metrics Beyond Headline EPS</h3>
\item {} <div class=``key-metrics''>
\item {} <ul>
\item {} <li><strong>[Metric 1]:</strong> [What consensus/management expects, why it matters. Be specific with numbers.]</li>
\item {} <li><strong>[Metric 2]:</strong> [Details]</li>
\item {} <li><strong>[Metric 3]:</strong> [Details]</li>
\item {} <!-- 3-5 items -->
\item {} </ul>
\item {} </div>
\item {} <!-- Themes to Watch -->
\item {} <h3 class=``subsection-title''>Themes to Watch</h3>
\item {} <div class=``themes''>
\item {} <ul>
\item {} <li><strong>[Theme 1]:</strong> [1-2 sentences max. Forward-looking, specific.]</li>
\item {} <li><strong>[Theme 2]:</strong> [Details]</li>
\item {} <li><strong>[Theme 3]:</strong> [Details]</li>
\item {} <!-- 3-5 themes -->
\item {} </ul>
\item {} </div>
\item {} <!-- Recent News \& Developments -->
\item {} <h3 class=``subsection-title''>Recent News \& Developments</h3>
\item {} <div class=``news-list''>
\item {} <ul>
\item {} <li><strong>[Date]:</strong> [Headline] — [Brief impact assessment, one line]</li>
\item {} <li><strong>[Date]:</strong> [Headline] — [Impact]</li>
\item {} <li><strong>[Date]:</strong> [Headline] — [Impact]</li>
\item {} <!-- 3-5 material items from last 60 days -->
\item {} </ul>
\item {} </div>
\item {} <div class=``source''>Source: Kensho</div>
\item {} </div>
\item {} <!-- -->
\item {} <!-- PAGES 3-5: FIGURES -->
\item {} <!-- All charts and tables, numbered sequentially -->
\item {} <!-- -->
\item {} <div class=``page-break''>
\item {} <h2 class=``section-title''>Financial \& Competitive Analysis</h2>
\item {} <!-- Figure 1: Quarterly Revenue \& Diluted EPS -->
\item {} <div class=``chart-row''>
\item {} <div class=``chart-container''>
\item {} <h4 class=``figure-title''>Figure 1: Quarterly Revenue \& Diluted EPS</h4>
\item {} <canvas id=``chart-rev-eps''></canvas>
\item {} <div class=``source''>Source: S\&P Capital IQ</div>
\item {} </div>
\item {} <!-- Figure 2: Margin Trends -->
\item {} <div class=``chart-container''>
\item {} <h4 class=``figure-title''>Figure 2: Margin Trends (Gross \& Operating \%)</h4>
\item {} <canvas id=``chart-margins''></canvas>
\item {} <div class=``source''>Source: S\&P Capital IQ</div>
\item {} </div>
\item {} </div>
\item {} <!-- Figure 3: Revenue Growth y/y \% -->
\item {} <div class=``chart-row''>
\item {} <div class=``chart-container chart-full''>
\item {} <h4 class=``figure-title''>Figure 3: Revenue Growth y/y (\%)</h4>
\item {} <canvas id=``chart-rev-growth'' style=``max-height: 200px;''></canvas>
\item {} <div class=``source''>Source: S\&P Capital IQ</div>
\item {} </div>
\item {} </div>
\item {} <!-- Figure 4: Business Segment Revenue -->
\item {} <h4 class=``figure-title''>Figure 4: Business Segment Revenue</h4>
\item {} <table>
\item {} <thead>
\item {} <tr>
\item {} <th>Segment</th>
\item {} <th class=``num''>Latest Q Rev (\$M)</th>
\item {} <th class=``num''>\% of Total</th>
\item {} <th class=``num''>y/y Change</th>
\item {} </tr>
\item {} </thead>
\item {} <tbody>
\item {} <!-- Populate from segment data. Color-code y/y change with pos/neg classes. -->
\item {} </tbody>
\item {} </table>
\item {} <div class=``source''>Source: S\&P Capital IQ</div>
\item {} </div>
\item {} <!-- Page break for stock \& competitor charts -->
\item {} <div class=``page-break''>
\item {} <!-- Figure 5: 1-Year Stock Price with Earnings Dates -->
\item {} <div class=``chart-row''>
\item {} <div class=``chart-container chart-full''>
\item {} <h4 class=``figure-title''>Figure 5: 1-Year Stock Price with Earnings Dates</h4>
\item {} <canvas id=``chart-price-annotated'' style=``max-height: 300px;''></canvas>
\item {} <div class=``source''>Source: S\&P Capital IQ</div>
\item {} </div>
\item {} </div>
\item {} <!-- Figure 6: Stock Performance vs. Competitors (Indexed to 100) -->
\item {} <div class=``chart-row''>
\item {} <div class=``chart-container chart-full''>
\item {} <h4 class=``figure-title''>Figure 6: Stock Performance vs. Competitors — 1 Year (Indexed to 100)</h4>
\item {} <canvas id=``chart-comp-perf'' style=``max-height: 300px;''></canvas>
\item {} <div class=``source''>Source: S\&P Capital IQ</div>
\item {} </div>
\item {} </div>
\item {} </div>
\item {} <div class=``page-break''>
\item {} <!-- Figure 7: LTM P/E vs. Competitors -->
\item {} <div class=``chart-row''>
\item {} <div class=``chart-container chart-full''>
\item {} <h4 class=``figure-title''>Figure 7: LTM P/E vs. Competitors</h4>
\item {} <canvas id=``chart-pe-comp'' style=``max-height: 280px;''></canvas>
\item {} <div class=``source''>Source: S\&P Capital IQ</div>
\item {} </div>
\item {} </div>
\item {} <!-- Figure 8: Competitor Comparison Table -->
\item {} <h4 class=``figure-title''>Figure 8: Competitor Comparison</h4>
\item {} <table>
\item {} <thead>
\item {} <tr>
\item {} <th>Ticker</th>
\item {} <th>Company</th>
\item {} <th class=``num''>Mkt Cap (\$B)</th>
\item {} <th class=``num''>LTM P/E</th>
\item {} <th class=``num''>NTM P/E</th>
\item {} <th class=``num''>YTD \%</th>
\item {} <th class=``num''>1-Yr \%</th>
\item {} </tr>
\item {} </thead>
\item {} <tbody>
\item {} <!-- Highlight the subject company row with class=``highlight-row'' -->
\item {} </tbody>
\item {} </table>
\item {} <div class=``source''>Source: S\&P Capital IQ</div>
\item {} </div>
\item {} <!-- -->
\item {} <!-- APPENDIX: DATA SOURCES \& CALCULATIONS -->
\item {} <!-- -->
\item {} <div class=``page-break appendix'' id=``appendix''>
\item {} <div class=``ai-disclaimer''>Analysis is AI-generated — please confirm all outputs</div>
\item {} <h2 class=``section-title''>Appendix: Data Sources \& Calculations</h2>
\item {} <p style=``font-size: 11px; color: \#666; margin-bottom: 12px;''>
\item {} Every claim in this report is hyperlinked to its entry below. Click any highlighted text to jump here.
\item {} </p>
\item {} <table>
\item {} <thead>
\item {} <tr>
\item {} <th style=``width: 40px;''>Ref</th>
\item {} <th style=``width: 170px;''>Fact</th>
\item {} <th style=``width: 75px;''>Value</th>
\item {} <th>Source \& Derivation</th>
\item {} </tr>
\item {} </thead>
\item {} <tbody>
\item {} <!-- Group: Quarterly Financials -->
\item {} <tr><td colspan=``4'' class=``appendix-group''>Quarterly Financials</td></tr>
\item {} <tr id=``ref-1''>
\item {} <td class=``ref-id''>1</td>
\item {} <td>[Q\# FY\#\#\#\# Revenue]</td>
\item {} <td class=``num''>\$[XX.X]B</td>
\item {} <td class=``source-detail''>
\item {} <span class=``src-label''>S\&P Capital IQ</span> — get\_\allowbreak{}financial\_\allowbreak{}line\_\allowbreak{}item\_\allowbreak{}from\_\allowbreak{}identifiers(identifier='[TICKER]', line\_\allowbreak{}item='revenue', period\_\allowbreak{}type='quarterly', period='[Q\# FY\#\#\#\#]')
\item {} </td>
\item {} </tr>
\item {} <tr id=``ref-2''>
\item {} <td class=``ref-id''>2</td>
\item {} <td>[Q\# FY\#\#\#\# Diluted EPS]</td>
\item {} <td class=``num''>\$[X.XX]</td>
\item {} <td class=``source-detail''>
\item {} <span class=``src-label''>S\&P Capital IQ</span> — get\_\allowbreak{}financial\_\allowbreak{}line\_\allowbreak{}item\_\allowbreak{}from\_\allowbreak{}identifiers(identifier='[TICKER]', line\_\allowbreak{}item='diluted\_\allowbreak{}eps', period\_\allowbreak{}type='quarterly', period='[Q\# FY\#\#\#\#]')
\item {} </td>
\item {} </tr>
\item {} <tr id=``ref-3''>
\item {} <td class=``ref-id''>3</td>
\item {} <td>[Q\# FY\#\#\#\# Gross Profit]</td>
\item {} <td class=``num''>\$[XX.X]B</td>
\item {} <td class=``source-detail''>
\item {} <span class=``src-label''>S\&P Capital IQ</span> — get\_\allowbreak{}financial\_\allowbreak{}line\_\allowbreak{}item\_\allowbreak{}from\_\allowbreak{}identifiers(identifier='[TICKER]', line\_\allowbreak{}item='gross\_\allowbreak{}profit', period\_\allowbreak{}type='quarterly', period='[Q\# FY\#\#\#\#]')
\item {} </td>
\item {} </tr>
\item {} <tr id=``ref-4''>
\item {} <td class=``ref-id''>4</td>
\item {} <td>[Q\# FY\#\#\#\# Gross Margin]</td>
\item {} <td class=``num''>[XX.X\%]</td>
\item {} <td class=``source-detail''>
\item {} <span class=``formula''><a href=``\#ref-3'' class=``data-ref''>Gross Profit \$XX.XB</a> / <a href=``\#ref-1'' class=``data-ref''>Revenue \$XX.XB</a> = XX.X\%</span><br>
\item {} <span class=``src-label''>S\&P Capital IQ</span> (calculated)
\item {} </td>
\item {} </tr>
\item {} <tr id=``ref-5''>
\item {} <td class=``ref-id''>5</td>
\item {} <td>[Q\# FY\#\#\#\# Revenue y/y Growth]</td>
\item {} <td class=``num''>[+/-X.X\%]</td>
\item {} <td class=``source-detail''>
\item {} <span class=``formula''>(<a href=``\#ref-1'' class=``data-ref''>[Q\# FY\#\# Rev \$XX.XB]</a> - <a href=``\#ref-N'' class=``data-ref''>[Q\# FY\#\# Rev \$XX.XB]</a>) / <a href=``\#ref-N'' class=``data-ref''>[Q\# FY\#\# Rev \$XX.XB]</a> = X.X\%</span><br>
\item {} <span class=``src-label''>S\&P Capital IQ</span> (calculated)
\item {} </td>
\item {} </tr>
\item {} <!-- Continue for all financial data points... -->
\item {} <!-- Group: Valuation -->
\item {} <tr><td colspan=``4'' class=``appendix-group''>Valuation</td></tr>
\item {} <tr id=``ref-N''>
\item {} <td class=``ref-id''>[N]</td>
\item {} <td>Current Stock Price — [TICKER]</td>
\item {} <td class=``num''>\$[XXX.XX]</td>
\item {} <td class=``source-detail''>
\item {} <span class=``src-label''>S\&P Capital IQ</span> — get\_\allowbreak{}prices\_\allowbreak{}from\_\allowbreak{}identifiers(identifier='[TICKER]', periodicity='day')
\item {} </td>
\item {} </tr>
\item {} <tr id=``ref-N''>
\item {} <td class=``ref-id''>[N]</td>
\item {} <td>Market Cap — [TICKER]</td>
\item {} <td class=``num''>\$[XXX.X]B</td>
\item {} <td class=``source-detail''>
\item {} <span class=``src-label''>S\&P Capital IQ</span> — get\_\allowbreak{}capitalization\_\allowbreak{}from\_\allowbreak{}identifiers(identifier='[TICKER]', capitalization='market\_\allowbreak{}cap')
\item {} </td>
\item {} </tr>
\item {} <tr id=``ref-N''>
\item {} <td class=``ref-id''>[N]</td>
\item {} <td>LTM P/E — [TICKER]</td>
\item {} <td class=``num''>[XX.X]x</td>
\item {} <td class=``source-detail''>
\item {} <span class=``formula''><a href=``\#ref-20'' class=``data-ref''>Price \$XXX.XX</a> / (<a href=``\#ref-8'' class=``data-ref''>Q1 EPS \$X.XX</a> + <a href=``\#ref-9'' class=``data-ref''>Q2 EPS \$X.XX</a> + <a href=``\#ref-10'' class=``data-ref''>Q3 EPS \$X.XX</a> + <a href=``\#ref-11'' class=``data-ref''>Q4 EPS \$X.XX</a>) = XX.Xx</span><br>
\item {} <span class=``src-label''>S\&P Capital IQ</span> (calculated)
\item {} </td>
\item {} </tr>
\item {} <tr id=``ref-N''>
\item {} <td class=``ref-id''>[N]</td>
\item {} <td>NTM P/E — [TICKER]</td>
\item {} <td class=``num''>[XX.X]x</td>
\item {} <td class=``source-detail''>
\item {} <span class=``formula''><a href=``\#ref-20'' class=``data-ref''>Price \$XXX.XX</a> / (<a href=``\#ref-N'' class=``data-ref''>Q4'25E \$X.XX</a> + <a href=``\#ref-N'' class=``data-ref''>Q1'26E \$X.XX</a> + <a href=``\#ref-N'' class=``data-ref''>Q2'26E \$X.XX</a> + <a href=``\#ref-N'' class=``data-ref''>Q3'26E \$X.XX</a>) = XX.Xx</span><br>
\item {} <span class=``src-label''>S\&P Capital IQ</span> — get\_\allowbreak{}consensus\_\allowbreak{}estimates\_\allowbreak{}from\_\allowbreak{}identifiers(identifier='[TICKER]', period\_\allowbreak{}type='quarterly', num\_\allowbreak{}periods\_\allowbreak{}forward=4). NTM EPS = sum of next 4 quarterly consensus mean EPS estimates.
\item {} </td>
\item {} </tr>
\item {} <!-- Group: Transcript Claims -->
\item {} <tr><td colspan=``4'' class=``appendix-group''>Transcript Claims</td></tr>
\item {} <tr id=``ref-N''>
\item {} <td class=``ref-id''>[N]</td>
\item {} <td>[Fact, e.g., ``Management guided comp sales +3-4\%'']</td>
\item {} <td class=``num''>N/A</td>
\item {} <td class=``source-detail''>
\item {} <span class=``excerpt''>``We expect comp sales growth of 3-4\% in Q4, driven by continued strength in grocery and health \&amp; wellness.''</span><br>
\item {} <span class=``src-label''>Source:</span> <span class=``transcript-ref''>[Q\# FY\#\#\#\# Earnings Call Transcript]</span> (key\_\allowbreak{}dev\_\allowbreak{}id: [ID]) — [Speaker Name], [Title]
\item {} </td>
\item {} </tr>
\item {} <!-- Group: Estimates \& Consensus -->
\item {} <tr><td colspan=``4'' class=``appendix-group''>Estimates \& Consensus</td></tr>
\item {} <tr id=``ref-N''>
\item {} <td class=``ref-id''>[N]</td>
\item {} <td>Consensus EPS — [Q\# FY\#\#\#\#]</td>
\item {} <td class=``num''>\$[X.XX]</td>
\item {} <td class=``source-detail''>
\item {} <span class=``excerpt''>``Consensus EPS estimate of \$X.XX, revised up from \$X.XX over the past 90 days.''</span><br>
\item {} <a href=``https://[source-url-from-kensho-search]'' target=``\_\allowbreak{}blank'' class=``src-url''>[Source Title / Publication Name]</a><br>
\item {} <span class=``src-label''>Query:</span> search(``[TICKER] earnings estimates consensus EPS revenue upcoming quarter'')
\item {} </td>
\item {} </tr>
\item {} <!-- Group: News \& Analyst Commentary -->
\item {} <tr><td colspan=``4'' class=``appendix-group''>News \& Analyst Commentary</td></tr>
\item {} <tr id=``ref-N''>
\item {} <td class=``ref-id''>[N]</td>
\item {} <td>[e.g., ``Barclays upgraded to Overweight'']</td>
\item {} <td class=``num''>N/A</td>
\item {} <td class=``source-detail''>
\item {} <span class=``excerpt''>``Barclays upgraded WMT to Overweight with a \$210 price target, citing accelerating eCommerce momentum.''</span><br>
\item {} <a href=``https://[source-url-from-kensho-search]'' target=``\_\allowbreak{}blank'' class=``src-url''>[Source Title / Publication, Date]</a><br>
\item {} <span class=``src-label''>Query:</span> search(``[TICKER] analyst ratings price target upgrades downgrades'')
\item {} </td>
\item {} </tr>
\item {} <!-- Group: Stock Performance -->
\item {} <tr><td colspan=``4'' class=``appendix-group''>Stock Performance</td></tr>
\item {} <tr id=``ref-N''>
\item {} <td class=``ref-id''>[N]</td>
\item {} <td>YTD Return — [TICKER]</td>
\item {} <td class=``num''>[+/-X.X\%]</td>
\item {} <td class=``source-detail''>
\item {} <span class=``formula''>(<a href=``\#ref-N'' class=``data-ref''>Current \$XXX.XX</a> - <a href=``\#ref-N'' class=``data-ref''>Dec 31 Close \$XXX.XX</a>) / <a href=``\#ref-N'' class=``data-ref''>Dec 31 Close \$XXX.XX</a> = X.X\%</span><br>
\item {} <span class=``src-label''>S\&P Capital IQ</span> (calculated from daily prices)
\item {} </td>
\item {} </tr>
\item {} </tbody>
\item {} </table>
\item {} </div>
\item {} <!-- -->
\item {} <!-- FOOTER -->
\item {} <!-- -->
\item {} <div class=``page-footer''>
\item {} <div class=``footer-disclaimer''>Analysis is AI-generated — please confirm all outputs</div>
\item {} <div class=``footer-meta''>Data: S\&P Capital IQ, Kensho | [Month Day, Year]</div>
\item {} </div>
\item {} </div>
\item {} <!-- -->
\item {} <!-- CHART.JS SCRIPTS -->
\item {} <!-- -->
\item {} <script>
\item {} // Register Annotation Plugin
\item {} // The CDN-loaded annotation plugin must be explicitly registered.
\item {} // It is available as a global after the script tag loads.
\item {} if (window['chartjs-plugin-annotation']) \{
\item {} Chart.register(window['chartjs-plugin-annotation']);
\item {} \}
\item {} // Chart Defaults
\item {} Chart.defaults.font.family = ``'Arial Narrow', Arial, sans-serif'';
\item {} Chart.defaults.font.size = 11;
\item {} Chart.defaults.color = '\#555';
\item {} Chart.defaults.plugins.legend.position = 'bottom';
\item {} Chart.defaults.plugins.legend.labels.boxWidth = 12;
\item {} // Color Palette
\item {} const COLORS = \{
\item {} navy: '\#1a1a4e',
\item {} blue: '\#3366cc',
\item {} teal: '\#0d9488',
\item {} orange: '\#e67e22',
\item {} red: '\#c0392b',
\item {} green: '\#27ae60',
\item {} purple: '\#8e44ad',
\item {} gray: '\#7f8c8d',
\item {} lightBlue:'\#85c1e9',
\item {} gold: '\#f0b429',
\item {} \};
\item {} const COMP\_\allowbreak{}COLORS = [
\item {} COLORS.navy, COLORS.blue, COLORS.teal,
\item {} COLORS.orange, COLORS.red, COLORS.green,
\item {} COLORS.purple, COLORS.gold
\item {} ];
\item {} // Helper: Revenue \& EPS Combo Chart
\item {} function createRevEpsChart(canvasId, labels, revenueData, epsData, revLabel) \{
\item {} const ctx = document.getElementById(canvasId);
\item {} if (!ctx) return;
\item {} new Chart(ctx, \{
\item {} type: 'bar',
\item {} data: \{
\item {} labels: labels,
\item {} datasets: [
\item {} \{
\item {} label: revLabel || 'Revenue (\$B)',
\item {} data: revenueData,
\item {} backgroundColor: COLORS.navy + 'cc',
\item {} borderColor: COLORS.navy,
\item {} borderWidth: 1,
\item {} yAxisID: 'y',
\item {} order: 2
\item {} \},
\item {} \{
\item {} label: 'Diluted EPS',
\item {} data: epsData,
\item {} type: 'line',
\item {} borderColor: COLORS.orange,
\item {} backgroundColor: COLORS.orange,
\item {} borderWidth: 2.5,
\item {} pointRadius: 4,
\item {} pointBackgroundColor: COLORS.orange,
\item {} tension: 0.3,
\item {} yAxisID: 'y1',
\item {} order: 1
\item {} \}
\item {} ]
\item {} \},
\item {} options: \{
\item {} responsive: true,
\item {} interaction: \{ mode: 'index', intersect: false \},
\item {} scales: \{
\item {} y: \{
\item {} position: 'left',
\item {} title: \{ display: true, text: revLabel || 'Revenue (\$B)', font: \{ size: 11 \} \},
\item {} grid: \{ color: '\#eee' \}
\item {} \},
\item {} y1: \{
\item {} position: 'right',
\item {} title: \{ display: true, text: 'EPS (\$)', font: \{ size: 11 \} \},
\item {} grid: \{ drawOnChartArea: false \}
\item {} \}
\item {} \}
\item {} \}
\item {} \});
\item {} \}
\item {} // Helper: Margin Trend Chart
\item {} function createMarginChart(canvasId, labels, grossMargins, opMargins) \{
\item {} const ctx = document.getElementById(canvasId);
\item {} if (!ctx) return;
\item {} new Chart(ctx, \{
\item {} type: 'line',
\item {} data: \{
\item {} labels: labels,
\item {} datasets: [
\item {} \{
\item {} label: 'Gross Margin \%',
\item {} data: grossMargins,
\item {} borderColor: COLORS.blue,
\item {} backgroundColor: COLORS.blue + '20',
\item {} borderWidth: 2.5,
\item {} pointRadius: 4,
\item {} fill: false,
\item {} tension: 0.3
\item {} \},
\item {} \{
\item {} label: 'Operating Margin \%',
\item {} data: opMargins,
\item {} borderColor: COLORS.teal,
\item {} backgroundColor: COLORS.teal + '20',
\item {} borderWidth: 2.5,
\item {} pointRadius: 4,
\item {} fill: false,
\item {} tension: 0.3
\item {} \}
\item {} ]
\item {} \},
\item {} options: \{
\item {} responsive: true,
\item {} scales: \{
\item {} y: \{
\item {} title: \{ display: true, text: 'Margin (\%)', font: \{ size: 11 \} \},
\item {} grid: \{ color: '\#eee' \},
\item {} ticks: \{ callback: v => v.toFixed(1) + '\%' \}
\item {} \}
\item {} \}
\item {} \}
\item {} \});
\item {} \}
\item {} // Helper: Revenue Growth Bar Chart
\item {} function createRevGrowthChart(canvasId, labels, growthData) \{
\item {} const ctx = document.getElementById(canvasId);
\item {} if (!ctx) return;
\item {} new Chart(ctx, \{
\item {} type: 'bar',
\item {} data: \{
\item {} labels: labels,
\item {} datasets: [\{
\item {} label: 'Revenue Growth y/y \%',
\item {} data: growthData,
\item {} backgroundColor: growthData.map(v => v >= 0 ? COLORS.green + 'cc' : COLORS.red + 'cc'),
\item {} borderColor: growthData.map(v => v >= 0 ? COLORS.green : COLORS.red),
\item {} borderWidth: 1
\item {} \}]
\item {} \},
\item {} options: \{
\item {} responsive: true,
\item {} scales: \{
\item {} y: \{
\item {} title: \{ display: true, text: 'Growth (\%)', font: \{ size: 11 \} \},
\item {} grid: \{ color: '\#eee' \},
\item {} ticks: \{ callback: v => v.toFixed(1) + '\%' \}
\item {} \}
\item {} \},
\item {} plugins: \{ legend: \{ display: false \} \}
\item {} \}
\item {} \});
\item {} \}
\item {} // Helper: Earnings-Annotated Stock Price Chart
\item {} function createAnnotatedPriceChart(canvasId, labels, prices, earningsDates, ticker) \{
\item {} const ctx = document.getElementById(canvasId);
\item {} if (!ctx) return;
\item {} const annotations = \{\};
\item {} earningsDates.forEach((ed, i) => \{
\item {} let xValue = ed.date;
\item {} const isNeg = ed.move.startsWith('-');
\item {} annotations['earnings' + i] = \{
\item {} type: 'line',
\item {} xMin: xValue,
\item {} xMax: xValue,
\item {} borderColor: isNeg ? '\#c0392b' : '\#0d7a3e',
\item {} borderWidth: 2,
\item {} borderDash: [6, 4],
\item {} label: \{
\item {} display: true,
\item {} content: ed.label + ' (' + ed.move + ')',
\item {} position: i \% 2 === 0 ? 'start' : 'end',
\item {} backgroundColor: isNeg ? '\#c0392b' : '\#0d7a3e',
\item {} color: '\#fff',
\item {} font: \{ size: 10, weight: 'bold' \},
\item {} padding: \{ top: 3, bottom: 3, left: 6, right: 6 \},
\item {} borderRadius: 3
\item {} \}
\item {} \};
\item {} \});
\item {} new Chart(ctx, \{
\item {} type: 'line',
\item {} data: \{
\item {} labels: labels,
\item {} datasets: [\{
\item {} label: ticker + ' Close',
\item {} data: prices,
\item {} borderColor: COLORS.navy,
\item {} backgroundColor: COLORS.navy + '15',
\item {} borderWidth: 1.5,
\item {} pointRadius: 0,
\item {} pointHitRadius: 4,
\item {} fill: true,
\item {} tension: 0.1
\item {} \}]
\item {} \},
\item {} options: \{
\item {} responsive: true,
\item {} interaction: \{ mode: 'index', intersect: false \},
\item {} scales: \{
\item {} x: \{ type: 'category', ticks: \{ maxTicksLimit: 12, font: \{ size: 10 \} \}, grid: \{ display: false \} \},
\item {} y: \{ title: \{ display: true, text: 'Price (\$)', font: \{ size: 11 \} \}, grid: \{ color: '\#eee' \} \}
\item {} \},
\item {} plugins: \{
\item {} annotation: \{ annotations: annotations \},
\item {} tooltip: \{ callbacks: \{ label: ctx => ticker + ': \$' + ctx.raw.toFixed(2) \} \}
\item {} \}
\item {} \}
\item {} \});
\item {} \}
\item {} // Helper: Competitor Indexed Performance Chart
\item {} // datasets: [\{ label: 'TICKER', data: [price1, price2, ...], color: '\#xxx', isSubject: true/false \}, ...]
\item {} function createCompPerfChart(canvasId, labels, datasets) \{
\item {} const ctx = document.getElementById(canvasId);
\item {} if (!ctx) return;
\item {} const chartDatasets = datasets.map((ds, i) => \{
\item {} const base = ds.data[0] || 1;
\item {} return \{
\item {} label: ds.label,
\item {} data: ds.data.map(v => (v / base) * 100),
\item {} borderColor: ds.color || COMP\_\allowbreak{}COLORS[i \% COMP\_\allowbreak{}COLORS.length],
\item {} backgroundColor: 'transparent',
\item {} borderWidth: ds.isSubject ? 3 : 1.5,
\item {} borderDash: ds.isSubject ? [] : [4, 2],
\item {} pointRadius: 0,
\item {} tension: 0.2
\item {} \};
\item {} \});
\item {} new Chart(ctx, \{
\item {} type: 'line',
\item {} data: \{ labels: labels, datasets: chartDatasets \},
\item {} options: \{
\item {} responsive: true,
\item {} interaction: \{ mode: 'index', intersect: false \},
\item {} scales: \{
\item {} y: \{ title: \{ display: true, text: 'Indexed (100 = Start)', font: \{ size: 11 \} \}, grid: \{ color: '\#eee' \} \},
\item {} x: \{ ticks: \{ maxTicksLimit: 12 \} \}
\item {} \},
\item {} plugins: \{
\item {} tooltip: \{ callbacks: \{ label: ctx => ctx.dataset.label + ': ' + ctx.raw.toFixed(1) \} \}
\item {} \}
\item {} \}
\item {} \});
\item {} \}
\item {} // Helper: LTM P/E Horizontal Bar Chart
\item {} // companies: [\{ label: 'TICKER', pe: 25.3, isSubject: true/false \}, ...]
\item {} function createPEChart(canvasId, companies) \{
\item {} const ctx = document.getElementById(canvasId);
\item {} if (!ctx) return;
\item {} companies.sort((a, b) => b.pe - a.pe);
\item {} new Chart(ctx, \{
\item {} type: 'bar',
\item {} data: \{
\item {} labels: companies.map(c => c.label),
\item {} datasets: [\{
\item {} label: 'LTM P/E',
\item {} data: companies.map(c => c.pe),
\item {} backgroundColor: companies.map(c => c.isSubject ? COLORS.navy : COLORS.lightBlue),
\item {} borderColor: companies.map(c => c.isSubject ? COLORS.navy : COLORS.blue),
\item {} borderWidth: 1
\item {} \}]
\item {} \},
\item {} options: \{
\item {} indexAxis: 'y',
\item {} responsive: true,
\item {} scales: \{
\item {} x: \{
\item {} title: \{ display: true, text: 'LTM P/E', font: \{ size: 11 \} \},
\item {} grid: \{ color: '\#eee' \}
\item {} \},
\item {} y: \{
\item {} ticks: \{ font: \{ size: 12, weight: 'bold' \} \}
\item {} \}
\item {} \},
\item {} plugins: \{
\item {} legend: \{ display: false \},
\item {} tooltip: \{
\item {} callbacks: \{ label: ctx => 'P/E: ' + ctx.raw.toFixed(1) + 'x' \}
\item {} \}
\item {} \}
\item {} \}
\item {} \});
\item {} \}
\item {} //
\item {} // HELPER FUNCTIONS DEFINED ABOVE — DO NOT REWRITE THEM.
\item {} // Use ONLY these functions to create charts.
\item {} // DO NOT write custom inline Chart.js code.
\item {} //
\item {} </script>
\item {} <!-- -->
\item {} <!-- CHART DATA — EACH CHART IN ITS OWN SCRIPT + TRY-CATCH -->
\item {} <!-- A syntax error in one chart must NOT break the others. -->
\item {} <!-- MANDATORY: Use the helper functions above. No custom code. -->
\item {} <!-- -->
\item {} <!-- Figure 1: Revenue \& EPS -->
\item {} <script>
\item {} try \{
\item {} createRevEpsChart('chart-rev-eps',
\item {} ['Q1 FY24','Q2 FY24','Q3 FY24','Q4 FY24','Q1 FY25','Q2 FY25','Q3 FY25','Q4 FY25'],
\item {} [152.3, 161.6, 160.8, 173.4, 161.5, 169.3, 165.8, 178.0], // revenue in \$B
\item {} [1.47, 1.84, 1.53, 1.80, 1.56, 1.92, 1.60, 1.90], // diluted EPS
\item {} 'Revenue (\$B)'
\item {} );
\item {} \} catch(e) \{ console.error('Figure 1 error:', e); \}
\item {} </script>
\item {} <!-- Figure 2: Margin Trends -->
\item {} <script>
\item {} try \{
\item {} createMarginChart('chart-margins',
\item {} ['Q1 FY24','Q2 FY24','Q3 FY24','Q4 FY24','Q1 FY25','Q2 FY25','Q3 FY25','Q4 FY25'],
\item {} [24.0, 24.4, 24.2, 23.8, 24.5, 24.8, 24.6, 24.1], // gross margin \%
\item {} [4.2, 5.1, 4.5, 4.8, 4.6, 5.3, 4.7, 5.0] // operating margin \%
\item {} );
\item {} \} catch(e) \{ console.error('Figure 2 error:', e); \}
\item {} </script>
\item {} <!-- Figure 3: Revenue Growth y/y — ONLY quarters where y/y can be computed (most recent 4) -->
\item {} <script>
\item {} try \{
\item {} createRevGrowthChart('chart-rev-growth',
\item {} ['Q1 FY25','Q2 FY25','Q3 FY25','Q4 FY25'], // Only 4 labels — quarters with y/y data
\item {} [6.0, 4.8, 3.1, 2.7] // y/y revenue growth \% for those 4 quarters
\item {} );
\item {} \} catch(e) \{ console.error('Figure 3 error:', e); \}
\item {} </script>
\item {} <!-- Figure 5: Annotated Stock Price -->
\item {} <script>
\item {} try \{
\item {} createAnnotatedPriceChart('chart-price-annotated',
\item {} ['2025-02-18','2025-02-19'], // ... daily date labels for 1 year
\item {} [170.5, 171.2], // ... daily closing prices
\item {} [
\item {} \{ date: '2025-05-15', label: 'Q1 FY26', move: '+3.2\%' \},
\item {} \{ date: '2025-08-15', label: 'Q2 FY26', move: '-1.8\%' \}
\item {} ],
\item {} 'WMT'
\item {} );
\item {} \} catch(e) \{ console.error('Figure 5 error:', e); \}
\item {} </script>
\item {} <!-- Figure 6: Competitor Indexed Performance -->
\item {} <script>
\item {} try \{
\item {} createCompPerfChart('chart-comp-perf',
\item {} ['2025-02-18','2025-03-18'], // ... date labels
\item {} [
\item {} \{ label: 'WMT', data: [170.5, 172.3], isSubject: true \},
\item {} \{ label: 'COST', data: [580.2, 595.1], isSubject: false \},
\item {} \{ label: 'TGT', data: [142.0, 138.5], isSubject: false \}
\item {} ]
\item {} );
\item {} \} catch(e) \{ console.error('Figure 6 error:', e); \}
\item {} </script>
\item {} <!-- Figure 7: LTM P/E Comparison -->
\item {} <script>
\item {} try \{
\item {} createPEChart('chart-pe-comp', [
\item {} \{ label: 'COST', pe: 52.3, isSubject: false \},
\item {} \{ label: 'WMT', pe: 28.1, isSubject: true \},
\item {} \{ label: 'TGT', pe: 15.6, isSubject: false \},
\item {} \{ label: 'BJ', pe: 22.4, isSubject: false \}
\item {} ]);
\item {} \} catch(e) \{ console.error('Figure 7 error:', e); \}
\item {} </script>
\item {} </body>
\item {} </html>
\end{itemize}

\subparagraph{Chart.js Implementation Notes}\mbox{}\par

\subparagraph{Figure 2: Revenue \& EPS Chart}\mbox{}\par
\begin{itemize}[leftmargin=*,itemsep=2pt,topsep=2pt]
\item {} \textbf{Type}: Combo bar + line
\item {} \textbf{Bars}: Quarterly revenue on left y-axis
\item {} \textbf{Line}: Diluted EPS on right y-axis
\item {} \textbf{Labels}: Quarter identifiers (e.g., ``Q1 FY24'')
\item {} Use 8 quarters of data
\end{itemize}

\subparagraph{Figure 3: Margin Trend Chart}\mbox{}\par
\begin{itemize}[leftmargin=*,itemsep=2pt,topsep=2pt]
\item {} \textbf{Type}: Dual line chart
\item {} \textbf{Lines}: Gross margin \% and operating margin \%
\item {} \textbf{Y-axis}: Percentage with 1 decimal place
\end{itemize}

\subparagraph{Figure 3: Revenue Growth Chart}\mbox{}\par
\begin{itemize}[leftmargin=*,itemsep=2pt,topsep=2pt]
\item {} \textbf{Type}: Bar chart with conditional coloring
\item {} \textbf{Green bars}: Positive growth quarters
\item {} \textbf{Red bars}: Negative growth quarters
\item {} \textbf{IMPORTANT}: Only include quarters where y/y growth can be computed (i.e., where both the current quarter AND the year-ago quarter exist in \emph{financials.csv}). With 8 quarters of raw data, this typically yields 4 bars — NOT 8. Do not pass labels for quarters without y/y data.
\item {} No legend needed (self-explanatory)
\end{itemize}

\subparagraph{Figure 4: Business Segment Revenue}\mbox{}\par
\begin{itemize}[leftmargin=*,itemsep=2pt,topsep=2pt]
\item {} Use HTML table (not a chart)
\item {} Columns: Segment | Latest Q Rev (\$M) | \% of Total | y/y Change
\item {} Color-code y/y change cells with pos/neg classes
\end{itemize}

\subparagraph{Figure 5: Earnings-Annotated Stock Price Chart}\mbox{}\par
\begin{itemize}[leftmargin=*,itemsep=2pt,topsep=2pt]
\item {} \textbf{Type}: Line chart with annotation plugin vertical lines
\item {} \textbf{Data}: 1 year of daily closing prices
\item {} \textbf{Annotations}: Vertical dashed lines at each earnings date
\item {} \textbf{Labels}: Quarter name + 1-day post-earnings stock move
\item {} \textbf{Colors}: Green for positive reactions, red for negative
\item {} \textbf{Calculating the 1-day move}: Compare closing price on earnings date to next trading day close
\item {} \textbf{CRITICAL}: The annotation plugin MUST be registered before creating charts: \emph{Chart.register(window['chartjs-plugin-annotation'])} — this is already in the template script block
\end{itemize}

\subparagraph{Figure 7: Competitor Indexed Performance Chart}\mbox{}\par
\begin{itemize}[leftmargin=*,itemsep=2pt,topsep=2pt]
\item {} \textbf{Type}: Multi-line chart, rebased to 100
\item {} \textbf{Subject company}: Solid thick line (borderWidth: 3)
\item {} \textbf{Competitors}: Thinner dashed lines (borderWidth: 1.5, borderDash)
\item {} This visual hierarchy makes the subject company immediately identifiable
\end{itemize}

\subparagraph{Figure 8: LTM P/E Comparison Chart}\mbox{}\par
\begin{itemize}[leftmargin=*,itemsep=2pt,topsep=2pt]
\item {} \textbf{Type}: Horizontal bar chart
\item {} \textbf{Subject company}: Highlighted in navy (\#1a1a4e)
\item {} \textbf{Competitors}: Light blue (\#85c1e9)
\item {} \textbf{Sorted}: Descending by P/E
\item {} Shows at a glance whether the company trades at a premium or discount to peers
\end{itemize}

\subparagraph{Figure 8: Competitor Comparison Table}\mbox{}\par
\begin{itemize}[leftmargin=*,itemsep=2pt,topsep=2pt]
\item {} HTML table with highlight-row for subject company
\item {} Columns: Ticker | Company | Mkt Cap (\$B) | LTM P/E | NTM P/E | YTD \% | 1-Yr \%
\item {} Color-code returns with pos/neg classes
\end{itemize}

\subparagraph{Formatting Conventions}\mbox{}\par

\subparagraph{Numbers}\mbox{}\par
\begin{itemize}[leftmargin=*,itemsep=2pt,topsep=2pt]
\item {} Revenue: 1 decimal place for \$B (e.g., ``\$152.3B''), no decimals for \$M (e.g., ``\$4,521M'')
\item {} EPS: 2 decimal places (e.g., ``\$1.47'')
\item {} Margins: 1 decimal place with \% sign (e.g., ``24.5\%'')
\item {} Growth rates: 1 decimal place with +/- sign (e.g., ``+5.2\%'', ``-3.1\%'')
\item {} Market cap: 1 decimal place for \$B (e.g., ``\$562.1B'')
\item {} Stock prices: 2 decimal places (e.g., ``\$172.35'')
\item {} P/E ratios: 1 decimal place with 'x' suffix (e.g., ``25.3x'')
\end{itemize}

\subparagraph{Color Coding}\mbox{}\par
\begin{itemize}[leftmargin=*,itemsep=2pt,topsep=2pt]
\item {} Positive values: \emph{class=``pos''} -- green (\#0d7a3e)
\item {} Negative values: \emph{class=``neg''} -- red (\#c0392b)
\item {} Neutral/flat: \emph{class=``neutral''} -- gray (\#555)
\item {} Subject company row: \emph{class=``highlight-row''} -- light blue background
\end{itemize}

\subparagraph{Figure Labels}\mbox{}\par
\begin{itemize}[leftmargin=*,itemsep=2pt,topsep=2pt]
\item {} Number all figures sequentially: ``Figure 1:'', ``Figure 2:'', etc.
\item {} Figures 1-8 are on Pages 3-5 (the Consensus Estimates table on Page 2 is not numbered)
\item {} Include source attribution under every chart and table: ``Source: S\&P Capital IQ''
\end{itemize}

\subparagraph{Hyperlinked Claims}\mbox{}\par
\begin{itemize}[leftmargin=*,itemsep=2pt,topsep=2pt]
\item {} Every factual claim in the report body — numbers AND qualitative statements — must be wrapped in \emph{<a href=``\#ref-N'' class=``data-ref''>CLAIM TEXT</a>}
\item {} The \emph{ref-N} ID must match a row in the Appendix table
\item {} This applies to: narrative text, bullet points, table cells, blockquotes — anywhere a fact appears
\item {} Chart axis labels and tooltips do NOT need hyperlinks (only report body text)
\item {} Assign reference IDs sequentially (\emph{ref-1}, \emph{ref-2}, ...) as you write the report
\item {} Multiple references to the same underlying claim should share the same ref ID
\item {} For qualitative claims, wrap the key phrase: \emph{<a href=``\#ref-25'' class=``data-ref''>management flagged tariff headwinds</a>}
\end{itemize}

\subparagraph{Appendix}\mbox{}\par
\begin{itemize}[leftmargin=*,itemsep=2pt,topsep=2pt]
\item {} \textbf{MUST begin with}: \emph{<div class=``ai-disclaimer''>Analysis is AI-generated — please confirm all outputs</div>}
\item {} The Appendix is the final section of the report, after all figures
\item {} \textbf{4 columns}: Ref \# | Fact | Value | Source \& Derivation
\item {} One row per unique claim referenced in the report (numeric and non-numeric)
\item {} \textbf{Every number in the report body must be a clickable @@TOKEN0@@ link to its appendix row. No exceptions.}
\item {} Group rows by category: Quarterly Financials, Valuation, Transcript Claims, Estimates \& Consensus, News \& Analyst Commentary, Stock Performance
\item {} Use subheading rows (\emph{appendix-group} class) to separate groups
\item {} \textbf{Source \& Derivation column} must include specific, detailed sourcing for EVERY row:
\begin{itemize}[leftmargin=*,itemsep=2pt,topsep=2pt]
\item {} For raw S\&P data (revenue, EPS, prices, market cap, etc.): \emph{<span class=``src-label''>S\&P Capital IQ</span>} followed by the specific MCP function call with parameters (e.g., \emph{get\_\allowbreak{}financial\_\allowbreak{}line\_\allowbreak{}item\_\allowbreak{}from\_\allowbreak{}identifiers(identifier='WMT', line\_\allowbreak{}item='revenue', period\_\allowbreak{}type='quarterly', period='Q3 FY2026')}). \textbf{Never write just ``S\&P Capital IQ'' with no detail.}
\item {} For calculated values (margins, growth rates, P/E, returns): the full formula with \emph{<a class=``data-ref''>} hyperlinks to each component row (use \emph{formula} CSS class). \textbf{Every number in the formula must be a clickable link to its own appendix row.}
\item {} For transcript claims: the verbatim excerpt sentence in italics (\emph{excerpt} CSS class) + transcript name with \emph{transcript-ref} class + \emph{key\_\allowbreak{}dev\_\allowbreak{}id}
\item {} For Kensho results: the key finding (\emph{excerpt} class) + \textbf{clickable source URL} as \emph{<a href=``[URL]'' target=``\_\allowbreak{}blank'' class=``src-url''>[Source Title]</a>} + the search query used. \textbf{Every Kensho-sourced claim must have a clickable URL to the original source.}
\end{itemize}
\item {} Source labels use the \emph{src-label} CSS class (bold navy)
\item {} External source URLs use the \emph{src-url} CSS class (blue, underlined, clickable)
\end{itemize}

\subparagraph{Style Rules}\mbox{}\par
\begin{itemize}[leftmargin=*,itemsep=2pt,topsep=2pt]
\item {} \textbf{NO EMOJIS} anywhere in the report. No emoji in headings, tables, chart labels, or body text. This is a professional research document.
\item {} Font: Arial Narrow throughout (body, headings, tables, charts).
\item {} Management quotes: Integrate as \emph{<blockquote>} elements within the executive thesis narrative. Never under a separate heading.
\item {} Keep all text concise. Target 4-5 printed pages total (appendix is additional).
\end{itemize}

\vspace{0.8em}\hrule\vspace{0.8em}

\subsection{Skill: funding-digest}\label{file:spglobal-skills-funding-digest-SKILL.md}

\paragraph{File Metadata}\mbox{}\par
\begingroup
\small
\setlength{\tabcolsep}{2pt}
\begin{longtable}{>{\RaggedRight\arraybackslash\hspace{0pt}}p{0.480\textwidth}>{\RaggedRight\arraybackslash\hspace{0pt}}p{0.480\textwidth}}
\toprule
{} \textbf{Field} & \textbf{Value} \\
\midrule
{} name & funding-digest \\
{} description & Generate a polished one-page PowerPoint slide summarizing key takeaways from recent funding rounds and notable capital markets activity across a user's watched sectors or companies. Use this skill when the user asks for a deal flow summary, weekly recap, funding digest, transaction roundup, or capital markets briefing. Triggers on: 'deal flow digest', 'weekly funding recap', 'deal roundup', 'transaction summary this week', 'what happened in [sector] this week', 'capital markets update', or any request to compile recent funding activity into a briefing slide. Produces a professional single-slide PPTX with key takeaways, valuation data, and Capital IQ deal links. \\
\bottomrule
\end{longtable}
\endgroup


\textbf{AI DISCLAIMER (MANDATORY):} You MUST include the following disclaimer text in the powerpoint footer. This is not optional — the report is incomplete without it:


\begin{quote}
``\textbf{``Analysis is AI-generated — please confirm all outputs''}''
\end{quote}

\textbf{Footer} — At the bottom of the generated slide, as a prominent yellow banner: ``Analysis is AI-generated — please confirm all outputs''


\vspace{0.5em}\hrule\vspace{0.5em}

\paragraph{Weekly Deal Flow Digest}\mbox{}\par

Generate an analyst-quality \textbf{single-slide PowerPoint} that summarizes key takeaways from recent funding rounds across watched sectors or companies, using S\&P Global Capital IQ data. Each deal links back to its Capital IQ profile for quick drill-down.


\subparagraph{When to Use}\mbox{}\par

Trigger on any of these patterns:

\begin{itemize}[leftmargin=*,itemsep=2pt,topsep=2pt]
\item {} ``Give me a deal flow digest for this week''
\item {} ``Weekly funding recap for [sector]''
\item {} ``What deals closed in [sector/companies] recently?''
\item {} ``Transaction roundup'' or ``deal roundup''
\item {} ``Capital markets update for my coverage universe''
\item {} ``Summarize recent funding activity''
\item {} Any periodic briefing request about deals, raises, or rounds
\end{itemize}

\subparagraph{Nested Skills}\mbox{}\par

This skill produces a one-slide PPTX briefing:

\begin{itemize}[leftmargin=*,itemsep=2pt,topsep=2pt]
\item {} \textbf{Read} \emph{/mnt/skills/public/pptx/SKILL.md} before generating the PowerPoint (and its sub-reference \emph{pptxgenjs.md} for creating from scratch)
\end{itemize}

\subparagraph{Entity Resolution \& Tool Robustness}\mbox{}\par

S\&P Global's identifier system resolves company names to legal entities. This works well for most companies but has known failure modes that cause empty results. \textbf{Apply these rules throughout the workflow to avoid silent data loss.}


\subparagraph{Rule 0: Pre-validate ALL identifiers before querying funding}\mbox{}\par

\textbf{Before} calling any funding tools, run every identifier through \emph{get\_\allowbreak{}info\_\allowbreak{}from\_\allowbreak{}identifiers}. This is the cheapest and most reliable way to catch problems early. Check two things in the response:


\begin{enumerate}[leftmargin=*,itemsep=2pt,topsep=2pt]
\item {} \textbf{Did it resolve at all?} If the identifier returns empty/error, the name doesn't exist in S\&P Global. Try the alias from \emph{references/sector-seeds.md}, the legal entity name, or the \emph{company\_\allowbreak{}id} directly.
\item {} \textbf{What is the @@TOKEN0@@ field?}
\begin{itemize}[leftmargin=*,itemsep=2pt,topsep=2pt]
\item {} \emph{``Operating''} → Safe to query for funding rounds.
\item {} \emph{``Operating Subsidiary''} → The company exists but is owned by a parent. It will return \textbf{zero funding rounds}. Note this in the digest as context (e.g., ``acquired by [Parent]'') but do not query for funding.
\item {} Any other status (e.g., closed, inactive) → The company is no longer operating. Historical data may exist but no new activity.
\end{itemize}
\end{enumerate}

\textbf{This single pre-validation step prevents the majority of empty-result issues.} Batch all candidates into a single \emph{get\_\allowbreak{}info\_\allowbreak{}from\_\allowbreak{}identifiers} call (it handles large batches well) and triage before proceeding.


\subparagraph{Rule 1: Never trust empty results without a fallback}\mbox{}\par

If \emph{get\_\allowbreak{}rounds\_\allowbreak{}of\_\allowbreak{}funding\_\allowbreak{}from\_\allowbreak{}identifiers} returns empty for a company you expect to have data:

\begin{enumerate}[leftmargin=*,itemsep=2pt,topsep=2pt]
\item {} \textbf{Try the legal entity name or company\_\allowbreak{}id.} Brand names usually work, but some don't. See the alias table in \emph{references/sector-seeds.md} for known mismatches. Common pattern: ``[Brand] AI'' → ``[Legal Name], Inc.'' (e.g., Together AI → ``Together Computer, Inc.'', Character.ai → ``Character Technologies, Inc.'', Runway ML → ``Runway AI, Inc.'').
\item {} \textbf{Verify the company exists in S\&P.} If you skipped Rule 0, call \emph{get\_\allowbreak{}info\_\allowbreak{}from\_\allowbreak{}identifiers(identifiers=[``Company''])} now — if this also returns empty, the company may be too early-stage or not yet indexed.
\end{enumerate}

\subparagraph{Rule 2: Subsidiaries have no funding rounds}\mbox{}\par

Companies that are divisions or wholly-owned subsidiaries of larger companies (e.g., DeepMind under Alphabet, GitHub under Microsoft, BeReal under Voodoo) will return \textbf{zero funding rounds}. Their capital events are tracked at the parent level.


\textbf{How to detect:} The \emph{status} field from \emph{get\_\allowbreak{}info\_\allowbreak{}from\_\allowbreak{}identifiers} will show \emph{``Operating Subsidiary''}. The \emph{references/sector-seeds.md} file also flags known subsidiaries with ⚠ warnings. Skip these for funding queries.


\subparagraph{Rule 3: Use \emph{get\_\allowbreak{}rounds\_\allowbreak{}of\_\allowbreak{}funding\_\allowbreak{}from\_\allowbreak{}identifiers} as the primary tool, not \emph{get\_\allowbreak{}funding\_\allowbreak{}summary\_\allowbreak{}from\_\allowbreak{}identifiers}}\mbox{}\par

The summary tool is faster but less reliable — it can return errors or incomplete data even when detailed rounds exist. Always use the detailed rounds tool as the primary data source. The summary tool is acceptable only for quick aggregate checks (total raised, round count) and should be verified against the rounds tool if results seem low.


\subparagraph{Rule 4: Batch carefully and validate}\mbox{}\par

When processing large company universes (50+ companies), batch in groups of 15–20. After each batch, check for companies that returned empty results and run them through the fallback steps in Rule 1 before moving on.


\subparagraph{Rule 5: The \emph{role} parameter is critical}\mbox{}\par

\begin{itemize}[leftmargin=*,itemsep=2pt,topsep=2pt]
\item {} \emph{company\_\allowbreak{}raising\_\allowbreak{}funds} → ``What rounds did X raise?'' (company perspective)
\item {} \emph{company\_\allowbreak{}investing\_\allowbreak{}in\_\allowbreak{}round\_\allowbreak{}of\_\allowbreak{}funding} → ``What did investor Y invest in?'' (investor perspective)
\end{itemize}

Using the wrong role returns empty results silently. For deal flow digests, you almost always want \emph{company\_\allowbreak{}raising\_\allowbreak{}funds}. Only use the investor role when specifically analyzing an investor's portfolio activity.


\subparagraph{Rule 6: Identifier resolution is case-insensitive but spelling-sensitive}\mbox{}\par

S\&P Global handles case variations (``openai'' = ``OpenAI'') but is strict on spelling and punctuation. ``Character AI'' may fail where ``Character.ai'' succeeds. When in doubt, use the \emph{company\_\allowbreak{}id} (e.g., \emph{C\_\allowbreak{}1829047235}) which is guaranteed to resolve.


\subparagraph{Workflow}\mbox{}\par

\subparagraph{Step 1: Establish Coverage \& Period}\mbox{}\par

Determine what the digest should cover. There are two setups:


\textbf{Returning user (has a watchlist):} If the user has previously defined sectors or companies to track, use that list. Check conversation history for prior watchlists.


\textbf{New user:} Ask for:


\begingroup
\small
\setlength{\tabcolsep}{2pt}
\begin{longtable}{>{\RaggedRight\arraybackslash\hspace{0pt}}p{0.320\textwidth}>{\RaggedRight\arraybackslash\hspace{0pt}}p{0.320\textwidth}>{\RaggedRight\arraybackslash\hspace{0pt}}p{0.320\textwidth}}
\toprule
{} \textbf{Parameter} & \textbf{Default} & \textbf{Notes} \\
\midrule
{} \textbf{Sectors} & \emph{(at least one)} & e.g., ``AI, Fintech, Biotech'' \\
{} \textbf{Specific companies} & Optional & Supplement sector-level coverage \\
{} \textbf{Time period} & Last 7 days & ``This week'', ``last 2 weeks'', ``this month'' \\
\bottomrule
\end{longtable}
\endgroup

Calculate the exact \emph{start\_\allowbreak{}date} and \emph{end\_\allowbreak{}date} from the time period.


\subparagraph{Step 2: Build the Company Universe}\mbox{}\par

For each sector specified, build a company universe using a validated bootstrapping approach:


\begin{enumerate}[leftmargin=*,itemsep=2pt,topsep=2pt]
\item {} \textbf{Seed companies} from domain knowledge (see \emph{references/sector-seeds.md})
\begin{itemize}[leftmargin=*,itemsep=2pt,topsep=2pt]
\item {} Pay attention to the ⚠ warnings and alias notes in the seeds file — some well-known companies are subsidiaries, have been acquired, or require a specific legal name to resolve.
\item {} The seeds file includes \emph{company\_\allowbreak{}id} values for known alias mismatches. Use these directly if the brand name fails.
\end{itemize}
\item {} \textbf{Pre-validate all seeds immediately} (Rule 0): ``\emph{ get\_\allowbreak{}info\_\allowbreak{}from\_\allowbreak{}identifiers(identifiers=[all\_\allowbreak{}seeds\_\allowbreak{}for\_\allowbreak{}this\_\allowbreak{}sector]) }`` Triage the results into two buckets:
\begin{itemize}[leftmargin=*,itemsep=2pt,topsep=2pt]
\item {} ✅ \textbf{Resolved \& Operating} (\emph{status} = ``Operating'') → proceed to competitor expansion
\item {} ❌ \textbf{Unresolved or Subsidiary} → retry with alias/legal name from seeds file; subsidiaries are noted for context but excluded from funding queries
\end{itemize}
\item {} \textbf{Expand via competitors} (using only the ✅ resolved seeds): ``\emph{ get\_\allowbreak{}competitors\_\allowbreak{}from\_\allowbreak{}identifiers(identifiers=[resolved\_\allowbreak{}seeds], competitor\_\allowbreak{}source=``all'') }``
\item {} \textbf{Validate expanded universe:} ``\emph{ get\_\allowbreak{}info\_\allowbreak{}from\_\allowbreak{}identifiers(identifiers=[new\_\allowbreak{}competitors]) }`\emph{ Apply the same triage. Filter by }simple\_\allowbreak{}industry` matching the target sector. Drop any unresolved names or subsidiaries.
\end{enumerate}

If the user provides specific companies, add those directly but still run them through the pre-validation triage. Never skip validation — even well-known brand names can fail silently.


Keep the universe manageable — aim for 15–40 \textbf{resolved, operating} companies per sector. For a multi-sector digest, this might total 50–100+ companies.


\subparagraph{Step 3: Pull Funding Rounds}\mbox{}\par

For all companies in the universe:


\textbf{Structured Template}
\begin{itemize}[leftmargin=*,itemsep=2pt,topsep=2pt]
\item {} get\_\allowbreak{}rounds\_\allowbreak{}of\_\allowbreak{}funding\_\allowbreak{}from\_\allowbreak{}identifiers(
\item {} identifiers=[batch],
\item {} role=``company\_\allowbreak{}raising\_\allowbreak{}funds'',
\item {} start\_\allowbreak{}date=``YYYY-MM-DD'',
\item {} end\_\allowbreak{}date=``YYYY-MM-DD''
\item {} )
\end{itemize}

Process in batches of 15–20 if the universe is large.


\textbf{After each batch, identify companies with empty results.} For any company expected to have activity:

\begin{enumerate}[leftmargin=*,itemsep=2pt,topsep=2pt]
\item {} Retry with the legal entity name or alternate identifier (see Entity Resolution rules above).
\item {} Log the company as ``no data'' only after exhausting fallbacks.
\end{enumerate}

Collect all \emph{transaction\_\allowbreak{}id} values from successful results, then enrich with detailed round info:


\textbf{Structured Template}
\begin{itemize}[leftmargin=*,itemsep=2pt,topsep=2pt]
\item {} get\_\allowbreak{}rounds\_\allowbreak{}of\_\allowbreak{}funding\_\allowbreak{}info\_\allowbreak{}from\_\allowbreak{}transaction\_\allowbreak{}ids(
\item {} transaction\_\allowbreak{}ids=[all\_\allowbreak{}funding\_\allowbreak{}ids]
\item {} )
\end{itemize}

Pass ALL transaction IDs in a single call (or small number of calls) rather than one per transaction — the tool handles batches efficiently.


\textbf{Extract the following from each round (critical for the slide):}

\begin{itemize}[leftmargin=*,itemsep=2pt,topsep=2pt]
\item {} \emph{transaction\_\allowbreak{}id} — needed for the Capital IQ deal link
\item {} \textbf{Announcement date} — when the round was publicly announced
\item {} \textbf{Close date} — when the round officially closed
\item {} Amount raised
\item {} \textbf{Pre-money valuation} (if disclosed)
\item {} \textbf{Post-money valuation} (if disclosed)
\item {} Lead investors
\item {} Round type (Series A, B, C, etc.)
\item {} Security terms
\item {} Advisors
\item {} Pricing trend (up-round / down-round / flat)
\end{itemize}

\begin{quote}
``\textbf{Dates are required.} The announcement and close dates must always appear in the final slide's deal table. If only one date is available, show it and mark the other as ``—''.''
\end{quote}

\subparagraph{Step 4: Pull Company Context for Notable Deals}\mbox{}\par

For any company involved in a significant deal (large round, notable valuation shift), get a brief description:


\textbf{Structured Template}
\begin{itemize}[leftmargin=*,itemsep=2pt,topsep=2pt]
\item {} get\_\allowbreak{}company\_\allowbreak{}summary\_\allowbreak{}from\_\allowbreak{}identifiers(identifiers=[notable\_\allowbreak{}companies])
\end{itemize}

This adds context to the narrative (e.g., ``The company, an AI infrastructure startup founded in 2021, is expanding into...'').


\subparagraph{Step 5: Identify Highlights \& Trends}\mbox{}\par

Before designing the slide, analyze the data to surface the story:


\textbf{Flag as ``Notable'':}

\begin{itemize}[leftmargin=*,itemsep=2pt,topsep=2pt]
\item {} Rounds ≥ \$100M
\item {} Down rounds (pricing trend = down)
\item {} New unicorns (post-money valuation crossing \$1B)
\item {} Significant valuation jumps (post-money ≥ 2x the last known valuation)
\item {} Repeat raisers (same company raising again within 6 months)
\item {} Unusually large investor syndicates
\end{itemize}

\textbf{Identify Trends:}

\begin{itemize}[leftmargin=*,itemsep=2pt,topsep=2pt]
\item {} Total capital deployed this period vs. typical (if historical data available)
\item {} Which sub-sectors are hottest (most rounds, most capital)
\item {} Round stage distribution (is early-stage or late-stage dominating?)
\item {} Most active investors across the digest
\item {} Geographic concentration
\item {} Valuation trends (are pre-money valuations compressing or expanding?)
\end{itemize}

\textbf{Select Key Takeaways (3–5):} Distill the most important signals into 3–5 concise bullet-style takeaways. These are the centerpiece of the slide. Each takeaway should be one sentence, punchy, and data-backed.


Examples:

\begin{itemize}[leftmargin=*,itemsep=2pt,topsep=2pt]
\item {} ``AI sector raised \$2.4B across 8 rounds — 3x the prior week, led by a \$800M mega-round from [Company] at a \$12B post-money valuation.''
\item {} ``[Company] closed a \$200M Series D at \$3.5B pre-money, up from \$1.8B in its Series C — signaling strong demand for AI developer tools.''
\item {} ``Down-round activity ticked up: 2 of 6 late-stage rounds priced below prior valuations.''
\end{itemize}

\subparagraph{Step 6: Generate Company Logos}\mbox{}\par

For each company featured in the key takeaways or notable deals, generate a logo using a two-tier local pipeline. \textbf{Do not use Clearbit} (\emph{logo.clearbit.com}) — it is deprecated and consistently fails. External logo CDNs (Brandfetch, logo.dev, Google Favicons) require API keys or are blocked by network restrictions. Instead, use the following approach:


\subparagraph{Tier 1: \emph{simple-icons} npm Package (3,300+ Brand SVGs, No Network Required)}\mbox{}\par

The \emph{simple-icons} package bundles high-quality SVG icons for thousands of well-known brands. It works entirely offline — no API keys, no network calls. Install it alongside \emph{sharp} for SVG → PNG conversion:


\textbf{Structured Template}
\begin{itemize}[leftmargin=*,itemsep=2pt,topsep=2pt]
\item {} npm install simple-icons sharp
\end{itemize}

\textbf{Lookup strategy:}


\textbf{Structured Template}
\begin{itemize}[leftmargin=*,itemsep=2pt,topsep=2pt]
\item {} const si = require('simple-icons');
\item {} const sharp = require('sharp');
\item {} // Find an icon by exact title match (case-insensitive)
\item {} function findSimpleIcon(companyName) \{
\item {} // Try exact match first
\item {} for (const [key, val] of Object.entries(si)) \{
\item {} if (!key.startsWith('si') || !val || !val.title) continue;
\item {} if (val.title.toLowerCase() === companyName.toLowerCase()) return val;
\item {} \}
\item {} // Try without common suffixes (AI, Inc., Corp.)
\item {} const stripped = companyName.replace(/\textbackslash{}s*(AI|Inc\textbackslash{}.?|Corp\textbackslash{}.?|Ltd\textbackslash{}.?)\$/i, '').trim();
\item {} if (stripped !== companyName) \{
\item {} for (const [key, val] of Object.entries(si)) \{
\item {} if (!key.startsWith('si') || !val || !val.title) continue;
\item {} if (val.title.toLowerCase() === stripped.toLowerCase()) return val;
\item {} \}
\item {} \}
\item {} return null;
\item {} \}
\item {} // Convert SVG to PNG with the brand's official color
\item {} async function simpleIconToPng(icon, outputPath) \{
\item {} const coloredSvg = icon.svg.replace('<svg', <svg fill=``\#\$\{icon.hex\}'');
\item {} await sharp(Buffer.from(coloredSvg))
\item {} .resize(128, 128, \{ fit: 'contain', background: \{ r: 255, g: 255, b: 255, alpha: 0 \} \})
\item {} .png()
\item {} .toFile(outputPath);
\item {} \}
\end{itemize}

\textbf{Coverage:} \textasciitilde{}43\% of typical deal flow companies (strong for major tech brands like Stripe, Anthropic, Databricks, Snowflake, Discord, Shopify, SpaceX, Mistral AI, Hugging Face; weaker for niche fintech, biotech, or early-stage companies).


\subparagraph{Tier 2: Initial-Based Fallback via \emph{sharp} (100\% Coverage)}\mbox{}\par

For companies not found in \emph{simple-icons}, generate a clean initial-based logo as a PNG:


\textbf{Structured Template}
\begin{itemize}[leftmargin=*,itemsep=2pt,topsep=2pt]
\item {} async function generateInitialLogo(companyName, outputPath) \{
\item {} const initial = companyName.charAt(0).toUpperCase();
\item {} const svg =
\item {} <svg width=``128'' height=``128'' xmlns=``http://www.w3.org/2000/svg''>
\item {} <circle cx=``64'' cy=``64'' r=``64'' fill=``\#BDBDBD''/>
\item {} <text x=``64'' y=``64'' font-family=``Arial, Helvetica, sans-serif''
\item {} font-size=``56'' font-weight=``bold'' fill=``\#FFFFFF''
\item {} text-anchor=``middle'' dominant-baseline=``central''>\$\{initial\}</text>
\item {} </svg>;
\item {} await sharp(Buffer.from(svg)).png().toFile(outputPath);
\item {} \}
\end{itemize}

\subparagraph{Complete Pipeline}\mbox{}\par

\textbf{Structured Template}
\begin{itemize}[leftmargin=*,itemsep=2pt,topsep=2pt]
\item {} async function fetchLogo(companyName, outputDir) \{
\item {} const fileName = companyName.toLowerCase().replace(/[\textbackslash{}s.]+/g, '-') + '.png';
\item {} const outPath = path.join(outputDir, fileName);
\item {} // Tier 1: Try simple-icons
\item {} const icon = findSimpleIcon(companyName);
\item {} if (icon) \{
\item {} await simpleIconToPng(icon, outPath);
\item {} return \{ path: outPath, source: 'simple-icons' \};
\item {} \}
\item {} // Tier 2: Generate initial-based fallback
\item {} await generateInitialLogo(companyName, outPath);
\item {} return \{ path: outPath, source: 'initial-fallback' \};
\item {} \}
\end{itemize}

\textbf{Logo guidelines:}

\begin{itemize}[leftmargin=*,itemsep=2pt,topsep=2pt]
\item {} Save all logos to \emph{/home/claude/logos/[company-name].png}
\item {} All logos are 128×128 PNG with transparent backgrounds
\item {} On the slide, display logos at 0.35``–0.5'' tall — they're accents, not focal points
\item {} Initial-fallback circles use gray (\emph{BDBDBD}) fill with white text — consistent with the monochrome palette
\item {} Never mix logo styles randomly — if most companies resolve to brand icons, the few fallbacks should blend in naturally
\end{itemize}

\subparagraph{Step 7: Generate the One-Page PPTX}\mbox{}\par

Read \emph{/mnt/skills/public/pptx/SKILL.md} and \emph{/mnt/skills/public/pptx/pptxgenjs.md} before creating the slide.


Create a \textbf{single-slide} PowerPoint using \emph{pptxgenjs}. The slide should be information-dense but visually clean — think ``executive dashboard'' not ``wall of text.''


\subparagraph{Slide Layout}\mbox{}\par

\textbf{Structured Template}
\begin{itemize}[leftmargin=*,itemsep=2pt,topsep=2pt]
\item {} DEAL FLOW DIGEST
\item {} [Period] · [Sectors] [Date]
\item {} \$X.XB N \$X.XB \$X.XB
\item {} Raised Rounds Avg Pre Largest
\item {} KEY TAKEAWAYS
\item {} [Logo] Takeaway 1 text goes here...
\item {} [Logo] Takeaway 2 text goes here...
\item {} [Logo] Takeaway 3 text goes here...
\item {} [Logo] Takeaway 4 text goes here...
\item {} TOP DEALS
\item {} Company Type Announced Closed Amount Pre-\$ Post-\$ Lead 🔗
\item {} ... ... ... ... ... ... ... ... 🔗
\item {} [Footer: Deal Flow Digest · Sources: S\&P Global Capital IQ]
\item {} [Footer: AI Disclaimer]
\end{itemize}

\subparagraph{Design Specifications}\mbox{}\par

\textbf{Color philosophy: Minimal, monochrome-first.} The slide should feel like a high-end financial brief — black, white, and gray dominate. Color is used \textbf{only} where it carries meaning (e.g., a red indicator for a down round, a green indicator for a standout metric) or where the reader would naturally expect it (company logos). Never use color for purely decorative purposes like background fills, accent bars, or gradient effects.


\textbf{Color palette — Monochrome Executive:}

\begin{itemize}[leftmargin=*,itemsep=2pt,topsep=2pt]
\item {} Primary background: \emph{FFFFFF} (white) — clean, open slide background
\item {} Header bar: \emph{1A1A1A} (near-black) — strong contrast for the title region
\item {} Primary text: \emph{1A1A1A} (near-black) — all body text, stat numbers, takeaways
\item {} Secondary text: \emph{6B6B6B} (medium gray) — labels, captions, footer, date stamps
\item {} Borders \& dividers: \emph{D0D0D0} (light gray) — subtle structural lines, card outlines, table borders
\item {} Card backgrounds: \emph{F5F5F5} (off-white / very light gray) — stat card fills, alternating table rows
\item {} Link text: \emph{2B5797} (muted blue) — Capital IQ deal links in the table (the only blue on the slide)
\item {} \textbf{Semantic color (sparingly):}
\begin{itemize}[leftmargin=*,itemsep=2pt,topsep=2pt]
\item {} Down rounds or negative signals: \emph{C0392B} (muted red) — use only as a small dot, tag, or single-word highlight, never as a fill or background
\item {} Standout positive metrics (new unicorn, outsized round): \emph{2E7D32} (muted green) — same minimal usage: a dot, a small tag, or a single highlighted number
\item {} If no data points warrant a color indicator, \textbf{use no color at all}. A fully monochrome slide is perfectly correct.
\end{itemize}
\end{itemize}

\textbf{Typography:}

\begin{itemize}[leftmargin=*,itemsep=2pt,topsep=2pt]
\item {} Title: 28–32pt, bold, white on near-black header bar
\item {} Stat numbers: 36–44pt, bold, near-black
\item {} Stat labels: 10–12pt, medium gray (\emph{6B6B6B})
\item {} Takeaway text: 12–14pt, near-black, left-aligned
\item {} Table text: 9–11pt, near-black with gray (\emph{6B6B6B}) for secondary columns
\item {} Link text: 9–10pt, muted blue (\emph{2B5797})
\item {} Footer: 8pt, medium gray
\end{itemize}

\textbf{Stat Cards (top row):}

\begin{itemize}[leftmargin=*,itemsep=2pt,topsep=2pt]
\item {} 4 key metrics as large-number callouts: Total Raised, \# Rounds, Avg Pre-Money Valuation, Largest Round
\item {} Each in a card with \emph{F5F5F5} fill and a thin \emph{D0D0D0} border — no shadow, no color fills
\item {} If a stat is surprising or extreme (e.g., 3x normal volume, a record deal), a small colored dot or underline may be placed next to that single number — otherwise keep fully monochrome
\item {} If pre-money valuations are mostly undisclosed, substitute with a different metric (e.g., Median Round Size, \# New Unicorns)
\end{itemize}

\textbf{Key Takeaways (middle section):}

\begin{itemize}[leftmargin=*,itemsep=2pt,topsep=2pt]
\item {} 3–5 one-line takeaways, each prefixed with the relevant company logo (small, \textasciitilde{}0.35`` tall)
\item {} If no logo available, use a \textbf{gray circle} with the company initial in white — not a colored circle
\item {} Left-aligned, with enough spacing to breathe
\item {} Down-round or negative takeaways may use a small red dot prefix; otherwise no color
\item {} Include valuation context where available (e.g., ``at a \$5B post-money valuation'')
\end{itemize}

\textbf{Top Deals Table (bottom section):}

\begin{itemize}[leftmargin=*,itemsep=2pt,topsep=2pt]
\item {} Compact table showing the 4–6 most notable deals
\item {} Columns: Company, Type (Series X), Announced (date), Closed (date), Amount (\$M), Pre-Money (\$M), Post-Money (\$M), Lead Investor, Deal Link
\item {} \textbf{Announced} and \textbf{Closed} columns show dates in \emph{MMM DD} format (e.g., ``Jan 15''). These columns are required and must always be present. If a date is not available, show ``—''.
\item {} The \textbf{Deal Link} column contains a clickable ``View →'' text linking to Capital IQ: ``\emph{ https://www.capitaliq.spglobal.com/web/client?\#offering/capitalOfferingProfile?id=<transaction\_\allowbreak{}id> }`\emph{ where }<transaction\_\allowbreak{}id>\emph{ is the }transaction\_\allowbreak{}id\emph{ from }get\_\allowbreak{}rounds\_\allowbreak{}of\_\allowbreak{}funding\_\allowbreak{}from\_\allowbreak{}identifiers`.
\item {} If pre-money or post-money valuation is not disclosed, show ``—'' in that cell
\item {} Header row with near-black (\emph{1A1A1A}) fill and white text; alternating rows in \emph{F5F5F5} and \emph{FFFFFF}
\item {} \textbf{Center the table horizontally} on the slide. Calculate the table's total width, then set \emph{x} so it is centered within the slide width: \emph{x = (slideWidth - tableWidth) / 2}. For a 16:9 layout (13.33`` wide), if the table is 12'' wide, use \emph{x = 0.67}. Never left-align the table to the slide edge.
\item {} Keep it tight — this is a reference, not the focal point
\item {} No colored fills in table cells. If a deal is a down round, a small red text tag ``(↓ down)'' may appear next to the amount — that is the only permitted color in the table.
\end{itemize}

\textbf{Deal Link Implementation (pptxgenjs):} In pptxgenjs, hyperlinks are added to table cells using the \emph{options.hyperlink} property on the cell object:

\textbf{Structured Template}
\begin{itemize}[leftmargin=*,itemsep=2pt,topsep=2pt]
\item {} // Table cell with Capital IQ deal link
\item {} \{
\item {} text: ``View →'',
\item {} options: \{
\item {} hyperlink: \{
\item {} url: https://www.capitaliq.spglobal.com/web/client?\#offering/capitalOfferingProfile?id=\$\{transactionId\}
\item {} \},
\item {} color: ``2B5797'',
\item {} fontSize: 9,
\item {} fontFace: ``Arial''
\item {} \}
\item {} \}
\end{itemize}

\textbf{Table Centering (pptxgenjs):} Always center the deal table on the slide. Calculate the x position dynamically:

\textbf{Structured Template}
\begin{itemize}[leftmargin=*,itemsep=2pt,topsep=2pt]
\item {} const SLIDE\_\allowbreak{}W = 13.33; // 16:9 slide width
\item {} const TABLE\_\allowbreak{}W = 12.5; // total table width (sum of all column widths)
\item {} const TABLE\_\allowbreak{}X = (SLIDE\_\allowbreak{}W - TABLE\_\allowbreak{}W) / 2; // ≈ 0.42``
\item {} slide.addTable(tableRows, \{
\item {} x: TABLE\_\allowbreak{}X,
\item {} y: tableY,
\item {} w: TABLE\_\allowbreak{}W,
\item {} colW: [1.8, 0.9, 0.9, 0.9, 1.0, 1.1, 1.2, 1.6, 0.7], // Company, Type, Announced, Closed, Amount, Pre-\$, Post-\$, Lead, Link
\item {} // ... other options
\item {} \});
\end{itemize}
Adjust \emph{colW} values as needed, but always recompute \emph{TABLE\_\allowbreak{}X} from \emph{(SLIDE\_\allowbreak{}W - sum(colW)) / 2} to keep the table centered.


\textbf{Footer:}

\begin{itemize}[leftmargin=*,itemsep=2pt,topsep=2pt]
\item {} Small text in medium gray: ``Deal Flow Digest · [Period] · Sources: S\&P Global Capital IQ · Generated [Date]''
\end{itemize}

\textbf{General color rules (enforce strictly):}

\begin{itemize}[leftmargin=*,itemsep=2pt,topsep=2pt]
\item {} Company logos are the only ``full color'' elements on the slide — they appear as-is from the source.
\item {} Deal links use muted blue (\emph{2B5797}) — this is the only non-monochrome text color besides semantic red/green.
\item {} Outside of logos and links, the slide should look correct printed on a black-and-white printer.
\item {} Never apply color to backgrounds, accent bars, decorative shapes, or section dividers.
\item {} When in doubt, leave it gray.
\end{itemize}

\subparagraph{Code Structure}\mbox{}\par

\textbf{Structured Template}
\begin{itemize}[leftmargin=*,itemsep=2pt,topsep=2pt]
\item {} const pptxgen = require(``pptxgenjs'');
\item {} const pres = new pptxgen();
\item {} pres.layout = ``LAYOUT\_\allowbreak{}16x9'';
\item {} pres.title = ``Deal Flow Digest'';
\item {} const slide = pres.addSlide();
\item {} const SLIDE\_\allowbreak{}W = 13.33; // 16:9 slide width in inches
\item {} // 1. Dark header bar with title and period
\item {} // 2. Stat cards row (4 cards: Total Raised, \# Rounds, Avg Pre-Money, Largest Round)
\item {} // 3. Key takeaways section with logos (include valuation context)
\item {} // 4. Top deals table with Announced, Closed, Pre-Money, Post-Money columns and Capital IQ deal links
\item {} // - Center the table: x = (SLIDE\_\allowbreak{}W - tableWidth) / 2
\item {} // 5. Footer
\item {} pres.writeFile(\{ fileName: ``/home/claude/deal-flow-digest.pptx'' \});
\end{itemize}

Use factory functions (not shared objects) for shadows and repeated styles per the pptxgenjs pitfalls guidance.


\subparagraph{Step 8: QA the Slide}\mbox{}\par

Follow the QA process from the PPTX skill:


\begin{enumerate}[leftmargin=*,itemsep=2pt,topsep=2pt]
\item {} \textbf{Content QA:} \emph{python -m markitdown deal-flow-digest.pptx} — verify all text, numbers, company names, valuation figures, and deal links are correct
\item {} \textbf{Visual QA:} Convert to image and inspect: ``\emph{bash python /mnt/skills/public/pptx/scripts/office/soffice.py --headless --convert-to pdf deal-flow-digest.pptx pdftoppm -jpeg -r 200 deal-flow-digest.pdf slide }`` Check for overlapping elements, text overflow, alignment issues, low-contrast text, logo sizing problems, and that deal link text is visible.
\item {} \textbf{Link QA:} Verify that the Capital IQ URLs in the table are correctly formatted with the right transaction IDs.
\item {} \textbf{Fix and re-verify} — at least one fix-and-verify cycle before declaring done.
\end{enumerate}

\subparagraph{Step 9: Present Results}\mbox{}\par

\begin{enumerate}[leftmargin=*,itemsep=2pt,topsep=2pt]
\item {} Copy the final \emph{.pptx} to \emph{/mnt/user-data/outputs/}
\item {} Use \emph{present\_\allowbreak{}files} to share the slide
\item {} Provide a 2–3 sentence verbal summary:
\begin{itemize}[leftmargin=*,itemsep=2pt,topsep=2pt]
\item {} ``Your digest covers X rounds totaling \$Y raised across [sectors].''
\item {} Call out the single most notable deal and its valuation
\item {} Flag any concerning trends (down rounds, valuation compression, etc.)
\end{itemize}
\end{enumerate}

\subparagraph{Error Handling}\mbox{}\par

\subparagraph{Entity Resolution Failures}\mbox{}\par
\begin{itemize}[leftmargin=*,itemsep=2pt,topsep=2pt]
\item {} \textbf{Empty results for a known company:} First check \emph{get\_\allowbreak{}info\_\allowbreak{}from\_\allowbreak{}identifiers} — if that fails, try the alias from \emph{references/sector-seeds.md} or the \emph{company\_\allowbreak{}id} directly. Common brand→legal mismatches: Together AI → ``Together Computer, Inc.'', Character.ai → ``Character Technologies, Inc.'', Runway ML → ``Runway AI, Inc.''.
\item {} \textbf{Subsidiary companies:} DeepMind, GitHub, Instagram, WhatsApp, YouTube, BeReal, etc. are subsidiaries — they have zero independent funding rounds. Note these as ``acquired/subsidiary'' in context but do not report them as ``no activity.''
\item {} \textbf{Defunct companies:} Companies like Convoy (shut down Oct 2023) still resolve in S\&P Global but will never have new activity. The \emph{references/sector-seeds.md} file flags these — check it before including a company.
\item {} \textbf{@@TOKEN0@@ errors or returns zeros:} Fall back to \emph{get\_\allowbreak{}rounds\_\allowbreak{}of\_\allowbreak{}funding\_\allowbreak{}from\_\allowbreak{}identifiers} — the summary tool is less reliable. Never rely on the summary tool as the sole data source.
\item {} \textbf{Wrong @@TOKEN0@@ parameter:} If investor-perspective queries return empty, verify you're using \emph{company\_\allowbreak{}investing\_\allowbreak{}in\_\allowbreak{}round\_\allowbreak{}of\_\allowbreak{}funding}, not \emph{company\_\allowbreak{}raising\_\allowbreak{}funds} (and vice versa).
\end{itemize}

\subparagraph{Data Quality Issues}\mbox{}\par
\begin{itemize}[leftmargin=*,itemsep=2pt,topsep=2pt]
\item {} \textbf{No activity in period:} If a sector had zero funding rounds, note this explicitly on the slide (``No transactions recorded in [Sector] during the period'') — absence of activity is itself informative.
\item {} \textbf{Sparse valuation data:} If pre-money and post-money valuations are undisclosed for most transactions, note the data limitation in a footer annotation and use ``—'' in the table. Adjust the stat card to show a different metric (e.g., Median Round Size) instead of Avg Pre-Money.
\item {} \textbf{Logo retrieval failures:} The \emph{simple-icons} npm package provides \textasciitilde{}43\% coverage for typical deal flow companies. For the remainder, use the \emph{sharp}-generated initial-based fallback. Keep a consistent icon style — don't mix random approaches. If \emph{simple-icons} or \emph{sharp} fail to install, fall back to pptxgenjs shape-based initials (gray ellipse + white text overlay) which require no external dependencies.
\item {} \textbf{Too many deals for one slide:} If there are more than 6 notable deals, show the top 6 in the table and add a footnote: ``+N additional deals not shown.'' Prioritize by deal size.
\item {} \textbf{Large universes:} For multi-sector digests with 100+ companies, batch all API calls in groups of 15–20. Prioritize depth on notable deals over completeness on minor ones.
\item {} \textbf{Stale seeds:} If competitor expansion returns very few results for a sector, the seed companies may be too niche. Broaden by adding 2–3 more well-known names and re-expanding.
\item {} \textbf{Invalid transaction IDs for links:} If a \emph{transaction\_\allowbreak{}id} from the funding tool doesn't produce a valid Capital IQ URL, omit the link cell for that row rather than including a broken link.
\end{itemize}

\subparagraph{Example Prompts}\mbox{}\par

\begin{itemize}[leftmargin=*,itemsep=2pt,topsep=2pt]
\item {} ``Give me a weekly deal flow digest for AI and fintech''
\item {} ``Summarize this week's funding in biotech''
\item {} ``Deal roundup for my coverage — cybersecurity, cloud infrastructure, and dev tools — last 2 weeks''
\item {} ``What happened in venture this week across all sectors I follow?''
\item {} ``Quick deal flow slide for climate tech this month''
\end{itemize}

\vspace{0.8em}\hrule\vspace{0.8em}

\subsection{Skill Resource: spglobal/skills/funding-digest/references/sector-seeds.md}\label{file:spglobal-skills-funding-digest-references-sector-seeds.md}

\paragraph{Sector Seed Companies Reference}\mbox{}\par

When the user specifies a sector but not specific companies, use these seed lists to bootstrap the company universe. These are starting points — always expand via \emph{get\_\allowbreak{}competitors\_\allowbreak{}from\_\allowbreak{}identifiers} and validate with \emph{get\_\allowbreak{}info\_\allowbreak{}from\_\allowbreak{}identifiers}.


\begin{quote}
``\textbf{All seeds below have been validated against S\&P Global's identifier system.} If a seed fails to resolve, try the alias listed in parentheses before dropping it.''
\end{quote}

\subparagraph{Technology / Software}\mbox{}\par

\subparagraph{AI / Machine Learning}\mbox{}\par
Seeds: OpenAI, Anthropic, Databricks, Scale AI, Cohere, Hugging Face, Mistral AI, xAI, Perplexity AI, Runway ML (alias: ``Runway AI, Inc.''), Together AI (alias: ``Together Computer, Inc.''), Character.ai (alias: ``Character Technologies, Inc.''), Groq, Stability AI, Aleph Alpha, Magic AI


⚠ \textbf{Excluded (do not use as seeds):}

\begin{itemize}[leftmargin=*,itemsep=2pt,topsep=2pt]
\item {} \emph{Inflection AI} — Core team absorbed by Microsoft (Mar 2024). Historical rounds exist but no new activity.
\item {} \emph{Adept AI} — Largely absorbed by Amazon (2024). Same as above.
\item {} \emph{DeepMind} — Subsidiary of Alphabet. No independent funding rounds.
\end{itemize}

\subparagraph{Cybersecurity}\mbox{}\par
Seeds: CrowdStrike, Palo Alto Networks, Wiz, Snyk, SentinelOne, Abnormal Security, Netskope


\subparagraph{Cloud Infrastructure / DevTools}\mbox{}\par
Seeds: Snowflake, HashiCorp, Datadog, Confluent, Vercel, Supabase, PlanetScale


\subparagraph{Fintech}\mbox{}\par
Seeds: Stripe, Plaid, Brex, Ramp, Mercury, Affirm, Marqeta, Navan


\subparagraph{Vertical SaaS}\mbox{}\par
Seeds: ServiceTitan, Toast, Procore, Veeva Systems, Blend Labs


\subparagraph{Healthcare / Life Sciences}\mbox{}\par

\subparagraph{Biotech / Pharma}\mbox{}\par
Seeds: Moderna, BioNTech, Recursion Pharmaceuticals, Tempus AI, Insitro, AbCellera


\subparagraph{Digital Health}\mbox{}\par
Seeds: Teladoc, Hims \& Hers, Ro, Noom, Color Health


⚠ \textbf{Excluded (do not use as seeds):}

\begin{itemize}[leftmargin=*,itemsep=2pt,topsep=2pt]
\item {} \emph{Cerebral} — Still operating but has faced significant regulatory issues; include only if the user specifically requests it.
\end{itemize}

\subparagraph{Medical Devices}\mbox{}\par
Seeds: Intuitive Surgical, Butterfly Network, Outset Medical


⚠ \textbf{Excluded (do not use as seeds):}

\begin{itemize}[leftmargin=*,itemsep=2pt,topsep=2pt]
\item {} \emph{Shockwave Medical} — Acquired by Johnson \& Johnson (May 2024). Now a subsidiary with no independent funding rounds.
\end{itemize}

\subparagraph{Energy / Climate}\mbox{}\par

\subparagraph{Climate Tech}\mbox{}\par
Seeds: Redwood Materials, Form Energy, Commonwealth Fusion, Sila Nanotechnologies, Climeworks


\subparagraph{Clean Energy}\mbox{}\par
Seeds: Enphase Energy, First Solar, Rivian, QuantumScape, Sunnova


\subparagraph{Consumer}\mbox{}\par

\subparagraph{E-Commerce / Marketplace}\mbox{}\par
Seeds: Shopify, Faire, Whatnot, Fanatics


⚠ \textbf{Excluded (do not use as seeds):}

\begin{itemize}[leftmargin=*,itemsep=2pt,topsep=2pt]
\item {} \emph{Temu (PDD Holdings)} — PDD Holdings is a massive public conglomerate; its funding activity is captured via equity markets, not venture rounds.
\end{itemize}

\subparagraph{Consumer Social / Media}\mbox{}\par
Seeds: Discord, Reddit, Substack


⚠ \textbf{Excluded (do not use as seeds):}

\begin{itemize}[leftmargin=*,itemsep=2pt,topsep=2pt]
\item {} \emph{BeReal} — Acquired by Voodoo (Jun 2024). Now a subsidiary.
\item {} \emph{Lemon8} — The brand name ``Lemon8'' resolves in S\&P Global to a small Dutch company (Lemon8 B.V.), \textbf{not} the ByteDance social media app. ByteDance's apps are subsidiaries and do not have independent funding rounds. Do not use.
\end{itemize}

\subparagraph{Industrials / Logistics}\mbox{}\par

\subparagraph{Logistics / Supply Chain}\mbox{}\par
Seeds: Flexport, Samsara, Project44, FourKites


⚠ \textbf{Excluded (do not use as seeds):}

\begin{itemize}[leftmargin=*,itemsep=2pt,topsep=2pt]
\item {} \emph{Convoy} — Shut down operations (Oct 2023). The identifier still resolves and historical rounds are available, but no new activity will appear.
\end{itemize}

\subparagraph{Robotics / Automation}\mbox{}\par
Seeds: Figure AI, Agility Robotics, Locus Robotics, Symbotic, Covariant


\subparagraph{Space / Aerospace}\mbox{}\par
Seeds: SpaceX, Relativity Space, Rocket Lab, Planet Labs, Astra


\subparagraph{Identifier Alias Reference}\mbox{}\par

Some well-known brand names don't match S\&P Global's legal entity names. If a brand name returns empty results from \emph{get\_\allowbreak{}info\_\allowbreak{}from\_\allowbreak{}identifiers}, try the alias:


\begingroup
\small
\setlength{\tabcolsep}{2pt}
\begin{longtable}{>{\RaggedRight\arraybackslash\hspace{0pt}}p{0.320\textwidth}>{\RaggedRight\arraybackslash\hspace{0pt}}p{0.320\textwidth}>{\RaggedRight\arraybackslash\hspace{0pt}}p{0.320\textwidth}}
\toprule
{} \textbf{Brand Name} & \textbf{S\&P Global Legal Name} & \textbf{company\_\allowbreak{} id} \\
\midrule
{} Together AI & Together Computer, Inc. & C\_\allowbreak{} 1860042219 \\
{} Character.ai & Character Technologies, Inc. & C\_\allowbreak{} 1829047235 \\
{} Runway ML & Runway AI, Inc. & C\_\allowbreak{} 633706980 \\
{} Adept AI & Adept AI Labs Inc. & C\_\allowbreak{} 1780739313 \\
{} xAI & X.AI LLC & C\_\allowbreak{} 1863863313 \\
\bottomrule
\end{longtable}
\endgroup

\begin{quote}
``\textbf{Tip:} When a brand name fails, try \emph{get\_\allowbreak{}info\_\allowbreak{}from\_\allowbreak{}identifiers} with the legal name. If that also fails, the company may not be indexed yet. As a last resort, use the \emph{company\_\allowbreak{}id} directly as the identifier.''
\end{quote}

\subparagraph{Notes}\mbox{}\par

\begin{itemize}[leftmargin=*,itemsep=2pt,topsep=2pt]
\item {} These lists skew toward US-based companies. For geographic filtering (Europe, Asia, etc.), the competitor expansion step is especially important.
\item {} For niche sub-sectors not listed here, ask the user for 2–3 example companies to use as seeds.
\item {} Always validate seeds are still active/relevant — companies pivot, merge, or shut down.
\item {} \textbf{Refresh cadence:} These seeds should be reviewed quarterly. AI sector seeds in particular change rapidly due to acquisitions and new entrants.
\item {} Seeds marked as subsidiaries or acquired will still resolve in \emph{get\_\allowbreak{}info\_\allowbreak{}from\_\allowbreak{}identifiers} (status = ``Operating Subsidiary'') but will return zero funding rounds. Skip these for funding queries.
\end{itemize}

\vspace{0.8em}\hrule\vspace{0.8em}

\subsection{Skill: tear-sheet}\label{file:spglobal-skills-tear-sheet-SKILL.md}

\paragraph{File Metadata}\mbox{}\par
\begingroup
\small
\setlength{\tabcolsep}{2pt}
\begin{longtable}{>{\RaggedRight\arraybackslash\hspace{0pt}}p{0.480\textwidth}>{\RaggedRight\arraybackslash\hspace{0pt}}p{0.480\textwidth}}
\toprule
{} \textbf{Field} & \textbf{Value} \\
\midrule
{} name & tear-sheet \\
{} description & Generate professional company tear sheets using S\&P Capital IQ data via the Kensho LLM-ready API MCP server. Use this skill whenever the user asks for a tear sheet, company one-pager, company profile, fact sheet, company snapshot, or company overview document — especially when they mention a specific company name or ticker. Also trigger when users ask for equity research summaries, M\&A company profiles, corporate development target profiles, sales/ BD meeting prep documents, or any concise single-company financial summary. This skill supports four audience types: equity research, investment banking/ M\&A, corporate development, and sales/ business development. If the user doesn't specify an audience, ask. Works for both public and private companies. \\
\bottomrule
\end{longtable}
\endgroup


\paragraph{Financial Tear Sheet Generator}\mbox{}\par

Generate audience-specific company tear sheets by pulling live data from S\&P Capital IQ via the S\&P Global MCP tools and formatting the result as a professional Word document.


\subparagraph{Style Configuration}\mbox{}\par

These are sensible defaults. To customize for your firm's brand, modify this section — common changes include swapping the color palette, changing the font (Calibri is standard at many banks), and updating the disclaimer text.


\textbf{Colors:}

\begin{itemize}[leftmargin=*,itemsep=2pt,topsep=2pt]
\item {} Primary (header banner background, section header text): \#1F3864
\item {} Accent (signature section highlights): \#2E75B6
\item {} Table header row fill: \#D6E4F0
\item {} Table alternating row fill: \#F2F2F2
\item {} Table borders: \#CCCCCC
\item {} Header banner text: \#FFFFFF
\end{itemize}

\textbf{Typography (sizes in half-points for docx-js):}

\begin{itemize}[leftmargin=*,itemsep=2pt,topsep=2pt]
\item {} Font family: Arial
\item {} Company name: 18pt bold (size: 36)
\item {} Section headers: 11pt bold (size: 22), Primary color
\item {} Body text: 9pt (size: 18)
\item {} Table text: 8.5pt (size: 17)
\item {} Footer/disclaimer: 7pt italic (size: 14)
\item {} Per-template overrides are specified in each reference file's Formatting Notes.
\end{itemize}

\textbf{Company Header Banner:}

\begin{itemize}[leftmargin=*,itemsep=2pt,topsep=2pt]
\item {} The header is a navy (\#1F3864) banner spanning the full page width with company name in white.
\item {} \textbf{Below the banner, key-value pairs MUST be rendered in a two-column borderless table spanning the full page width.} Left column: company identifiers (ticker, HQ, founded, employees, sector). Right column: financial identifiers (market cap, EV, stock price, shares outstanding). Each cell contains a bold label and regular-weight value on the same line (e.g., ``\textbf{Market Cap} \$124.7B''). Do not left-justify all fields in a single column — this wastes horizontal space and looks unprofessional. The two-column spread is the single most important visual signal that distinguishes a professional tear sheet from a default document.
\begin{itemize}[leftmargin=*,itemsep=2pt,topsep=2pt]
\item {} \textbf{Implementation:} Create a 2-column table with \emph{borders: none} and \emph{shading: none} on all cells. Set column widths to 50\% each. Place left-column fields (ticker, HQ, founded, employees) as separate paragraphs in the left cell. Place right-column fields (market cap, EV, stock price, shares outstanding) in the right cell. Each field is a single paragraph: bold run for the label, regular run for the value.
\item {} The specific fields in each column vary by audience — see the reference file's header spec. The principle is always: spread across the page, not clumped left.
\end{itemize}
\item {} \textbf{Do not use a bordered table for the header key-value block.} Bordered tables are reserved for financial data only.
\item {} Key metrics in the header (market cap, EV, stock price) should be displayed as inline key-value pairs, not in a separate bordered table.
\end{itemize}

\textbf{Section Headers:}

\begin{itemize}[leftmargin=*,itemsep=2pt,topsep=2pt]
\item {} Each section header gets a horizontal rule (thin line, \#CCCCCC, 0.5pt) directly beneath it to create clean visual separation between sections.
\item {} \textbf{Render the rule as a bottom border on the header paragraph itself} — do not insert a separate paragraph element for the rule. A separate paragraph adds its own before/after spacing and causes excessive whitespace below section titles.
\item {} \textbf{Implementation:} In docx-js, apply a bottom border to the section header paragraph via \emph{paragraph.borders.bottom = \{ style: BorderStyle.SINGLE, size: 1, color: ``CCCCCC'' \}}. Do not use \emph{doc.addParagraph()} with a separate horizontal rule element. Do not use \emph{thematicBreak}. The border must be on the heading paragraph itself with 0pt spacing after, so the rule sits tight against the header text.
\item {} Spacing: 12pt before the header paragraph, 0pt after the header paragraph, 4pt before the next content element.
\end{itemize}

\textbf{Bullet Formatting:}

\begin{itemize}[leftmargin=*,itemsep=2pt,topsep=2pt]
\item {} Use a single bullet character (•) for all bulleted content across all tear sheet types. Do not mix •, -, ▸, or numbered lists within or across tear sheets.
\item {} \textbf{Synthesis/analysis bullets} (Earnings Highlights, Strategic Fit, Integration Considerations, Conversation Starters): indented block-style formatting with left indent 360 DXA (0.25``) and a hanging indent for the bullet character. These should be visually offset from body text — they're interpretive content and should look distinct from data tables and prose paragraphs.
\item {} \textbf{Informational bullets} within relationship sections: standard body indent (180 DXA), no hanging indent.
\item {} \textbf{Do not apply left-border accents to any bullet sections.} Left-border styling renders inconsistently in docx-js and creates visual artifacts. Use indentation and text size differentiation to distinguish signature sections instead.
\end{itemize}

\textbf{Tables (financial data only):}

\begin{itemize}[leftmargin=*,itemsep=2pt,topsep=2pt]
\item {} Header row: Table Header Fill (\#D6E4F0) with bold dark text
\item {} Body rows: alternating white / Table Alternating Fill (\#F2F2F2)
\item {} Borders: Table Border color (\#CCCCCC), thin (BorderStyle.SINGLE, size 1)
\item {} Cell padding: top/bottom 40 DXA, left/right 80 DXA
\item {} Right-align all numeric columns
\item {} Always use ShadingType.CLEAR (never SOLID — SOLID causes black backgrounds)
\end{itemize}

\textbf{Layout:}

\begin{itemize}[leftmargin=*,itemsep=2pt,topsep=2pt]
\item {} US Letter portrait, 0.75`` margins (1080 DXA all sides)
\end{itemize}

\textbf{Number formatting:}

\begin{itemize}[leftmargin=*,itemsep=2pt,topsep=2pt]
\item {} Currency: USD. Use millions unless company revenue > \$50B (then billions, one decimal). Label units in column headers (e.g., ``Revenue (\$M)''), not in individual cells.
\item {} \textbf{Table cells: plain numbers with commas, no dollar signs.} Example: a revenue cell shows ``4,916'' not ``\$4,916''. The column header carries the unit.
\item {} Fiscal years: actual years (FY2022, FY2023, FY2024), never relative labels (FY-2, FY-1).
\item {} Negatives: parentheses, e.g., (2.3\%)
\item {} Percentages: one decimal place
\item {} Large numbers: commas as thousands separators
\end{itemize}

\textbf{Footer (document footer, not inline):} Place the source attribution and disclaimer in the actual document footer (repeated on every page), not as inline body text at the bottom. The footer is exactly two lines, centered, on every page:

\begin{itemize}[leftmargin=*,itemsep=2pt,topsep=2pt]
\item {} Line 1: ``Data: S\&P Capital IQ via Kensho | Analysis: AI-generated | [Month Day, Year]''
\item {} Line 2: ``For informational purposes only. Not investment advice.''
\item {} Style: 7pt italic, centered, \#666666 text color
\item {} This footer text must be identical across all tear sheet types for the same company. Do not vary the wording by audience.
\item {} \textbf{This footer is required on every tear sheet, every audience type, every page.} Do not omit it.
\end{itemize}

\subparagraph{Component Functions}\mbox{}\par

\textbf{You MUST use these exact functions to create document elements. Do NOT write custom docx-js styling code.} Copy these functions into your generated Node script and call them. The Style Configuration prose above remains as documentation; these functions are the enforcement mechanism.


\textbf{Structured Template}
\begin{itemize}[leftmargin=*,itemsep=2pt,topsep=2pt]
\item {} const docx = require(``docx'');
\item {} const \{
\item {} Document, Paragraph, TextRun, Table, TableRow, TableCell,
\item {} WidthType, AlignmentType, BorderStyle, ShadingType,
\item {} Header, Footer, PageNumber, HeadingLevel, TableLayoutType,
\item {} convertInchesToTwip
\item {} \} = docx;
\item {} // Color constants
\item {} const COLORS = \{
\item {} PRIMARY: ``1F3864'',
\item {} ACCENT: ``2E75B6'',
\item {} TABLE\_\allowbreak{}HEADER\_\allowbreak{}FILL: ``D6E4F0'',
\item {} TABLE\_\allowbreak{}ALT\_\allowbreak{}ROW: ``F2F2F2'',
\item {} TABLE\_\allowbreak{}BORDER: ``CCCCCC'',
\item {} HEADER\_\allowbreak{}TEXT: ``FFFFFF'',
\item {} FOOTER\_\allowbreak{}TEXT: ``666666'',
\item {} \};
\item {} const FONT = ``Arial'';
\item {} // 1. createHeaderBanner
\item {} // Returns an array of docx elements: [banner paragraph, key-value table]
\item {} function createHeaderBanner(companyName, leftFields, rightFields) \{
\item {} // leftFields / rightFields: arrays of \{ label: string, value: string \}
\item {} const banner = new Paragraph(\{
\item {} children: [
\item {} new TextRun(\{
\item {} text: companyName,
\item {} bold: true,
\item {} size: 36, // 18pt
\item {} color: COLORS.HEADER\_\allowbreak{}TEXT,
\item {} font: FONT,
\item {} \}),
\item {} ],
\item {} shading: \{ type: ShadingType.CLEAR, color: ``auto'', fill: COLORS.PRIMARY \},
\item {} spacing: \{ after: 0 \},
\item {} alignment: AlignmentType.LEFT,
\item {} \});
\item {} function buildCellParagraphs(fields) \{
\item {} return fields.map(
\item {} (f) =>
\item {} new Paragraph(\{
\item {} children: [
\item {} new TextRun(\{ text: f.label + `` '', bold: true, size: 18, font: FONT \}),
\item {} new TextRun(\{ text: f.value, size: 18, font: FONT \}),
\item {} ],
\item {} spacing: \{ after: 40 \},
\item {} \})
\item {} );
\item {} \}
\item {} const noBorder = \{ style: BorderStyle.NONE, size: 0, color: ``FFFFFF'' \};
\item {} const noBorders = \{ top: noBorder, bottom: noBorder, left: noBorder, right: noBorder \};
\item {} const noShading = \{ type: ShadingType.CLEAR, color: ``auto'', fill: ``FFFFFF'' \};
\item {} const kvTable = new Table(\{
\item {} rows: [
\item {} new TableRow(\{
\item {} children: [
\item {} new TableCell(\{
\item {} children: buildCellParagraphs(leftFields),
\item {} width: \{ size: 50, type: WidthType.PERCENTAGE \},
\item {} borders: noBorders,
\item {} shading: noShading,
\item {} \}),
\item {} new TableCell(\{
\item {} children: buildCellParagraphs(rightFields),
\item {} width: \{ size: 50, type: WidthType.PERCENTAGE \},
\item {} borders: noBorders,
\item {} shading: noShading,
\item {} \}),
\item {} ],
\item {} \}),
\item {} ],
\item {} width: \{ size: 100, type: WidthType.PERCENTAGE \},
\item {} \});
\item {} return [banner, kvTable];
\item {} \}
\item {} // 2. createSectionHeader
\item {} // Returns a single Paragraph with bottom border rule
\item {} function createSectionHeader(text) \{
\item {} return new Paragraph(\{
\item {} children: [
\item {} new TextRun(\{
\item {} text: text,
\item {} bold: true,
\item {} size: 22, // 11pt
\item {} color: COLORS.PRIMARY,
\item {} font: FONT,
\item {} \}),
\item {} ],
\item {} spacing: \{ before: 240, after: 0 \}, // 12pt before, 0pt after
\item {} border: \{
\item {} bottom: \{ style: BorderStyle.SINGLE, size: 1, color: COLORS.TABLE\_\allowbreak{}BORDER \},
\item {} \},
\item {} \});
\item {} \}
\item {} // 3. createTable
\item {} // headers: string[], rows: string[][], options: \{ accentHeader?, fontSize? \}
\item {} function createTable(headers, rows, options = \{\}) \{
\item {} const fontSize = options.fontSize || 17; // 8.5pt default
\item {} const headerFill = options.accentHeader ? COLORS.ACCENT : COLORS.TABLE\_\allowbreak{}HEADER\_\allowbreak{}FILL;
\item {} const headerTextColor = options.accentHeader ? COLORS.HEADER\_\allowbreak{}TEXT : ``000000'';
\item {} const cellBorders = \{
\item {} top: \{ style: BorderStyle.SINGLE, size: 1, color: COLORS.TABLE\_\allowbreak{}BORDER \},
\item {} bottom: \{ style: BorderStyle.SINGLE, size: 1, color: COLORS.TABLE\_\allowbreak{}BORDER \},
\item {} left: \{ style: BorderStyle.SINGLE, size: 1, color: COLORS.TABLE\_\allowbreak{}BORDER \},
\item {} right: \{ style: BorderStyle.SINGLE, size: 1, color: COLORS.TABLE\_\allowbreak{}BORDER \},
\item {} \};
\item {} const cellMargins = \{ top: 40, bottom: 40, left: 80, right: 80 \};
\item {} function isNumeric(val) \{
\item {} if (typeof val !== ``string'') return false;
\item {} const cleaned = val.replace(/[,\$\%()]/g, ``'').trim();
\item {} return cleaned !== ``'' \&\& !isNaN(cleaned);
\item {} \}
\item {} // Header row
\item {} const headerRow = new TableRow(\{
\item {} children: headers.map(
\item {} (h) =>
\item {} new TableCell(\{
\item {} children: [
\item {} new Paragraph(\{
\item {} children: [
\item {} new TextRun(\{
\item {} text: h,
\item {} bold: true,
\item {} size: fontSize,
\item {} color: headerTextColor,
\item {} font: FONT,
\item {} \}),
\item {} ],
\item {} \}),
\item {} ],
\item {} shading: \{ type: ShadingType.CLEAR, color: ``auto'', fill: headerFill \},
\item {} borders: cellBorders,
\item {} margins: cellMargins,
\item {} \})
\item {} ),
\item {} \});
\item {} // Data rows with alternating shading
\item {} const dataRows = rows.map((row, rowIdx) => \{
\item {} const fill = rowIdx \% 2 === 1 ? COLORS.TABLE\_\allowbreak{}ALT\_\allowbreak{}ROW : ``FFFFFF'';
\item {} return new TableRow(\{
\item {} children: row.map((cell, colIdx) => \{
\item {} const align = colIdx > 0 \&\& isNumeric(cell)
\item {} ? AlignmentType.RIGHT
\item {} : AlignmentType.LEFT;
\item {} return new TableCell(\{
\item {} children: [
\item {} new Paragraph(\{
\item {} children: [
\item {} new TextRun(\{ text: cell, size: fontSize, font: FONT \}),
\item {} ],
\item {} alignment: align,
\item {} \}),
\item {} ],
\item {} shading: \{ type: ShadingType.CLEAR, color: ``auto'', fill: fill \},
\item {} borders: cellBorders,
\item {} margins: cellMargins,
\item {} \});
\item {} \}),
\item {} \});
\item {} \});
\item {} return new Table(\{
\item {} rows: [headerRow, ...dataRows],
\item {} width: \{ size: 100, type: WidthType.PERCENTAGE \},
\item {} \});
\item {} \}
\item {} // 4. createBulletList
\item {} // items: string[], style: ``synthesis'' | ``informational''
\item {} function createBulletList(items, style = ``synthesis'') \{
\item {} const indent =
\item {} style === ``synthesis''
\item {} ? \{ left: 360, hanging: 180 \} // 360 DXA left, hanging indent for bullet
\item {} : \{ left: 180 \}; // 180 DXA, no hanging
\item {} return items.map(
\item {} (item) =>
\item {} new Paragraph(\{
\item {} children: [
\item {} new TextRun(\{ text: ``• '', font: FONT, size: 18 \}),
\item {} new TextRun(\{ text: item, font: FONT, size: 18 \}),
\item {} ],
\item {} indent: indent,
\item {} spacing: \{ after: 60 \},
\item {} \})
\item {} );
\item {} \}
\item {} // 5. createFooter
\item {} // date: string (e.g., ``February 23, 2026'')
\item {} function createFooter(date) \{
\item {} return new Footer(\{
\item {} children: [
\item {} new Paragraph(\{
\item {} children: [
\item {} new TextRun(\{
\item {} text: Data: S\&P Capital IQ via Kensho | Analysis: AI-generated | \$\{date\},
\item {} italics: true,
\item {} size: 14, // 7pt
\item {} color: COLORS.FOOTER\_\allowbreak{}TEXT,
\item {} font: FONT,
\item {} \}),
\item {} ],
\item {} alignment: AlignmentType.CENTER,
\item {} \}),
\item {} new Paragraph(\{
\item {} children: [
\item {} new TextRun(\{
\item {} text: ``For informational purposes only. Not investment advice.'',
\item {} italics: true,
\item {} size: 14,
\item {} color: COLORS.FOOTER\_\allowbreak{}TEXT,
\item {} font: FONT,
\item {} \}),
\item {} ],
\item {} alignment: AlignmentType.CENTER,
\item {} \}),
\item {} ],
\item {} \});
\item {} \}
\end{itemize}

\textbf{Usage in generated scripts:}

\begin{enumerate}[leftmargin=*,itemsep=2pt,topsep=2pt]
\item {} Copy all functions and constants above into the generated Node.js script
\item {} Call \emph{createHeaderBanner(...)} instead of manually building banner paragraphs and tables
\item {} Call \emph{createSectionHeader(...)} for every section title — never manually set paragraph borders
\item {} Call \emph{createTable(...)} for \textbf{all} tabular data — financial summaries, trading comps, M\&A activity, relationship tables, funding history, etc. Pass \emph{\{ accentHeader: true \}} for M\&A activity tables (IB/M\&A template). For non-numeric tables (e.g., relationships, ownership), the function still works correctly — it only right-aligns cells that contain numeric values.
\item {} Call \emph{createBulletList(items, ``synthesis'')} for earnings highlights, strategic fit, integration considerations, and conversation starters
\item {} Call \emph{createBulletList(items, ``informational'')} for relationship entries
\item {} Pass \emph{createFooter(date)} to the Document constructor's \emph{footers.default} property
\end{enumerate}

\textbf{What these functions eliminate:}

\begin{itemize}[leftmargin=*,itemsep=2pt,topsep=2pt]
\item {} Black background tables (enforces \emph{ShadingType.CLEAR} everywhere)
\item {} Separate horizontal rule paragraphs under section headers (enforces \emph{border.bottom} on the paragraph itself)
\item {} Bordered key-value tables in headers (enforces \emph{borders: none})
\item {} Inconsistent bullet styles (enforces \emph{•} character only)
\item {} Missing footers (provides the exact footer structure)
\end{itemize}

\subparagraph{Workflow}\mbox{}\par

\subparagraph{Step 1: Identify Inputs}\mbox{}\par

Gather up to four things before proceeding:


\begin{enumerate}[leftmargin=*,itemsep=2pt,topsep=2pt]
\item {} \textbf{Company} — name or ticker. If only a ticker, resolve the full company name with an initial query (e.g., use the company info tool).
\item {} \textbf{Audience} — one of four types:
\begin{itemize}[leftmargin=*,itemsep=2pt,topsep=2pt]
\item {} \textbf{Equity Research} — for buy-side/sell-side analysts evaluating an investment
\item {} \textbf{IB / M\&A} — for bankers profiling a company in transaction context
\item {} \textbf{Corp Dev} — for internal strategic teams evaluating an acquisition target
\item {} \textbf{Sales / BD} — for commercial teams preparing for a client meeting
\end{itemize}
\item {} \textbf{Comparable companies} (optional) — if the user has specific comps in mind, note them. Otherwise the skill will identify peers from S\&P Global data. This matters for Equity Research, IB/M\&A, and Corp Dev tear sheets.
\item {} \textbf{Page length preference} (optional) — defaults vary by audience (see below), but the user can override.
\end{enumerate}

If the user doesn't specify an audience, ask.


\subparagraph{Step 2: Read the Audience-Specific Reference}\mbox{}\par

Read the corresponding reference file from this skill's directory:


\begin{itemize}[leftmargin=*,itemsep=2pt,topsep=2pt]
\item {} Equity Research → \emph{references/equity-research.md}
\item {} IB / M\&A → \emph{references/ib-ma.md}
\item {} Corp Dev → \emph{references/corp-dev.md}
\item {} Sales / BD → \emph{references/sales-bd.md}
\end{itemize}

Each reference defines sections, a query plan, formatting guidance, and page length defaults.


\subparagraph{Step 3: Pull Data via S\&P Global MCP}\mbox{}\par

\textbf{First:} Create the intermediate file directory:

\textbf{Structured Template}
\begin{itemize}[leftmargin=*,itemsep=2pt,topsep=2pt]
\item {} mkdir -p /tmp/tear-sheet/
\end{itemize}

Use the \textbf{S\&P Global} MCP tools (also known as the Kensho LLM-ready API). Claude will have access to structured tools for financial data, company information, market data, consensus estimates, earnings transcripts, M\&A transactions, and business relationships. The query plans in each reference file describe what data to retrieve for each section — map these to the appropriate S\&P Global tools available in the conversation.


\textbf{After each query step, immediately write the retrieved data to the intermediate file(s) specified in the reference file's query plan.} Do not defer writes — data written to disk is protected from context degradation in long conversations.


\textbf{Query strategy:} Each reference file includes a query plan with 4-6 data retrieval steps. These are starting points, not rigid constraints. Prioritize data completeness over minimizing calls:


\begin{itemize}[leftmargin=*,itemsep=2pt,topsep=2pt]
\item {} \textbf{Always pull 4 fiscal years of financial data}, even though only 3 years are displayed. The fourth (earliest) year is needed to compute YoY revenue growth for the first displayed year. Without it, the earliest year's growth rate will show ``N/A'' — which looks like missing data, not a design choice.
\item {} Execute the query plan as written, using whichever S\&P Global tools match the data needed.
\item {} If a tool call returns incomplete results, try alternative tools or narrower queries. For example, if company summary doesn't include segment detail, try the segments tool directly.
\item {} If a data point isn't returned after a targeted retry, move on — label it ``N/A'' or ``Not disclosed.''
\item {} Never fabricate data. If the tools don't return a number, do not estimate from training knowledge.
\end{itemize}

\textbf{User-specified comps:} If the user provided comparable companies, query financials and multiples for each comp explicitly. If no comps were provided, use whatever peer data the tools return, or identify peers from the company's sector using the competitors tool.


\textbf{Optional context from the user:} Listen for additional context the user provides naturally. If they mention who the acquirer is (``we're looking at this for our platform''), what they sell (``we sell data analytics to banks''), or who the likely buyers are (``this would be interesting to Salesforce or Microsoft''), incorporate that context into the relevant synthesis sections (Strategic Fit, Conversation Starters, Deal Angle). Don't prompt for this information — just use it if offered.


\textbf{Private company handling:} CIQ includes private company data, so query the same way. However, expect sparser results. When generating for a private company:

\begin{itemize}[leftmargin=*,itemsep=2pt,topsep=2pt]
\item {} Skip: stock price, 52-week range, beta, stock performance, consensus estimates, trading comps
\item {} Lean into: business overview, relationships, ownership structure, whatever financials are available
\item {} Note ``Private Company'' prominently in the header
\end{itemize}

\subparagraph{Step 3b: Calculate Derived Metrics}\mbox{}\par

After all data collection is complete and intermediate files are written, compute all derived metrics in a single dedicated pass. This is a calculation-only step — no new MCP queries.


\textbf{Read all intermediate files back into context}, then compute:


\begin{itemize}[leftmargin=*,itemsep=2pt,topsep=2pt]
\item {} \textbf{Margins:} Gross Margin \%, EBITDA Margin \%, FCF Margin \%, Operating Margin \%
\item {} \textbf{Growth rates:} YoY revenue growth, YoY segment revenue growth, YoY EPS growth
\item {} \textbf{Efficiency ratios:} FCF Conversion (FCF/EBITDA), R\&D as \% of Revenue, Capex as \% of Revenue
\item {} \textbf{Capital structure:} Net Debt (Total Debt − Cash \& Equivalents), Net Debt / EBITDA
\item {} \textbf{Segment mix:} Each segment's revenue as \% of consolidated total revenue (use consolidated revenue as denominator per Data Integrity Rule 8)
\end{itemize}

\textbf{Validation (moved from Arithmetic Validation):} During this calculation pass, enforce all arithmetic checks:


\begin{itemize}[leftmargin=*,itemsep=2pt,topsep=2pt]
\item {} \textbf{Margin calculations:} Verify EBITDA Margin = EBITDA / Revenue, Gross Margin = Gross Profit / Revenue, etc. If the computed margin doesn't match the raw numbers, use the computation from raw components.
\item {} \textbf{Growth rates:} Verify YoY growth = (Current − Prior) / Prior. Don't rely on pre-computed growth rates if you have the underlying values.
\item {} \textbf{Segment totals:} If showing revenue by segment, verify segments sum to total revenue (within rounding tolerance). If they don't, omit the total row rather than publishing inconsistent math.
\item {} \textbf{Percentage columns:} Verify ``\% of Total'' columns sum to \textasciitilde{}100\%.
\item {} \textbf{Valuation cross-checks:} If you show both EV and EV/Revenue, verify EV / Revenue ≈ the stated multiple.
\end{itemize}

If a validation fails: attempt recalculation from raw data. If still inconsistent, flag the metric as ``N/A'' rather than publishing incorrect numbers. Quiet math errors in a tear sheet destroy credibility.


\textbf{Write results} to \emph{/tmp/tear-sheet/calculations.csv} with columns: \emph{metric,value,formula,components}


Example rows:

\textbf{Structured Template}
\begin{itemize}[leftmargin=*,itemsep=2pt,topsep=2pt]
\item {} metric,value,formula,components
\item {} gross\_\allowbreak{}margin\_\allowbreak{}fy2024,72.4\%,gross\_\allowbreak{}profit/revenue,``9524/13159''
\item {} revenue\_\allowbreak{}growth\_\allowbreak{}fy2024,12.3\%,(current-prior)/prior,``13159/11716''
\item {} net\_\allowbreak{}debt\_\allowbreak{}fy2024,2150,total\_\allowbreak{}debt-cash,``4200-2050''
\end{itemize}

\subparagraph{Step 3c: Verify Data Files}\mbox{}\par

Before generating the document, verify that all intermediate files are present and populated.


\textbf{Read each intermediate file} via separate read operations and print a verification summary:


\textbf{Structured Template}
\begin{itemize}[leftmargin=*,itemsep=2pt,topsep=2pt]
\item {} === Tear Sheet Data Verification ===
\item {} company-profile.txt: ✓ (12 fields)
\item {} financials.csv: ✓ (36 rows)
\item {} segments.csv: ✓ (8 rows)
\item {} valuation.csv: ✓ (5 rows)
\item {} calculations.csv: ✓ (18 rows)
\item {} earnings.txt: ✓ (populated)
\item {} relationships.txt: ⚠ MISSING
\item {} peer-comps.csv: ✓ (12 rows)
\end{itemize}

\textbf{Soft gate:} If any file expected for the current audience type is missing or empty, print a warning but continue. The tear sheet handles missing data gracefully with ``N/A'' and section skipping. However, the warning ensures visibility into what data was lost.


\textbf{Critical rule: The files — not your memory of earlier conversation — are the single source of truth for every number in the document.} When generating the DOCX in Step 4, read values from the intermediate files. Do not rely on conversation context for financial data.


\subparagraph{Step 4: Format as DOCX}\mbox{}\par

Read \emph{/mnt/skills/public/docx/SKILL.md} for docx creation mechanics (docx-js via Node). Apply the Style Configuration above plus the section-specific formatting in the reference file.


\textbf{Page length defaults (user can override):}

\begin{itemize}[leftmargin=*,itemsep=2pt,topsep=2pt]
\item {} Equity Research: 1 page (density is the convention)
\item {} IB / M\&A: 1-2 pages
\item {} Corp Dev: 1-2 pages
\item {} Sales / BD: 1-2 pages
\end{itemize}

If content exceeds the target, each reference file specifies which sections to cut first.


\textbf{Output filename:} \emph{[CompanyName]\_\allowbreak{}TearSheet\_\allowbreak{}[Audience]\_\allowbreak{}[YYYYMMDD].docx} Example: \emph{Nvidia\_\allowbreak{}TearSheet\_\allowbreak{}CorpDev\_\allowbreak{}20260220.docx}


Save to \emph{/mnt/user-data/outputs/} and present to the user.


\subparagraph{Data Integrity Rules}\mbox{}\par

These override everything else:

\begin{enumerate}[leftmargin=*,itemsep=2pt,topsep=2pt]
\item {} \textbf{S\&P Global tools are the only source for financial data.} Do not fill gaps with training knowledge — it may be stale or wrong.
\item {} \textbf{Label what you can't find.} Use ``N/A'' or ``Not disclosed'' rather than omitting a row silently.
\item {} \textbf{Dates matter.} Note the fiscal year end or reporting period. Don't assume calendar year = fiscal year. Market data (stock prices, market cap) should include an ``as of'' date.
\item {} \textbf{Don't mix reporting periods.} If you have FY2023 revenue and LTM EBITDA, label them distinctly.
\item {} \textbf{Prefer MCP-returned fields over manual computation.} If the S\&P Global tools return a pre-computed field (e.g., net debt, EBITDA, FCF), use that value directly rather than computing it from components. Only compute derived metrics manually when the tools do not return the field. This reduces discrepancies.
\item {} \textbf{Ensure consistency across tear sheet types.} If generating multiple tear sheets for the same company (e.g., equity research and IB/M\&A in the same session), the same underlying data points must produce identical values across all outputs. Net debt, revenue, EBITDA, margins, and growth rates must match exactly. Do not re-query or re-compute independently per report — reuse the same retrieved values.
\item {} \textbf{Never downgrade known transaction values.} If the M\&A tools return a deal value for a transaction, that value must appear in the output. Do not replace a known deal value with ``Undisclosed.'' Use ``Undisclosed'' only when the tools genuinely return no value for a transaction.
\item {} \textbf{Use consolidated revenue as the denominator for segment percentages.} When computing ``\% of Total'' for segment tables, divide each segment's revenue by consolidated total revenue (as reported on the income statement), not by the sum of segment revenues. The sum of segments often exceeds consolidated revenue due to intersegment eliminations. Using consolidated revenue ensures percentages align with the total revenue figure shown elsewhere in the document.
\item {} \textbf{Always include forward (NTM) multiples when available.} If the tools return both trailing and forward valuation multiples, both must appear in the output. Forward multiples are the primary valuation reference for equity research, IB/M\&A, and corp dev audiences. Never show only trailing multiples when forward data is available.
\item {} \textbf{No S\&P Global tool returns executive or management data.} Do not populate management names, titles, or biographical details from training data — this violates Rule 1 and produces stale information. If a management section appears in a template, omit it entirely. Ownership structure (institutional holders, insider \%, PE sponsor) may be included only if returned by the tools — gate with ``data permitting.''
\end{enumerate}

\subparagraph{Intermediate File Rule}\mbox{}\par

All data retrieved from MCP tools must be persisted to structured intermediate files before document generation. These files — not conversation context — are the single source of truth for every number in the document.


\textbf{Setup:} At the start of Step 3, create the working directory:

\textbf{Structured Template}
\begin{itemize}[leftmargin=*,itemsep=2pt,topsep=2pt]
\item {} mkdir -p /tmp/tear-sheet/
\end{itemize}

\textbf{Write-after-query mandate:} After each MCP query step completes, immediately write the retrieved data to the appropriate intermediate file(s). Do not wait until all queries finish. Each reference file's query plan specifies which file(s) to write after each step.


\textbf{File schemas:}


\begingroup
\small
\setlength{\tabcolsep}{2pt}
\begin{longtable}{>{\RaggedRight\arraybackslash\hspace{0pt}}p{0.240\textwidth}>{\RaggedRight\arraybackslash\hspace{0pt}}p{0.240\textwidth}>{\RaggedRight\arraybackslash\hspace{0pt}}p{0.240\textwidth}>{\RaggedRight\arraybackslash\hspace{0pt}}p{0.240\textwidth}}
\toprule
{} \textbf{File} & \textbf{Format} & \textbf{Columns /  Structure} & \textbf{Used By} \\
\midrule
{} \emph{/ tmp/ tear-sheet/ company-profile.txt} & Key-value text & name, ticker, exchange, HQ, sector, industry, founded, employees, market\_\allowbreak{} cap, enterprise\_\allowbreak{} value, stock\_\allowbreak{} price, 52wk\_\allowbreak{} high, 52wk\_\allowbreak{} low, shares\_\allowbreak{} outstanding, beta, ownership & All \\
{} \emph{/ tmp/ tear-sheet/ financials.csv} & CSV & \emph{period,line\_\allowbreak{} item,value,source} & All \\
{} \emph{/ tmp/ tear-sheet/ segments.csv} & CSV & \emph{period,segment\_\allowbreak{} name,revenue,source} & ER, IB, CD \\
{} \emph{/ tmp/ tear-sheet/ valuation.csv} & CSV & \emph{metric,trailing,forward,source} & ER, IB, CD \\
{} \emph{/ tmp/ tear-sheet/ consensus.csv} & CSV & \emph{metric,fy\_\allowbreak{} year,value,source} & ER \\
{} \emph{/ tmp/ tear-sheet/ earnings.txt} & Structured text & Quarter, date, key quotes, guidance, key drivers & ER, IB, Sales \\
{} \emph{/ tmp/ tear-sheet/ relationships.txt} & Structured text & Customers, suppliers, partners, competitors — each with descriptors & IB, CD, Sales \\
{} \emph{/ tmp/ tear-sheet/ peer-comps.csv} & CSV & \emph{ticker,metric,value,source} & ER, IB, CD \\
{} \emph{/ tmp/ tear-sheet/ ma-activity.csv} & CSV & \emph{date,target,deal\_\allowbreak{} value,type,rationale,source} & IB, CD \\
{} \emph{/ tmp/ tear-sheet/ calculations.csv} & CSV & \emph{metric,value,formula,components} & All (written in Step 3b) \\
\bottomrule
\end{longtable}
\endgroup

\textbf{Abbreviations:} ER = Equity Research, IB = IB/M\&A, CD = Corp Dev, Sales = Sales/BD.


Not every audience type uses every file — the reference files define which query steps apply. Files not relevant to the current audience type need not be created.


\textbf{Raw values only.} Intermediate files store raw values as returned by the tools. Do not pre-compute margins, growth rates, or other derived metrics in these files — that happens in Step 3b.


\textbf{Page budget enforcement:} Each reference file specifies a default page length and a numbered cut order. If the rendered document exceeds the target, apply cuts in the order specified — do not attempt to shrink font sizes or margins below the template minimums. The cut order is a strict priority stack: cut section 1 completely before touching section 2.


\subparagraph{Content Quality Rules}\mbox{}\par

\begin{enumerate}[leftmargin=*,itemsep=2pt,topsep=2pt]
\item {} \textbf{Rewrite every narrative section for the audience.} The CIQ company summary is an input, not an output. Each audience type needs a different description: concise and thesis-oriented for equity research, pitchbook prose for IB, product-focused for Corp Dev, plain language for Sales/BD. Never paste the CIQ summary verbatim into any tear sheet.
\item {} \textbf{Differentiate earnings highlights by audience.} The same earnings call produces different takeaways for different readers. Equity research wants segment-level performance and consensus beat/miss. IB wants margin trajectory and strategic commentary. Sales/BD wants strategic themes that create conversation angles. Do not reuse the same bullets across tear sheet types.
\item {} \textbf{Synthesis sections are the differentiator.} Strategic Fit Analysis, Integration Considerations, Conversation Starters, and Business Overview paragraphs are where the tear sheet earns its value. These sections require analytical reasoning that connects data points into a narrative — listing company names without context is not synthesis.
\item {} \textbf{Flag pending divestitures in segment tables.} If a company has announced a pending divestiture of a segment or business unit, add a footnote or parenthetical to the segment table noting the pending transaction (e.g., ``Mobility\emph{ — }Pending divestiture, expected mid-2026''). For Corp Dev and IB/M\&A tear sheets, include a one-line note below the segment table showing pro-forma revenue and revenue mix excluding the divested segment. This helps the reader evaluate the ``go-forward'' business without doing the math themselves.
\end{enumerate}

\subparagraph{Arithmetic Validation}\mbox{}\par

\textbf{→ Arithmetic validation is now enforced in Step 3b (Calculate Derived Metrics).} All margin calculations, growth rates, segment totals, percentage columns, and valuation cross-checks are validated during the dedicated calculation pass, before document generation begins. See Step 3b for the full validation checklist.


\vspace{0.8em}\hrule\vspace{0.8em}

\subsection{Skill Resource: spglobal/skills/tear-sheet/references/corp-dev.md}\label{file:spglobal-skills-tear-sheet-references-corp-dev.md}

\paragraph{Corporate Development Tear Sheet}\mbox{}\par

\subparagraph{Purpose}\mbox{}\par
A target evaluation profile for an internal Corp Dev team assessing a potential acquisition. Unlike the IB/M\&A profile (which tells a strategic \emph{story} for a pitch), this is an analytical document focused on strategic fit: how the target's products, customers, technology, and operations overlap with or complement the acquirer's. The Corp Dev audience already knows their own company — they need to understand the target through the lens of ``should we buy this, and what would integration look like?''


\textbf{Default page length:} 1-2 pages. Corp Dev teams are comfortable with denser profiles and won't penalize a second page if the content is substantive.


\subparagraph{Query Plan}\mbox{}\par

Start with these queries. Break apart and follow up if results are incomplete.


\textbf{Query 1 — Target profile:} ``[Company] business description products services technology platform headquarters founded employees sector industry'' → Header, Business \& Product Overview → \textbf{Immediately write} to \emph{/tmp/tear-sheet/company-profile.txt}


\textbf{Query 2 — Financials:} ``[Company] annual income statement revenue gross profit EBITDA operating income net income capex free cash flow R\&D expense total debt cash and equivalents last 4 fiscal years'' → Financial Summary (pull 4 years; display 3; use the earliest year only for YoY growth computation. R\&D is especially important for Corp Dev — they want to understand the target's investment in IP) → \textbf{Immediately write} raw values to \emph{/tmp/tear-sheet/financials.csv}


\textbf{Query 3 — Segments + customers:} ``[Company] revenue by segment business unit last 2 fiscal years key customers end markets customer concentration'' → Revenue Mix (need 2 years for YoY growth), Customer Analysis → \textbf{Immediately write} segment data to \emph{/tmp/tear-sheet/segments.csv} (skip if no segment data returned)


\textbf{Query 4 — Relationships + competitive landscape:} ``[Company] key customers suppliers partners competitors business relationships technology vendors'' → Strategic Fit Analysis, Ecosystem Map → \textbf{Immediately write} to \emph{/tmp/tear-sheet/relationships.txt}


\textbf{Query 5 — Valuation + peers:} If user provided specific comps: ``[Comp company] enterprise value EV/Revenue EV/EBITDA revenue growth EBITDA margin'' Otherwise: Use the competitors tool to identify 3-5 public peers, then pull EV/Revenue (NTM), EV/EBITDA (NTM), revenue growth \%, and EBITDA margin \% for each. Also: ``[Company] enterprise value market capitalization valuation multiples'' → Valuation Context → \textbf{Immediately write} company multiples to \emph{/tmp/tear-sheet/valuation.csv} → \textbf{Immediately write} peer data to \emph{/tmp/tear-sheet/peer-comps.csv}


\textbf{Query 6 — Ownership (data permitting):} ``[Company] ownership structure investors institutional ownership insider ownership'' → Ownership snapshot (no intermediate file — included in company-profile.txt if available)


\subparagraph{Sections}\mbox{}\par

Listed in priority order. See ``Page Budget \& Cut Order'' below for explicit cut guidance.


\subparagraph{1. Company Header}\mbox{}\par
Similar to IB/M\&A but adds a ``Strategic Relevance'' tagline.


\begingroup
\small
\setlength{\tabcolsep}{2pt}
\begin{longtable}{>{\RaggedRight\arraybackslash\hspace{0pt}}p{0.480\textwidth}>{\RaggedRight\arraybackslash\hspace{0pt}}p{0.480\textwidth}}
\toprule
{} \textbf{Field} & \textbf{Notes} \\
\midrule
{} Company name & Large, bold \\
{} Ticker \& Exchange & If public; omit for private \\
{} HQ Location &  \\
{} Founded /  Employees &  \\
{} Enterprise Value & Market cap for public; latest valuation for private \\
{} Ownership & Public /  PE-backed /  Founder-led /  etc. \\
{} \textbf{Strategic Relevance} & One sentence: why this target is attractive to an acquirer (synthesized from data) \\
\bottomrule
\end{longtable}
\endgroup

The ``Strategic Relevance'' line is unique to Corp Dev. Write from a neutral analytical perspective, not from the target's marketing language. Frame what an acquirer \emph{gets} — not what the target \emph{is}. Good: ``Dominant position in commodity price benchmarks with regulatory moat and 70\%+ gross margins — a capital-light, high-barrier data asset.'' Bad: ``The world's leading provider of financial data and analytics.'' The former tells a corp dev team why this is attractive; the latter is a press release.


\subparagraph{2. Business \& Product Overview}\mbox{}\par
\textbf{Do not paste the CIQ company summary.} Rewrite in 4-6 sentences (one paragraph, not two) with a product and technology focus — this audience evaluates acquisitions, not investments:

\begin{itemize}[leftmargin=*,itemsep=2pt,topsep=2pt]
\item {} Core products/services and what problem they solve
\item {} Technology or platform architecture (at a high level)
\item {} Key differentiators and competitive moat
\item {} Customer segments — who buys this and why
\item {} Current strategic direction
\end{itemize}

Write with an acquirer's lens. The reader is thinking: ``How would this fit with what we already have?'' The description should make it easy to spot overlaps and complements.


\subparagraph{3. Revenue Mix \& Customer Profile}\mbox{}\par
Two sub-sections:


\textbf{a) Revenue by Segment} (if available):


\begingroup
\small
\setlength{\tabcolsep}{2pt}
\begin{longtable}{>{\RaggedRight\arraybackslash\hspace{0pt}}p{0.240\textwidth}>{\RaggedRight\arraybackslash\hspace{0pt}}p{0.240\textwidth}>{\RaggedRight\arraybackslash\hspace{0pt}}p{0.240\textwidth}>{\RaggedRight\arraybackslash\hspace{0pt}}p{0.240\textwidth}}
\toprule
{} \textbf{Segment} & \textbf{Revenue (\$M)} & \textbf{\% of Total} & \textbf{YoY Growth} \\
\midrule
\bottomrule
\end{longtable}
\endgroup

\textbf{b) Customer Analysis:} Below the segment table (or in place of it if segments aren't available), write a qualitative paragraph covering:

\begin{itemize}[leftmargin=*,itemsep=2pt,topsep=2pt]
\item {} Customer concentration: are the top 5 customers >30\% of revenue? (Critical for acquisition risk)
\item {} Customer type: enterprise vs. SMB vs. consumer
\item {} Contract structure: subscription vs. transactional vs. license (if known)
\item {} Geographic mix
\end{itemize}

Customer concentration is the single most important signal here — if the tools return any data on top customers or revenue concentration, highlight it prominently. If data is sparse, write what you can infer from the business description and relationships. Keep this to 2-3 sentences — it's context for the segment table, not a standalone analysis.


\subparagraph{4. Strategic Fit Analysis}\mbox{}\par
\textbf{This is the signature section of the Corp Dev tear sheet.} It is required — do not skip or compress it. It uses the Business Relationships data from S\&P Capital IQ to map overlaps and complements.


Organize as three buckets. \textbf{Each bucket must be 2-3 sentences of analytical reasoning, not a list of company names.} Listing ``Customers: Microsoft, JPMorgan, Google'' without context is useless. Instead, explain \emph{what the overlap or complement means} for an acquisition.


\textbf{Customer Overlap:} Do the target's customers also buy from the acquirer (or the acquirer's industry)? Shared customers = easier cross-sell, but also risk of channel conflict. Name shared relationships and explain the implications.


\textbf{Technology / Product Complement:} Does the target fill a gap in the acquirer's portfolio? What does their tech stack look like? Who are their technology vendors? If they're already using the acquirer's technology, explain why that's a strong integration signal.


\textbf{Competitive Displacement:} Is the target a competitor, or does acquiring them remove a competitor from the market? Name the target's key competitors and explain how the competitive landscape would change post-acquisition.


This section requires synthesis — combine relationship data with the business overview to draw connections. If the user has told you who the acquirer is, tailor to that acquirer. If not, write it generically but still analytically.


\textbf{Length discipline:} The 2-3 sentence target per bucket above is a hard constraint, not a suggestion. Each bullet should front-load the insight, then support with one specific data point. The actual outputs frequently balloon to 5-7 sentences per bullet — this is the primary driver of page 3 spillover. If you find yourself writing more than 3 sentences, move the excess detail to Integration Considerations.


\subparagraph{5. Financial Summary (3-Year)}\mbox{}\par
Similar to IB/M\&A but with added emphasis on R\&D and capital efficiency.


\begingroup
\small
\setlength{\tabcolsep}{2pt}
\begin{longtable}{>{\RaggedRight\arraybackslash\hspace{0pt}}p{0.240\textwidth}>{\RaggedRight\arraybackslash\hspace{0pt}}p{0.240\textwidth}>{\RaggedRight\arraybackslash\hspace{0pt}}p{0.240\textwidth}>{\RaggedRight\arraybackslash\hspace{0pt}}p{0.240\textwidth}}
\toprule
{} \textbf{Metric (\$M)} & \textbf{FY20XX} & \textbf{FY20XX} & \textbf{FY20XX} \\
\midrule
{} Revenue &  &  &  \\
{} Revenue Growth \% &  &  &  \\
{} Gross Margin \% &  &  &  \\
{} EBITDA &  &  &  \\
{} EBITDA Margin \% &  &  &  \\
{} R\&D Expense &  &  &  \\
{} R\&D as \% of Revenue &  &  &  \\
{} Capex &  &  &  \\
{} Free Cash Flow &  &  &  \\
{} FCF Conversion (FCF/ EBITDA) &  &  &  \\
\bottomrule
\end{longtable}
\endgroup

R\&D intensity signals how much of the target's value is in IP vs. services. FCF conversion signals integration complexity — high-capex businesses are harder to integrate. If R\&D expense isn't separately reported, note ``Not separately disclosed'' in the row rather than ``N/A'' — this itself is informative (suggests R\&D may be embedded in COGS or SG\&A).


\textbf{Capital Structure} (compact sub-table below the financial summary):


\begingroup
\small
\setlength{\tabcolsep}{2pt}
\begin{longtable}{>{\RaggedRight\arraybackslash\hspace{0pt}}p{0.480\textwidth}>{\RaggedRight\arraybackslash\hspace{0pt}}p{0.480\textwidth}}
\toprule
{} \textbf{Metric} & \textbf{Value} \\
\midrule
{} Total Debt & \$XXM \\
{} Cash \& Equivalents & \$XXM \\
{} Net Debt & \$XXM \\
{} Net Debt /  EBITDA & X.Xx \\
\bottomrule
\end{longtable}
\endgroup

Frame for acquisition context: below the table, include one sentence interpreting leverage (e.g., ``Net debt of \$X represents Y.Yx EBITDA leverage, which is manageable for an investment-grade acquirer but would compound leverage in an LBO scenario.''). Pull from the balance sheet data already retrieved.


\subparagraph{6. Valuation Context}\mbox{}\par
\textbf{This section is required for public companies.} Corp Dev teams need market context to frame a potential bid. If the tools return valuation multiples, this section must appear. Do not cut it for space — cut Ownership Snapshot (Section 7) first.


Not a formal valuation — Corp Dev teams do their own DCFs. This provides market context.


\textbf{a) Trading multiples} (if public):


\textbf{Peer context is essential, not optional.} If the user provided comps, show each in its own column. If no comps were provided, use the competitors tool to identify 3-5 public peers, then pull EV/Revenue (NTM), EV/EBITDA (NTM), revenue growth \%, and EBITDA margin \% for each.


\begingroup
\small
\setlength{\tabcolsep}{2pt}
\begin{longtable}{>{\RaggedRight\arraybackslash\hspace{0pt}}p{0.153\textwidth}>{\RaggedRight\arraybackslash\hspace{0pt}}p{0.153\textwidth}>{\RaggedRight\arraybackslash\hspace{0pt}}p{0.153\textwidth}>{\RaggedRight\arraybackslash\hspace{0pt}}p{0.153\textwidth}>{\RaggedRight\arraybackslash\hspace{0pt}}p{0.153\textwidth}>{\RaggedRight\arraybackslash\hspace{0pt}}p{0.153\textwidth}}
\toprule
{} \textbf{Metric} & \textbf{[Company]} & \textbf{[Peer 1]} & \textbf{[Peer 2]} & \textbf{[Peer 3]} & \textbf{Peer Median} \\
\midrule
{} EV/ Revenue (NTM) &  &  &  &  &  \\
{} EV/ EBITDA (NTM) &  &  &  &  &  \\
{} Revenue Growth \% &  &  &  &  &  \\
{} EBITDA Margin \% &  &  &  &  &  \\
\bottomrule
\end{longtable}
\endgroup

A company's multiples in isolation tell a corp dev team nothing — they need relative context to assess whether the target is cheap or rich vs. the peer set. This is the minimum viable valuation context for a target profile.


If peer multiples are unavailable from the tools, show the company's own multiples and list selected peer names below the table for the reader's reference.


\textbf{b) Precedent transactions (data permitting):} If the M\&A tools return recent transactions involving the company's peers or within the same sub-sector, display them:


\begingroup
\small
\setlength{\tabcolsep}{2pt}
\begin{longtable}{>{\RaggedRight\arraybackslash\hspace{0pt}}p{0.153\textwidth}>{\RaggedRight\arraybackslash\hspace{0pt}}p{0.153\textwidth}>{\RaggedRight\arraybackslash\hspace{0pt}}p{0.153\textwidth}>{\RaggedRight\arraybackslash\hspace{0pt}}p{0.153\textwidth}>{\RaggedRight\arraybackslash\hspace{0pt}}p{0.153\textwidth}>{\RaggedRight\arraybackslash\hspace{0pt}}p{0.153\textwidth}}
\toprule
{} \textbf{Date} & \textbf{Target} & \textbf{Acquirer} & \textbf{EV (\$M)} & \textbf{EV/ Rev} & \textbf{EV/ EBITDA} \\
\midrule
\bottomrule
\end{longtable}
\endgroup

Frame as context: ``Recent transactions in [sub-sector] have valued targets at X–Yx revenue.''


If precedent transaction data is not available from the tools, note ``Precedent transaction data not available from data source'' rather than omitting silently. Do not fabricate precedent multiples.


For private companies, skip trading multiples. Precedent transactions (data permitting) are still useful.


\subparagraph{7. Ownership Snapshot (data permitting — omit if tools return nothing)}\mbox{}\par
If the tools return ownership data, include a compact block: founder/family control \%, PE sponsor details, institutional concentration. Ownership structure directly impacts deal complexity and is relevant for integration considerations.


\textbf{Do not include a Management Team table.} No S\&P Global tool returns executive data — see Data Integrity Rule 10. Management names from training data will be stale.


\subparagraph{8. Integration Considerations}\mbox{}\par
\textbf{Required section — do not cut.} This is the analytical capstone of the Corp Dev tear sheet. Synthesized observations drawn from everything above. 3-4 substantive bullets:

\begin{itemize}[leftmargin=*,itemsep=2pt,topsep=2pt]
\item {} \textbf{Key risks:} Customer concentration, key-person dependency, technology debt indicators
\item {} \textbf{Integration complexity:} Geography (multi-country = harder), employee count, number of products
\item {} \textbf{Synergy signals:} Shared customers, complementary products, overlapping technology vendors
\item {} \textbf{Open questions:} What would a corp dev team need to diligence before making a bid? Focus on concrete, actionable diligence items — not open-ended strategic questions. Good examples: regulatory approval timeline and jurisdictions, change-of-control provisions in key contracts, key-person retention risk (especially founders/CTOs), customer contract transferability, IP ownership structure, outstanding litigation. Bad examples: ``How defensible is X?'' or ``What is the long-term market opportunity?'' — these are analyst questions, not diligence items.
\end{itemize}

Label clearly as analytical observations, not facts from the data source.


\subparagraph{Page Budget \& Cut Order}\mbox{}\par

\textbf{Target: 2 pages.} Corp Dev profiles can fill both pages with substantive content, but must not spill to a third unless the user explicitly requests a longer format. If content exceeds 2 pages, cut in this order (cut each completely before moving to the next):


\begin{enumerate}[leftmargin=*,itemsep=2pt,topsep=2pt]
\item {} Ownership Snapshot (section 7) — omit entirely
\item {} Customer Analysis prose under Revenue Mix (section 3) — reduce to 2 sentences max
\item {} Valuation Context prose (section 6) — keep peer comp table, cut interpretive paragraph below it
\item {} Precedent Transactions (section 6b) — keep the ``not available'' note, remove if space-critical
\item {} Capital Structure prose — keep table, cut interpretive paragraph below it
\item {} Strategic Fit Analysis (section 4) — compress each bullet to 2-3 sentences max (do not remove the section)
\end{enumerate}

Never cut: Business \& Product Overview (2), Financial Summary (5 core table), Integration Considerations (8).


\subparagraph{Formatting Notes}\mbox{}\par
\begin{itemize}[leftmargin=*,itemsep=2pt,topsep=2pt]
\item {} \textbf{Strategic Fit Analysis is the centerpiece.} Give it distinctive formatting:
\begin{itemize}[leftmargin=*,itemsep=2pt,topsep=2pt]
\item {} Bump text to 9.5pt (size: 19) — slightly larger than standard body text. This section should look like the most important content on the page.
\item {} Use the standard indented block-style bullets from the global style config (360 DXA left indent).
\item {} Do not apply left-border accents — they render inconsistently in docx-js.
\end{itemize}
\item {} Financial Summary should include R\&D and FCF conversion prominently.
\item {} Valuation Context with peer comps is required for public companies — do not cut before Ownership Snapshot.
\item {} For private targets, the profile leans heavily on qualitative sections (Business Overview, Strategic Fit, Relationships) — this is expected and still highly valuable.
\item {} Tone: analytical and evaluative, not promotional. This isn't a pitch — it's an assessment.
\item {} \textbf{Two-page discipline:} Before rendering, mentally estimate page count. The combination of Strategic Fit (3 long bullets) + Integration Considerations (4 bullets) alone can consume a full page — if so, compress Strategic Fit bullets first.
\end{itemize}

\vspace{0.8em}\hrule\vspace{0.8em}

\subsection{Skill Resource: spglobal/skills/tear-sheet/references/equity-research.md}\label{file:spglobal-skills-tear-sheet-references-equity-research.md}

\paragraph{Equity Research Tear Sheet}\mbox{}\par

\subparagraph{Purpose}\mbox{}\par
A dense snapshot for buy-side or sell-side analysts evaluating an investment. Emphasis on valuation, forward estimates, financial trajectory, and analyst consensus. Everything supports or challenges an investment thesis.


\textbf{Default page length:} 1 page. Equity research tear sheets are conventionally single-page — density is the point. If the user requests more space, extend to 2 pages.


\subparagraph{Query Plan}\mbox{}\par

Start with these queries. If results are incomplete, break into narrower follow-ups.


\textbf{Query 1 — Profile + market data:} ``[Company] company overview sector industry market capitalization enterprise value stock price 52-week high low shares outstanding beta'' → Header, Business Description → \textbf{Immediately write} to \emph{/tmp/tear-sheet/company-profile.txt}


\textbf{Query 2 — Historical financials:} ``[Company] annual income statement revenue gross profit EBITDA net income EPS free cash flow total debt cash and equivalents last 4 fiscal years'' → Financial Summary (pull 4 years; display 3; use the earliest year only for YoY growth computation) → \textbf{Immediately write} raw values to \emph{/tmp/tear-sheet/financials.csv}


\textbf{Query 2b — Segments (if available):} ``[Company] revenue by segment business unit last 2 fiscal years'' → Revenue \& Segment Breakdown (need 2 years for YoY growth). If segment data is unavailable, skip this section — do not leave a blank table. → \textbf{Immediately write} to \emph{/tmp/tear-sheet/segments.csv} (skip if no segment data returned)


\textbf{Query 3 — Valuation + consensus:} ``[Company] P/E EV/EBITDA EV/Revenue valuation multiples consensus revenue EPS estimates analyst recommendations price target'' → Valuation Snapshot, Consensus Estimates → \textbf{Immediately write} multiples to \emph{/tmp/tear-sheet/valuation.csv} → \textbf{Immediately write} estimates to \emph{/tmp/tear-sheet/consensus.csv}


\textbf{Query 4 — Earnings:} ``[Company] most recent earnings call key takeaways guidance'' → Earnings Highlights → \textbf{Immediately write} to \emph{/tmp/tear-sheet/earnings.txt}


\textbf{Query 5 — Stock performance:} ``[Company] stock price return 1 month 3 month 6 month 1 year YTD performance'' → Stock Performance (no intermediate file — data goes directly into document)


\textbf{Query 6 (if user provided comps):} ``[Comp company] market cap EV/Revenue EV/EBITDA revenue growth'' — repeat for each comp → Peer context for Valuation Snapshot → \textbf{Immediately write} to \emph{/tmp/tear-sheet/peer-comps.csv}


\subparagraph{Sections}\mbox{}\par

Listed in priority order. If constrained to one page, cut from the bottom.


\subparagraph{1. Company Header}\mbox{}\par
Compact key-value block rendered as a two-column borderless table per the global style config.


Left column: Ticker (exchange), sector / industry, HQ Right column: Stock price, 52-week range, market cap, EV, shares outstanding, beta


\subparagraph{2. Business Description}\mbox{}\par
2-3 tight sentences. \textbf{Rewrite in your own words for an analyst audience} — do not paste the CIQ company summary verbatim. The CIQ summary is an input; the output should be concise and thesis-oriented.


For well-known, widely covered companies, assume the analyst already knows the basic business. Lead with what's \emph{changing} — strategic pivots, portfolio reshaping, new growth vectors — not a Wikipedia-style overview. For example, for a Fortune 500 company, don't spend a sentence on ``provides financial data to institutions.'' Instead: ``Reshaping its portfolio toward higher-margin data and analytics businesses, with a pending Mobility spin-off and aggressive AI investment through its Kensho subsidiary.''


For lesser-known companies, a brief ``what they do'' sentence is appropriate before pivoting to the thesis-relevant dynamics.


\subparagraph{3. Valuation Snapshot}\mbox{}\par
Centerpiece of the equity tear sheet.


\begingroup
\small
\setlength{\tabcolsep}{2pt}
\begin{longtable}{>{\RaggedRight\arraybackslash\hspace{0pt}}p{0.320\textwidth}>{\RaggedRight\arraybackslash\hspace{0pt}}p{0.320\textwidth}>{\RaggedRight\arraybackslash\hspace{0pt}}p{0.320\textwidth}}
\toprule
{} \textbf{Metric} & \textbf{Trailing} & \textbf{Forward (NTM)} \\
\midrule
{} P/ E &  &  \\
{} EV/ EBITDA &  &  \\
{} EV/ Revenue &  &  \\
{} P/ FCF &  &  \\
{} Dividend Yield &  &  \\
\bottomrule
\end{longtable}
\endgroup

\textbf{Forward multiples are mandatory when available from the tools.} Showing only trailing multiples when forward data exists is a significant gap — analysts value companies on forward earnings. If forward multiples aren't available, show trailing only and note ``Fwd estimates N/A.''


If the user provided comparable companies, add columns for each comp (or a ``Peer Median'' column if 3+ comps). If no user comps, and the tools return peer data, include a Peer Median column. If neither, show the company's multiples only.


\subparagraph{4. Consensus Estimates}\mbox{}\par
What the Street expects — a key differentiator for this tear sheet type.


\textbf{Data availability note:} The depth of consensus data varies by company and coverage. Include whatever the tools return — analyst count, price target, Buy/Hold/Sell distribution — but do not fabricate or estimate any consensus figure. For thinly covered companies, this section may contain only revenue and EPS estimates, which is still valuable.


\begingroup
\small
\setlength{\tabcolsep}{2pt}
\begin{longtable}{>{\RaggedRight\arraybackslash\hspace{0pt}}p{0.320\textwidth}>{\RaggedRight\arraybackslash\hspace{0pt}}p{0.320\textwidth}>{\RaggedRight\arraybackslash\hspace{0pt}}p{0.320\textwidth}}
\toprule
{} \textbf{Metric} & \textbf{FY[year] Est.} & \textbf{FY[year+1] Est.} \\
\midrule
{} Revenue &  &  \\
{} EPS (normalized) &  &  \\
{} EBITDA &  &  \\
\bottomrule
\end{longtable}
\endgroup

Use normalized/adjusted EPS where available — GAAP EPS can be distorted by one-time items and is less useful for forward estimates.


Below the main estimates table, include a compact analyst consensus block (if data is available from the tools):


\begingroup
\small
\setlength{\tabcolsep}{2pt}
\begin{longtable}{>{\RaggedRight\arraybackslash\hspace{0pt}}p{0.480\textwidth}>{\RaggedRight\arraybackslash\hspace{0pt}}p{0.480\textwidth}}
\toprule
{} \textbf{Analyst Consensus} & \textbf{} \\
\midrule
{} Mean Price Target & \$XXX \\
{} \# of Estimates & XX \\
{} Buy /  Hold /  Sell & XX /  XX /  XX \\
\bottomrule
\end{longtable}
\endgroup

If price target or recommendation data is unavailable, omit the analyst consensus block rather than showing empty rows. Do not fabricate these figures.


If consensus data isn't available at all, skip this entire section and note ``Consensus estimates not available.'' Do not estimate.


\subparagraph{5. Financial Summary (3-Year)}\mbox{}\par
Dense single table using actual fiscal year labels.


\begingroup
\small
\setlength{\tabcolsep}{2pt}
\begin{longtable}{>{\RaggedRight\arraybackslash\hspace{0pt}}p{0.240\textwidth}>{\RaggedRight\arraybackslash\hspace{0pt}}p{0.240\textwidth}>{\RaggedRight\arraybackslash\hspace{0pt}}p{0.240\textwidth}>{\RaggedRight\arraybackslash\hspace{0pt}}p{0.240\textwidth}}
\toprule
{} \textbf{Metric (\$M)} & \textbf{FY20XX} & \textbf{FY20XX} & \textbf{FY20XX} \\
\midrule
{} Revenue &  &  &  \\
{} Revenue Growth \% &  &  &  \\
{} Gross Margin \% &  &  &  \\
{} EBITDA &  &  &  \\
{} EBITDA Margin \% &  &  &  \\
{} Net Income &  &  &  \\
{} EPS (Diluted) &  &  &  \\
{} Free Cash Flow &  &  &  \\
{} Net Debt &  &  &  \\
\bottomrule
\end{longtable}
\endgroup

Compute derived metrics (growth \%, margins) from raw data rather than querying separately.


\textbf{Capital Structure (separate compact sub-table below the Financial Summary):} Do NOT append these rows to the 3-year Financial Summary table — blank cells in prior-year columns look like missing data. Instead, render as a separate 2-column sub-table (Metric | Value) directly below, showing only the most recent fiscal year:


\begingroup
\small
\setlength{\tabcolsep}{2pt}
\begin{longtable}{>{\RaggedRight\arraybackslash\hspace{0pt}}p{0.480\textwidth}>{\RaggedRight\arraybackslash\hspace{0pt}}p{0.480\textwidth}}
\toprule
{} \textbf{Metric} & \textbf{FY[latest]} \\
\midrule
{} Total Debt & \$XXM \\
{} Cash \& Equivalents & \$XXM \\
{} Net Debt & \$XXM \\
\bottomrule
\end{longtable}
\endgroup

This matches the capital structure format used in Corp Dev and IB/M\&A tear sheets. Keep it tight — no extra spacing between the Financial Summary and this sub-table.


\subparagraph{6. Revenue \& Segment Breakdown}\mbox{}\par
If segment data is available from the tools, include a compact segment revenue table. This is essential context for equity research — analysts need to see which business lines are driving growth.


\begingroup
\small
\setlength{\tabcolsep}{2pt}
\begin{longtable}{>{\RaggedRight\arraybackslash\hspace{0pt}}p{0.240\textwidth}>{\RaggedRight\arraybackslash\hspace{0pt}}p{0.240\textwidth}>{\RaggedRight\arraybackslash\hspace{0pt}}p{0.240\textwidth}>{\RaggedRight\arraybackslash\hspace{0pt}}p{0.240\textwidth}}
\toprule
{} \textbf{Segment} & \textbf{Revenue (\$M)} & \textbf{\% of Total} & \textbf{YoY Growth} \\
\midrule
{} [Segment A] &  &  &  \\
{} [Segment B] &  &  &  \\
\bottomrule
\end{longtable}
\endgroup

Keep this compact for the one-pager: table only, no qualitative paragraph. If segment data is unavailable, skip this section entirely — do not include a blank or placeholder table.


\subparagraph{7. Recent Earnings Highlights}\mbox{}\par
Include the quarter and date (e.g., ``Q3 FY2025 — October 2025''). 3-4 bullets with an \textbf{investment lens} — this audience wants segment-level specifics, not strategic themes:


\textbf{Lead with guidance if provided.} If management issued forward guidance, make it the first bullet with a bold \textbf{Guidance:} prefix (e.g., ``\textbf{Guidance:} FY2026 EPS of \$19.40–\$19.65, implying 9–10\% growth; organic revenue growth of 6–8\%.''). Guidance is the single most actionable forward-looking data point for an analyst — do not bury it among operating highlights.


\textbf{Beat/miss context (data permitting):} If consensus estimate data is available for the reported period, frame the headline result bullet as a beat/miss: ``Revenue of \$X beat/missed consensus of \$Y by Z\%.'' If consensus data for the reported period is not available from the tools, describe results in absolute terms and growth rates only — do not fabricate consensus figures.


Then:

\begin{itemize}[leftmargin=*,itemsep=2pt,topsep=2pt]
\item {} Segment-level performance: which business lines accelerated or decelerated
\item {} Margin trajectory: any commentary on cost structure or profitability trends
\end{itemize}

Attribute to management: ``Management highlighted…'' or ``CFO noted…'' rather than declarative claims. This section should feel like an earnings recap, not a press release summary.


\subparagraph{8. Key Operating Metrics}\mbox{}\par
2-3 sector-relevant KPIs if available from the data (subscriber count, same-store sales, NIM, etc.). If nothing sector-specific, compute ROE and ROIC from the financials already retrieved. Optional — drop if space is tight.


\subparagraph{9. Stock Performance (cut first)}\mbox{}\par
Period return table. Low analytical value for this audience — cut this before anything else.


\begingroup
\small
\setlength{\tabcolsep}{2pt}
\begin{longtable}{>{\RaggedRight\arraybackslash\hspace{0pt}}p{0.480\textwidth}>{\RaggedRight\arraybackslash\hspace{0pt}}p{0.480\textwidth}}
\toprule
{} \textbf{Period} & \textbf{Return} \\
\midrule
{} 1 Month &  \\
{} 3 Month &  \\
{} YTD &  \\
{} 1 Year &  \\
\bottomrule
\end{longtable}
\endgroup

\subparagraph{Formatting Notes}\mbox{}\par
\begin{itemize}[leftmargin=*,itemsep=2pt,topsep=2pt]
\item {} \textbf{This is the densest tear sheet type.} It should feel noticeably more compressed than the other templates.
\item {} Financial Summary table: 8pt text (size: 16), row height 240 DXA. Tightest spacing of any template.
\item {} Valuation Snapshot and Consensus Estimates: 8.5pt text, occupy the visual center — analysts scan these first.
\item {} Body text: 8.5pt (size: 17) — smaller than the global 9pt default. Every half-point matters on a one-pager.
\item {} Section spacing: minimize. 6pt before section headers (not the global 12pt), 2pt after rules.
\item {} Bold any metric showing a notable inflection (growth turning positive, margin expansion/compression).
\item {} Prioritize information over whitespace at every turn.
\end{itemize}

\vspace{0.8em}\hrule\vspace{0.8em}

\subsection{Skill Resource: spglobal/skills/tear-sheet/references/ib-ma.md}\label{file:spglobal-skills-tear-sheet-references-ib-ma.md}

\paragraph{Investment Banking / M\&A Tear Sheet}\mbox{}\par

\subparagraph{Purpose}\mbox{}\par
A company profile for transaction context — potential target, acquirer, or pitchbook inclusion. Emphasis on strategic positioning, segment detail, business relationships, and M\&A activity. Valuation framed in transaction multiples, not investment thesis.


\textbf{Default page length:} 1-2 pages. IB company profiles routinely spill to a second page, especially when segment and M\&A data is rich.


\subparagraph{Query Plan}\mbox{}\par

Start with these queries. Break apart and follow up if results are incomplete.


\textbf{Query 1 — Profile + identification:} ``[Company] business description headquarters founded employees sector industry ownership structure major shareholders'' → Header, Business Overview → \textbf{Immediately write} to \emph{/tmp/tear-sheet/company-profile.txt}


\textbf{Query 2 — Segments + financials:} ``[Company] revenue by segment business unit last 2 fiscal years AND annual income statement revenue gross profit EBITDA operating income net income capex free cash flow total debt cash and equivalents last 4 fiscal years'' → Segment Breakdown (need 2 years to compute YoY growth), Financial Summary (pull 4 years; display 3; use the earliest year only for YoY growth computation) → \textbf{Immediately write} raw financials to \emph{/tmp/tear-sheet/financials.csv} → \textbf{Immediately write} segment data to \emph{/tmp/tear-sheet/segments.csv} (skip if no segment data returned)


\textbf{Query 3 — Valuation + comps:} If the user provided specific comparable companies, query each: ``[Comp company] enterprise value EV/Revenue EV/EBITDA revenue growth EBITDA margin'' Otherwise: Use the competitors tool to identify 3-5 public peers, then pull EV/Revenue (NTM), EV/EBITDA (NTM), revenue growth \%, and EBITDA margin \% for each. Also: ``[Company] enterprise value market capitalization EV/Revenue EV/EBITDA valuation multiples'' → Header (EV/market cap), Trading Comps → \textbf{Immediately write} company multiples to \emph{/tmp/tear-sheet/valuation.csv} → \textbf{Immediately write} peer data to \emph{/tmp/tear-sheet/peer-comps.csv}


\textbf{Query 4 — M\&A activity:} ``[Company] acquisitions divestitures completed transactions last 5 years deal value'' → Company's own deals → \textbf{Immediately write} to \emph{/tmp/tear-sheet/ma-activity.csv}


\textbf{Query 5 — Comparable transactions (data permitting):} Use the specific industry from Query 1 results (e.g., ``cloud observability software M\&A transactions'' not ``technology M\&A''). If the user specified comps, also try: ``[comp company] acquisition deal multiples.'' → Comparable Transactions → \textbf{Append} comparable transactions to \emph{/tmp/tear-sheet/ma-activity.csv} (add \emph{type=precedent} to distinguish from company's own deals)


\textbf{Query 6 — Relationships + ownership:} ``[Company] key customers suppliers partners business relationships institutional ownership insider ownership'' → Business Relationships, Ownership (data permitting) → \textbf{Immediately write} to \emph{/tmp/tear-sheet/relationships.txt}


\subparagraph{Sections}\mbox{}\par

Listed in priority order. See ``Page Budget \& Cut Order'' below for explicit cut guidance.


\subparagraph{1. Company Header}\mbox{}\par
Key-value block with transaction-relevant identifiers.


\begingroup
\small
\setlength{\tabcolsep}{2pt}
\begin{longtable}{>{\RaggedRight\arraybackslash\hspace{0pt}}p{0.480\textwidth}>{\RaggedRight\arraybackslash\hspace{0pt}}p{0.480\textwidth}}
\toprule
{} \textbf{Field} & \textbf{Notes} \\
\midrule
{} Company name & Large, bold \\
{} Ticker \& Exchange & If public; omit for private \\
{} HQ Location & City, State/ Country \\
{} Founded & Year \\
{} Employees & Approximate headcount \\
{} Market Cap /  Valuation & Market cap if public; latest known valuation if private \\
{} Enterprise Value &  \\
{} Ownership & Public /  PE-backed (name sponsor) /  Family /  Other \\
\bottomrule
\end{longtable}
\endgroup

For private companies: note ``Private Company'' prominently; skip ticker, exchange.


\subparagraph{2. Business Overview}\mbox{}\par
\textbf{This is a pitchbook paragraph, not a CIQ summary.} Do not paste the company description from the data tools. Rewrite in 4-6 sentences using pitchbook prose:

\begin{itemize}[leftmargin=*,itemsep=2pt,topsep=2pt]
\item {} Characterize the revenue model (recurring vs. transactional, subscription vs. license)
\item {} Name the customer base at scale (``serves 80\% of Fortune 500 banks'')
\item {} Identify competitive moats by name (``proprietary dataset of X'', ``only provider with Y'')
\item {} Frame growth vectors that would matter to a buyer or seller
\end{itemize}

The tone should be confident and specific. A banker reading this paragraph should immediately understand the company's strategic positioning.


\subparagraph{3. Revenue \& Segment Breakdown}\mbox{}\par
Revenue by business unit or product line.


\begingroup
\small
\setlength{\tabcolsep}{2pt}
\begin{longtable}{>{\RaggedRight\arraybackslash\hspace{0pt}}p{0.240\textwidth}>{\RaggedRight\arraybackslash\hspace{0pt}}p{0.240\textwidth}>{\RaggedRight\arraybackslash\hspace{0pt}}p{0.240\textwidth}>{\RaggedRight\arraybackslash\hspace{0pt}}p{0.240\textwidth}}
\toprule
{} \textbf{Segment} & \textbf{Revenue (\$M)} & \textbf{\% of Total} & \textbf{YoY Growth} \\
\midrule
{} [Segment A] &  &  &  \\
{} [Segment B] &  &  &  \\
{} \textbf{Total} &  & \textbf{100\%} &  \\
\bottomrule
\end{longtable}
\endgroup

\textbf{If segment data isn't available:} Replace the table with 2-3 sentences describing the revenue mix qualitatively (e.g., ``Revenue is primarily derived from subscription licenses (\textasciitilde{}70\%) and professional services (\textasciitilde{}30\%).''). A blank table looks broken — a qualitative description still conveys useful information.


\subparagraph{4. Financial Summary (3-Year + LTM)}\mbox{}\par
Actual fiscal years as column headers. \textbf{If LTM (last twelve months) data is available, add an LTM column to the right of the most recent fiscal year.} Label it ``LTM [quarter end date]'' (e.g., ``LTM Q3 2025''). Bankers work off LTM for valuation cross-checks — a fiscal-year-only view can be 6+ months stale.


\begingroup
\small
\setlength{\tabcolsep}{2pt}
\begin{longtable}{>{\RaggedRight\arraybackslash\hspace{0pt}}p{0.184\textwidth}>{\RaggedRight\arraybackslash\hspace{0pt}}p{0.184\textwidth}>{\RaggedRight\arraybackslash\hspace{0pt}}p{0.184\textwidth}>{\RaggedRight\arraybackslash\hspace{0pt}}p{0.184\textwidth}>{\RaggedRight\arraybackslash\hspace{0pt}}p{0.184\textwidth}}
\toprule
{} \textbf{Metric (\$M)} & \textbf{FY20XX} & \textbf{FY20XX} & \textbf{FY20XX} & \textbf{LTM [date]} \\
\midrule
{} Revenue &  &  &  &  \\
{} Revenue Growth \% &  &  &  &  \\
{} Gross Profit &  &  &  &  \\
{} Gross Margin \% &  &  &  &  \\
{} EBITDA &  &  &  &  \\
{} EBITDA Margin \% &  &  &  &  \\
{} Operating Income &  &  &  &  \\
{} Net Income &  &  &  &  \\
{} Capex &  &  &  &  \\
{} Free Cash Flow &  &  &  &  \\
{} FCF Margin \% &  &  &  &  \\
\bottomrule
\end{longtable}
\endgroup

In M\&A context, margins and FCF conversion matter as much as absolute scale — a buyer is evaluating efficiency and cash generation potential.


\textbf{Capital Structure} (compact sub-table below the financial summary):


\begingroup
\small
\setlength{\tabcolsep}{2pt}
\begin{longtable}{>{\RaggedRight\arraybackslash\hspace{0pt}}p{0.480\textwidth}>{\RaggedRight\arraybackslash\hspace{0pt}}p{0.480\textwidth}}
\toprule
{} \textbf{Metric} & \textbf{Value} \\
\midrule
{} Total Debt & \$XXM \\
{} Cash \& Equivalents & \$XXM \\
{} Net Debt & \$XXM \\
{} Net Debt /  EBITDA & X.Xx \\
{} S\&P Credit Rating & XX (if available) \\
\bottomrule
\end{longtable}
\endgroup

Pull from the balance sheet data already retrieved. This is standard in any banker's company profile — total debt, leverage ratio, and credit rating tell a buyer what the financing picture looks like. If credit rating is unavailable, omit that row.


\subparagraph{5. Trading Comps}\mbox{}\par
Company multiples \textbf{and operating metrics} — bankers use these together, not multiples in isolation.


\textbf{Forward multiples are mandatory when available.} The table must include NTM multiples if the tools return them. Showing only trailing multiples when forward data exists is a significant gap — bankers price deals on forward earnings.


\textbf{Peer columns are expected.} If the user provided comps, each gets its own column. If not, use the competitors tool to identify 3-5 public peers, pull their multiples and operating metrics, and include them. A company-only trading comps table without peer context is incomplete for any banking use case.


\begingroup
\small
\setlength{\tabcolsep}{2pt}
\begin{longtable}{>{\RaggedRight\arraybackslash\hspace{0pt}}p{0.153\textwidth}>{\RaggedRight\arraybackslash\hspace{0pt}}p{0.153\textwidth}>{\RaggedRight\arraybackslash\hspace{0pt}}p{0.153\textwidth}>{\RaggedRight\arraybackslash\hspace{0pt}}p{0.153\textwidth}>{\RaggedRight\arraybackslash\hspace{0pt}}p{0.153\textwidth}>{\RaggedRight\arraybackslash\hspace{0pt}}p{0.153\textwidth}}
\toprule
{} \textbf{Metric} & \textbf{[Company]} & \textbf{[Comp 1]} & \textbf{[Comp 2]} & \textbf{[Comp 3]} & \textbf{Peer Median} \\
\midrule
{} EV/ Revenue (NTM) &  &  &  &  &  \\
{} EV/ EBITDA (NTM) &  &  &  &  &  \\
{} Revenue Growth \% &  &  &  &  &  \\
{} EBITDA Margin \% &  &  &  &  &  \\
\bottomrule
\end{longtable}
\endgroup

If only the company's own multiples are available after attempting peer queries, show those and list selected peer names below the table.


For private companies, skip this section.


\subparagraph{6. M\&A Activity}\mbox{}\par
Two parts. \textbf{Deal values are expected in this section} — bankers notice when they're missing.


\textbf{Hard rule: If the M\&A tools return a transaction value, it must appear in the output.} Never downgrade a known deal value to ``Undisclosed'' — this is a data integrity violation. Use ``Undisclosed'' only when the tools genuinely return no value for a transaction.


\textbf{a) Company's Own Transactions:}


\textbf{The table must include a dedicated Deal Value column.} Do not merge deal values into a Notes or Rationale column — bankers scan for numbers, and a dedicated column enables that.


\begingroup
\small
\setlength{\tabcolsep}{2pt}
\begin{longtable}{>{\RaggedRight\arraybackslash\hspace{0pt}}p{0.184\textwidth}>{\RaggedRight\arraybackslash\hspace{0pt}}p{0.184\textwidth}>{\RaggedRight\arraybackslash\hspace{0pt}}p{0.184\textwidth}>{\RaggedRight\arraybackslash\hspace{0pt}}p{0.184\textwidth}>{\RaggedRight\arraybackslash\hspace{0pt}}p{0.184\textwidth}}
\toprule
{} \textbf{Date} & \textbf{Target /  Divested} & \textbf{Deal Value (\$M)} & \textbf{Type} & \textbf{Rationale} \\
\midrule
{}  &  &  & Acquisition /  Divestiture &  \\
\bottomrule
\end{longtable}
\endgroup

If none found: ``No disclosed transactions in the last 5 years.''


\textbf{b) Comparable Sector Transactions (data permitting):} Use the specific industry identified earlier (e.g., ``cloud infrastructure software M\&A'' not ``technology M\&A''). If the user specified comps, also look for transactions involving those specific companies.


\begingroup
\small
\setlength{\tabcolsep}{2pt}
\begin{longtable}{>{\RaggedRight\arraybackslash\hspace{0pt}}p{0.153\textwidth}>{\RaggedRight\arraybackslash\hspace{0pt}}p{0.153\textwidth}>{\RaggedRight\arraybackslash\hspace{0pt}}p{0.153\textwidth}>{\RaggedRight\arraybackslash\hspace{0pt}}p{0.153\textwidth}>{\RaggedRight\arraybackslash\hspace{0pt}}p{0.153\textwidth}>{\RaggedRight\arraybackslash\hspace{0pt}}p{0.153\textwidth}}
\toprule
{} \textbf{Date} & \textbf{Target} & \textbf{Acquirer} & \textbf{EV (\$M)} & \textbf{EV/ Rev} & \textbf{EV/ EBITDA} \\
\midrule
\bottomrule
\end{longtable}
\endgroup

If not available from the tools: ``Comparable transaction data not available from data source.'' Do not fabricate precedent multiples.


\subparagraph{7. Key Business Relationships}\mbox{}\par
Leverages S\&P Capital IQ relationship data. Group by type:


\textbf{Customers:} Top 3-5 named clients \textbf{Suppliers:} Key vendors and technology providers \textbf{Partners:} Strategic alliances, distribution, JVs \textbf{Competitors:} Key competitive names


For each entry, include the relationship type and a parenthetical descriptor explaining the nature of the relationship (e.g., ``AWS — cloud infrastructure (primary hosting provider)'', ``JPMorgan — customer (enterprise analytics platform)''). A name alone without context isn't useful. Give this section breathing room — it's strategically valuable and most competing tear sheets lack it.


\subparagraph{8. Ownership Snapshot (data permitting — omit if tools return nothing)}\mbox{}\par
If the tools return institutional ownership, insider ownership percentage, or top holders, include a compact block:


Top institutional holders (if available): list top 3-5 by ownership \%. Insider ownership: aggregate percentage.


For PE-backed companies, name the sponsor, acquisition year, and ownership stake. Ownership structure directly affects deal complexity — a widely held public company, a founder-controlled company, and a PE-backed company each require fundamentally different deal approaches.


If ownership data is not returned by the tools, omit this section entirely rather than showing placeholders.


\textbf{Do not include a Management Team table.} No S\&P Global tool returns executive data — see Data Integrity Rule 10. Management names from training data will be stale.


\subparagraph{Page Budget \& Cut Order}\mbox{}\par

\textbf{Target: 2 pages maximum.} IB profiles can use both pages fully, but must not spill to a third. If content exceeds 2 pages, cut in this order (cut each completely before moving to the next):


\begin{enumerate}[leftmargin=*,itemsep=2pt,topsep=2pt]
\item {} Ownership Snapshot (section 8) — omit entirely
\item {} Comparable Sector Transactions (section 6b) — keep company's own M\&A, cut precedents
\item {} Business Relationships (section 7) — compress to top 3 competitors + top 2 each of suppliers/partners/customers, drop descriptors to parenthetical only
\item {} Capital Structure (section 4) — merge Net Debt/EBITDA into Financial Summary as a single row
\item {} Trading Comparables (section 5) — reduce to 3 peers
\end{enumerate}

Never cut: Business Overview (2), Revenue \& Segment Breakdown (3), Financial Summary (4 core), M\&A Activity company transactions (6a).


\subparagraph{Formatting Notes}\mbox{}\par
\begin{itemize}[leftmargin=*,itemsep=2pt,topsep=2pt]
\item {} Business Overview and Segment Breakdown get the most visual weight.
\item {} \textbf{M\&A Activity is the signature section:} Use Accent color (\#2E75B6) fill with white bold text for M\&A table headers — this is the only table that uses a different header color from the standard \#D6E4F0. It should catch the eye on a scan.
\item {} \textbf{Business Relationships: give breathing room when space permits.} 6pt between each relationship group, full descriptors on each entry. But if hitting page 3, compress per the cut order above — this section is high-value but also the most compressible.
\item {} For private companies, the header, business overview, financials, and relationships form the core; other sections may be sparse.
\item {} \textbf{Two-page discipline:} Before rendering, mentally estimate page count. If Business Relationships + M\&A Activity together will exceed one full page, start cutting per the cut order.
\end{itemize}

\vspace{0.8em}\hrule\vspace{0.8em}

\subsection{Skill Resource: spglobal/skills/tear-sheet/references/sales-bd.md}\label{file:spglobal-skills-tear-sheet-references-sales-bd.md}

\paragraph{Sales / Business Development Tear Sheet}\mbox{}\par

\subparagraph{Purpose}\mbox{}\par
A briefing document to prepare a sales or BD team for a client meeting. Emphasis on understanding the prospect's business, what leadership is publicly prioritizing, and identifying angles to connect your product/service to their needs. Financials are simplified to scale and trajectory.


This is the most narrative-driven tear sheet type. It reads like a briefing memo, not a data sheet.


\textbf{Default page length:} 1-2 pages. No strict convention — prioritize readability over compression.


\subparagraph{Query Plan}\mbox{}\par

Start with these queries. Follow up if results are incomplete.


\textbf{Query 1 — Company profile + financials:} ``[Company] business description overview products services headquarters employees sector industry revenue gross margin last 2 fiscal years'' → Header, Company Overview, Financial Snapshot → \textbf{Immediately write} to \emph{/tmp/tear-sheet/company-profile.txt} → \textbf{Immediately write} raw financials to \emph{/tmp/tear-sheet/financials.csv}


\textbf{Query 2 — Strategy + earnings:} ``[Company] most recent earnings call strategic priorities CEO commentary guidance key initiatives'' → Strategic Priorities → \textbf{Immediately write} to \emph{/tmp/tear-sheet/earnings.txt}


\textbf{Query 3 — Relationships + news:} ``[Company] key customers suppliers partners competitors business relationships recent acquisitions partnerships announcements'' → Key Relationships, Recent News → \textbf{Immediately write} to \emph{/tmp/tear-sheet/relationships.txt}


Three queries is usually sufficient. For well-known companies, these return rich results. For smaller or private companies, Queries 2 and 3 may be sparse — that's fine. The qualitative sections still deliver value because Claude synthesizes whatever is available into useful framing.


\subparagraph{Sections}\mbox{}\par

Listed in priority order. Cut from the bottom if space is constrained.


\subparagraph{1. Company Header}\mbox{}\par
Streamlined — identification only, no valuation metrics.


\begingroup
\small
\setlength{\tabcolsep}{2pt}
\begin{longtable}{>{\RaggedRight\arraybackslash\hspace{0pt}}p{0.480\textwidth}>{\RaggedRight\arraybackslash\hspace{0pt}}p{0.480\textwidth}}
\toprule
{} \textbf{Field} & \textbf{Notes} \\
\midrule
{} Company name & Large, bold \\
{} Ticker \& Exchange & If public; omit for private \\
{} HQ & City, State/ Country \\
{} Industry /  Sector &  \\
{} Employees &  \\
{} Annual Revenue (latest) & Round to nearest \$B or \$M \\
{} Market Cap & If public; omit for private \\
{} FY Revenue Guidance & If available from earnings data; e.g., ``6–8\% organic growth'' \\
{} FY EPS or EPS Guidance & If public company; pull from \emph{earnings.txt} \\
{} Key Financial Metric & Adj. EPS, EBITDA margin, or similar — one headline number \\
\bottomrule
\end{longtable}
\endgroup

Forward guidance and EPS data come from Query 2 (earnings) — source from \emph{earnings.txt} when populating the header, not just \emph{company-profile.txt}.


One horizontal band at the top. No enterprise value, no beta — this audience doesn't need them.


\subparagraph{2. Company Overview}\mbox{}\par
The most important section. A sales rep should read this and walk into a meeting with a solid grasp of the prospect.


\textbf{Do not paste the CIQ company summary.} It reads like an SEC filing and will lose a non-finance audience. Rewrite in 3-5 sentences of plain language:

\begin{itemize}[leftmargin=*,itemsep=2pt,topsep=2pt]
\item {} What the company does (no investor jargon — say ``sells'' not ``monetizes'', say ``yearly revenue'' not ``top-line CAGR'')
\item {} Who their customers are
\item {} How they make money
\item {} Where they operate
\item {} What makes them distinctive
\end{itemize}

Write in prose paragraphs. A non-finance person should understand every sentence.


\subparagraph{3. Strategic Priorities \& Management Commentary}\mbox{}\par
What is leadership focused on right now? If you know what the CEO is talking about, you can align your pitch.


\textbf{This is NOT an earnings recap.} Do not reuse the same bullets as the equity research tear sheet. Extract 3-5 \emph{strategic themes} with thematic labels and supporting data points:


\begin{itemize}[leftmargin=*,itemsep=2pt,topsep=2pt]
\item {} \textbf{AI/Data Monetization:} Management is investing aggressively in AI capabilities, allocating \$X to R\&D in FY2024 and launching three new AI-powered products.
\item {} \textbf{Margin Expansion:} CEO cited operational efficiency as a top priority, targeting 200bps of operating margin improvement through automation.
\item {} \textbf{International Growth:} Opening new offices in EMEA; Europe now represents 25\% of revenue, up from 18\%.
\end{itemize}

Each bullet should answer: ``What does this company care about right now, and what evidence supports it?'' A sales rep should read these and immediately see angles to connect their product.


If earnings call data isn't available (common for private companies), replace with ``Recent Developments'' pulling from whatever news or announcements the tools return.


\subparagraph{4. Financial Snapshot}\mbox{}\par
Simplified. Single compact table — just scale, direction, and one profitability signal. A sales rep needs to know three things: how big is this company, is it growing, and is it healthy.


\begingroup
\small
\setlength{\tabcolsep}{2pt}
\begin{longtable}{>{\RaggedRight\arraybackslash\hspace{0pt}}p{0.240\textwidth}>{\RaggedRight\arraybackslash\hspace{0pt}}p{0.240\textwidth}>{\RaggedRight\arraybackslash\hspace{0pt}}p{0.240\textwidth}>{\RaggedRight\arraybackslash\hspace{0pt}}p{0.240\textwidth}}
\toprule
{} \textbf{Metric} & \textbf{FY[prior year]} & \textbf{FY[latest]} & \textbf{YoY Change} \\
\midrule
{} Revenue &  &  & +X\% \\
{} Gross Margin \% &  &  & +X.Xpp \\
{} Employees &  &  &  \\
\bottomrule
\end{longtable}
\endgroup

Three rows maximum. No EBITDA, no EV multiples, no balance sheet, no net income. Gross margin (not EBITDA, not net income) is the right single profitability metric for a sales audience — it's intuitive (``how much they keep after direct costs'') and signals business health without requiring finance literacy. If the rep needs deeper financials, they can request the equity research version.


For private companies, this table may have only revenue and employees — two data points showing scale and growth are enough for a sales meeting.


\subparagraph{5. Key Relationships \& Ecosystem}\mbox{}\par
Helps reps understand the prospect's world and identify angles.


Four labeled lists, 3-5 names each with one-line descriptors:


\textbf{Customers:} Who they sell to. Helps understand their go-to-market and shared target customers. Example: ``JPMorgan Chase — customer (enterprise analytics platform deployment)''


\textbf{Vendors / Technology:} What they buy. Identifies displacement or integration opportunities. Example: ``AWS — vendor (primary cloud infrastructure); Snowflake — vendor (data warehouse)''


\textbf{Partners:} Who they work with. Co-sell angles, mutual connections. Example: ``Deloitte — partner (implementation for enterprise clients)''


\textbf{Competitors:} Who they compete against. Positioning context. Example: ``Bloomberg — competitor (financial data terminal)''


Every entry needs a parenthetical descriptor so the rep understands \emph{what} the relationship means, not just \emph{who}.


If relationship data isn't available, note ``Relationship data not available'' and move on.


\subparagraph{6. Recent News \& Developments}\mbox{}\par
3-5 bullets, each one sentence, most recent first. Include date if available. \textbf{Lead with the latest earnings report} — it's the most important recent event for a sales meeting and sets the context for everything else. Then: M\&A, product launches, leadership changes, major wins.


\textbf{Include upcoming catalysts at the end of the section} with a ``\textbf{Coming Up:}'' label. Examples: upcoming earnings dates, pending divestitures, product launches, investor days, regulatory milestones. These create natural engagement windows for sales outreach and are often the best reason to reach out (``I noticed you have your Investor Day coming up in Q2...''). Pull from earnings call commentary and M\&A data (pending transactions).


\subparagraph{7. Conversation Starters (cut first)}\mbox{}\par
2-3 suggested talking points synthesized from the data above. \textbf{Each starter must reference a specific strategic priority from Section 3} — generic questions are useless. Frame as questions.


\textbf{Tag each starter with a suggested persona} — e.g., (CTO), (CFO), (CDO), (Head of Data), (VP Engineering), (Head of Procurement). A rep meeting with a CTO needs a different opener than one meeting with a CFO. If the user specified who they're meeting with, generate starters for that persona only.


Examples:

\begin{itemize}[leftmargin=*,itemsep=2pt,topsep=2pt]
\item {} \emph{(CTO/CDO)} ``You mentioned [specific AI initiative from Section 3] on your last earnings call — how is that impacting your [relevant need]?''
\item {} \emph{(CFO/COO)} ``With your [margin expansion target from Section 3], are you evaluating tools that could accelerate that timeline?''
\item {} \emph{(VP Strategy/Corp Dev)} ``I noticed [Company] recently acquired [target] — does that change your approach to [capability]?''
\end{itemize}

If the user described what product/service they're selling, tailor these to connect the prospect's priorities to that offering. If not, keep them grounded in the prospect's strategic themes but generic in solution framing.


\textbf{Validation test:} Before finalizing each conversation starter, verify it contains at least one specific number, date, product name, or initiative name drawn from the tear sheet. A question that could apply to any company in the sector is not a good conversation starter. ``How are you thinking about AI?'' is generic. ``Your CEO mentioned \$1B in AI investment since 2018 and 20\% adoption of the iLEVEL auto-ingestion feature — how is that changing your analysts' workflows?'' is specific and demonstrates preparation.


This section is Claude's synthesis, not raw data from the tools. Label it clearly.


\subparagraph{Formatting Notes}\mbox{}\par
\begin{itemize}[leftmargin=*,itemsep=2pt,topsep=2pt]
\item {} \textbf{This should be the warmest, most readable tear sheet.} Rigorous content, but the formatting should feel like a briefing memo, not a data terminal.
\item {} Body text: use the full 9pt (size: 18). No compression — readability over density.
\item {} Line spacing: 1.15x on body paragraphs (slightly more generous than default). This template is the one where whitespace is a feature, not waste.
\item {} Company Overview and Strategic Priorities should take \textasciitilde{}40-50\% of the page.
\item {} Financial Snapshot: visually compact, secondary prominence. Keep it small and clean.
\item {} Key Relationships: consider two-column layout (Customers + Vendors on left, Partners + Competitors on right) to save vertical space while keeping descriptors visible.
\item {} Conversation Starters: use the indented block-style bullets. These are the payoff of the document — they should stand out visually.
\end{itemize}

\vspace{0.8em}\hrule\vspace{0.8em}

\end{document}